%% A sharp denominator theorem for the cellular rational approximations to zeta(5)
%% Draft. Compile with: pdflatex sharp12 (twice)
\documentclass[11pt,reqno]{amsart}

\usepackage[margin=1.15in]{geometry}
\usepackage{amsmath,amssymb,amsthm,mathrsfs}
\usepackage{booktabs}
\usepackage{enumitem}
\usepackage{microtype}
\usepackage{xcolor}
\definecolor{linknavy}{rgb}{0.10,0.20,0.50}
\usepackage[colorlinks=true,linkcolor=linknavy,citecolor=linknavy,urlcolor=linknavy]{hyperref}

\theoremstyle{plain}
\newtheorem{theorem}{Theorem}[section]
\newtheorem{proposition}[theorem]{Proposition}
\newtheorem{lemma}[theorem]{Lemma}
\newtheorem{corollary}[theorem]{Corollary}
\theoremstyle{definition}
\newtheorem{definition}[theorem]{Definition}
\newtheorem{remark}[theorem]{Remark}
\newtheorem{example}[theorem]{Example}
\theoremstyle{remark}
\newtheorem{evidence}[theorem]{Evidence}

\numberwithin{equation}{section}

\newcommand{\ZZ}{\mathbb{Z}}
\newcommand{\QQ}{\mathbb{Q}}
\newcommand{\NN}{\mathbb{N}}
\newcommand{\FF}{\mathbb{F}}
\newcommand{\Zp}{\ZZ_p}
\newcommand{\Qp}{\QQ_p}
\newcommand{\ord}{\operatorname{ord}}
\newcommand{\den}{\operatorname{den}}
\newcommand{\rank}{\operatorname{rank}}
\newcommand{\Res}{\operatorname{Res}}
\newcommand{\lcm}{\operatorname{lcm}}
\newcommand{\wt}{\operatorname{wt}}
\newcommand{\car}{\operatorname{car}}
\newcommand{\ii}[1]{[\![#1]\!]}
\newcommand{\Tcal}{\mathcal{T}}
\newcommand{\Ecal}{\mathcal{E}}
\newcommand{\Dcal}{\mathcal{D}}
\newcommand{\Mcal}{\mathcal{M}}
\newcommand{\LBZ}{L_{\mathrm{BZ}}}
\newcommand{\wfive}{w_5}
\newcommand{\wthree}{\hat w_3}
\newcommand{\vfive}{v_5}
\newcommand{\PROVED}{\textsf{[PROVED]}}
\newcommand{\VERIFIED}{\textsf{[VERIFIED]}}
\newcommand{\CERTIFIED}{\textsf{[CERTIFIED]}}
\newcommand{\CERT}{\textsf{[CERT]}}
\newcommand{\LEAN}{\textsf{[LEAN]}}
\newcommand{\OPEN}{\textsf{[OPEN]}}

\begin{document}

\title[A sharp denominator theorem for cellular approximations to $\zeta(5)$]
{A sharp denominator theorem for the cellular\\ rational approximations to $\zeta(5)$}

%% AUTHOR PLACEHOLDER --- to be filled in before circulation
\author{}
\address{}
\email{}

\date{\today}

\subjclass[2020]{11J82, 11B65, 11A07, 33F10, 14H10}
\keywords{Cellular integrals, $\zeta(5)$, denominators, Lucas congruences, Kummer's theorem,
Wolstenholme's congruence, creative telescoping}

\begin{abstract}
Brown and Zudilin analyse a family of cellular integrals on the moduli space
$\mathcal M_{0,8}$ whose totally symmetric member produces simultaneous rational
approximations $I'_n=Q_n\zeta(5)-P_n$ and $I''_n=Q_n\zeta(3)-\hat P_n$. On the basis of
extensive computation they record the inclusions
$Q_n,\;d_n^2d_{2n}\hat P_n,\;d_n^5P_n\in\ZZ$, where $d_n=\lcm(1,\dots,n)$. We complete this
picture. The inclusions hold after multiplication by the single global factor
$12=2^2\cdot 3$ --- as their own printed value $P_2=1190161/384$ already shows, since
$d_2^5=32$ and $384=12\cdot 32$ --- and the factor $12$ is sharp.

Our main theorem is the $p\ge5$ half of the resulting law:
$\ord_p(P_n)\ge-5\lfloor\log_p n\rfloor=-\ord_p(d_n^5)$ for every prime $p\ge5$ and every
$n\ge1$. The proof is a digit-scaled induction on the number of base-$p$ digits of $n$,
%% FIXED [CRITICAL/MAJOR M-1,M-2, m-1, m-3]: the abstract now counts TWO unproved inputs for
%% FIXED Theorem \ref{thm:main} --- (T1-top) and (DEPTH) --- names Theorem B as a third input used
%% FIXED only by the weight-3 companion, states the exact-over-$\QQ$ ranges, and no longer calls
%% FIXED Wolstenholme and $\sum j^{-2}$ the only arithmetic inputs simpliciter.
applied to the graded remainder $W_n=P_n-H^{(5)}_nQ_n$, and it rests on exactly two inputs
that are not proved here. The first is the finite decomposition identity
$P_n=\sum_{k,l}T(n,k,l)\,w_5(n,k,l)$ over the Brown--Zudilin summand
$T(n,k,l)=\binom{n+k}{n}\binom{n}{k}^2\binom{n+l}{n}\binom{n}{l}^2\binom{n+k+l}{n}$ of
$Q_n=\sum_{k,l}T(n,k,l)$, with $w_5$ an explicit weight-$5$ form in harmonic letters; it is
verified exactly over $\QQ$ for $n\le34$ (for the representative used in the base case,
$n\le20$), agrees with $P_n$ at every $n\le360$ and satisfies the Brown--Zudilin operator at
$748$ consecutive values modulo two $25$-bit primes, but is not yet machine-certified. The
second is the linear-algebra certificate (DEPTH) cutting out the depth-conditioned family in
which $w_5$ is chosen: an exact rank computation, reproduced at two auxiliary primes with
identical pivot sets, but not a human proof. We state both statuses precisely and isolate what
a certificate for the first must contain.

Everything else in that chain is proved: an exact Lucas congruence $Q_{ap+r}\equiv Q_aQ_r\pmod p$ for the
bottom row (formalised in Lean~4 with no axioms beyond \texttt{propext},
\texttt{Classical.choice}, \texttt{Quot.sound}), a graded $H^{(5)}$-layer reduction, a
level-independent depth calculus for $w_5$, an exact fibre block factorisation of $T$ whose
only external input is Wolstenholme's congruence, a descent Kummer lemma, and a base case
that dissolves at the midpoint $n=(p+1)/2$ into a three-cell corner sum equal to
$-14\sum_{j<p}j^{-2}\equiv0\pmod p$. Those two classical facts --- Wolstenholme, and
$\sum_{j<p}j^{-2}\equiv0$ --- are the only arithmetic inputs of the chain that constrain the
prime (Wilson's theorem is used as well, but holds at every prime and so costs nothing), and
they are exactly why $p\ge5$ is the right hypothesis. We also prove two structural obstruction
theorems that explain why several natural routes cannot work; and we prove the weight-$3$
companion statement for $\hat P_n$ modulo a third unproved input, the weight-$3$ decomposition
identity $\hat P_n=\sum_{k,l}T\,\hat w_3$ (Theorem B, likewise verified but not certified),
which the $p\ge5$ law for $P_n$ does \emph{not} use. The residual factor $12$ at $p\in\{2,3\}$ is of a different nature; we record its
status honestly, including the fact that the geometric mechanism proposed for it provably
cannot produce the factor $3$.
\end{abstract}

\maketitle

\setcounter{tocdepth}{1}
\tableofcontents

%%%%%%%%%%%%%%%%%%%%%%%%%%%%%%%%%%%%%%%%%%%%%%%%%%%%%%%%%%%%%%%%%%%%%%%%%%%%%%%%
\section{Introduction}
%%%%%%%%%%%%%%%%%%%%%%%%%%%%%%%%%%%%%%%%%%%%%%%%%%%%%%%%%%%%%%%%%%%%%%%%%%%%%%%%

\subsection{The Brown--Zudilin family}

In \cite{BZ} Brown and Zudilin study the family of five-fold cellular integrals
\[
I(\boldsymbol a)=I(a_1,\dots,a_8)
\]
of \cite[Example 7.5]{Br16}, period integrals on the moduli space $\mathcal M_{0,8}$ of genus
zero curves with eight marked points. General results on periods of these moduli spaces
\cite{Br09} force $I(\boldsymbol a)$ to be a $\QQ$-linear combination of multiple zeta values
of weight $\le5$; the cellular nature of the integral kills the subleading weight and pins the
leading weight to the single period $\zeta(5)+2\zeta(3)\zeta(2)$, and the coefficient of
$\zeta(3)$ vanishes identically. The upshot \cite[(deco)]{BZ} is a decomposition
\[
I=Q\cdot\bigl(2\zeta(5)+4\zeta(3)\zeta(2)\bigr)-4\hat P\cdot\zeta(2)-2P,
\qquad Q\in\ZZ,\ P,\hat P\in\QQ,
\]
and, after ``setting $\zeta(2)$ to zero'', a \emph{pair} of simultaneous rational
approximations
\[
I'(\boldsymbol a)=Q\zeta(5)-P,\qquad I''(\boldsymbol a)=Q\zeta(3)-\hat P .
\]

A large transformation group $\mathfrak G$ of order $7!$ acts on the normalised integrals and on
$Q,P,\hat P$, and it is by exploiting that action that Brown and Zudilin compute a sharp
denominator; such Rhin--Viola groups are a recurring and arithmetically decisive feature of
rational approximations to $\zeta(2)$, $\zeta(3)$ and $\zeta(4)$ \cite{RV96,RV01,MZ20}, as they
are of the hypergeometric constructions for odd zeta values generally \cite{Zu04,Zu02a}. The
whole subject descends from Ap\'ery \cite{Ap79} and Beukers' re-reading of it \cite{Be79}.

We work throughout with the \emph{totally symmetric} member $a_1=\dots=a_8=n$ of the family,
and write $Q_n$, $P_n$, $\hat P_n$ for the resulting coefficients, normalised as in \cite{BZ}:
\begin{equation}\label{eq:normalisation}
Q_0=1,\quad Q_1=21,\quad Q_2=2989;\qquad
P_1=\tfrac{87}{4},\quad P_2=\tfrac{1190161}{384};\qquad
\hat P_1=\tfrac{101}{4},\quad \hat P_2=\tfrac{344923}{96}.
\end{equation}
All three sequences satisfy one and the same Ap\'ery-type recursion of order $3$
\cite{Zu02b}, which we call $\LBZ$:
\begin{equation}\label{eq:LBZ}
c_0(\nu)Y_\nu+c_1(\nu)Y_{\nu+1}+c_2(\nu)Y_{\nu+2}+c_3(\nu)Y_{\nu+3}=0,
\end{equation}
\begin{equation}\label{eq:LBZcoeffs}
\begin{aligned}
c_0(\nu)&=(\nu+1)^5(\nu+2)\,a_0(\nu+1), &\qquad c_1(\nu)&=-2(\nu+2)B_8(\nu),\\
c_2(\nu)&=-2B_9(\nu), &\qquad c_3(\nu)&=2(\nu+3)^5(2\nu+5)\,a_0(\nu),
\end{aligned}
\end{equation}
with $a_0(x)=41218x^3+198849x^2+320790x+173057$ and $B_8,B_9$ explicit polynomials making all
four $c_i$ of degree $9$.

The bottom row has a closed form \cite[(Q\_n)]{BZ}: with
\begin{equation}\label{eq:T}
T(n,k,l):=\binom{n+k}{n}\binom{n}{k}^{2}\binom{n+l}{n}\binom{n}{l}^{2}\binom{n+k+l}{n},
\end{equation}
one has
\begin{equation}\label{eq:BZQ}
Q_n=\sum_{k,l=0}^{n}T(n,k,l)\in\ZZ .
\end{equation}
We refer to $T$ as the \emph{Brown--Zudilin summand}, and to the triple $(Q_n,\hat P_n,P_n)$
as the \emph{bottom}, \emph{middle} and \emph{top row} of the ladder.

\subsection{The denominator law, and the factor \texorpdfstring{$12$}{12}}

On the basis of extensive computation, Brown and Zudilin record
\cite[(dn-totsym)]{BZ}
\begin{equation}\label{eq:BZincl}
Q_n,\qquad d_n^2d_{2n}\hat P_n,\qquad d_n^5P_n\ \in\ \ZZ
\qquad\text{for } n=0,1,2,\dots,
\end{equation}
where $d_n=\lcm(1,2,\dots,n)$; the general-parameter analogues are their inclusions
\cite[(incl), (sharp\_incl)]{BZ}, which they describe as an experimental observation that ``it
is in principle possible to prove, with considerable effort''.

The purpose of this paper is to complete \eqref{eq:BZincl}. The inclusions hold after
multiplication by one global constant, and that constant is $12=2^2\cdot3$. This is already
visible in \cite{BZ}'s own printed value: $P_2=1190161/384$ and $d_2^5=2^5=32$, so
$d_2^5P_2=1190161/12\notin\ZZ$, while $12\,d_2^5P_2\in\ZZ$. The corrected law is sharp in
both directions.

%% FIXED [M-3]: the $c=4$ row (the one that shows the factor 3 is real) has been added;
%% FIXED recomputed independently, 330/361 with 31 failures, first at n=2,6,8,18,20,24,26,54.
\begin{evidence}[the empirical law]\label{ev:law12}
$\VERIFIED$ Over the $361$ exact ladder values $n=0,\dots,360$, the number of $n$ with
$\den(P_n)\mid c\,d_n^5$ is
\[
\begin{array}{c|ccccccc}
c & 1 & 2 & 3 & 4 & 6 & 12 & 24\\\hline
\#\{n\le360\} & 1 & 1 & 1 & 330 & 1 & 361 & 361
\end{array}
\]
So the constant is exactly $12$: it suffices, and neither of its maximal proper divisors $6$
and $4$ does. The two halves fail very differently. Dropping the $2^2$ to $6$ (hence also $3$,
$2$, $1$) fails at all $360$ values $n\ge1$; dropping the $3$ to $4$ fails at only $31$ values,
the first being $n=2,6,8,18,20,24,26,54$ --- so the factor $3$ is real but rare, which is
consistent with its provenance being the least understood part of the constant
(\S\ref{sec:twothree}). Moreover $\log\den(P_n)/\log(d_n^5)\to1$ (values
$0.9953,0.9928,0.9947,0.9936$ at $n=80,200,320,355$), so the exponent $5$ is itself tight.
\end{evidence}

\begin{remark}
We stress the framing. Brown and Zudilin state \eqref{eq:BZincl} explicitly as an experimental
observation, not as a theorem, and their paper is careful to flag which of its assertions are
proved and which are computational. What we do here is supply the missing constant and prove
the arithmetic half of the resulting statement; we regard this as completing their analysis
rather than correcting it. The $2$-adic and $3$-adic content of the constant $12$ is, as we
explain in \S\ref{sec:twothree}, of a genuinely different flavour from the $p\ge5$ part, and
Brown and Zudilin's own heuristic --- that the construction obtains $I'$ by motivically
setting $\zeta(2):=0$, a row operation against the period $(2\pi i)^2$ whose integral
normalisation carries a Bernoulli denominator --- is precisely the mechanism one would expect
to leave a universal $2^2\cdot3$ behind.
\end{remark}

\subsection{The main theorem}

Throughout, $p$ denotes a prime, $v_p$ (equivalently $\ord_p$) the $p$-adic valuation on
$\QQ$, and
\[
L=L(n):=\lfloor\log_p n\rfloor=v_p(d_n),\qquad H^{(m)}_N:=\sum_{i=1}^{N}i^{-m}.
\]

%% FIXED [M-1]: (DEPTH) added to the hypothesis list here, in \S16's root row and in
%% FIXED Remark \ref{rem:flip}; it is [CERT], not [PROVED].
\begin{theorem}[Sharp-$12$, $p\ge5$ part]\label{thm:main}
Assume the decomposition identity \textup{(T1-top)} of \S\ref{sec:decomp} and the
linear-algebra certificate \textup{(DEPTH)} of \S\ref{sec:depth}. Then for every
prime $p\ge5$ and every $n\ge1$,
\[
\ord_p(P_n)\ \ge\ -5\lfloor\log_p n\rfloor\ =\ -\ord_p(d_n^5).
\]
Equivalently $d_n^5P_n\in\ZZ[1/6]$, which is the $p\ge5$ half of $\den(P_n)\mid 12\,d_n^5$.
\end{theorem}

The two hypotheses are both finite, and isolated as sharply as we know how:

\begin{itemize}[leftmargin=1.6em]
\item \textbf{(T1-top).} $P_n=\sum_{k,l=0}^{n}T(n,k,l)\,\wfive(n,k,l)$, where $\wfive$ is an
  explicit $\QQ$-combination of weight-$5$ monomials in the harmonic letters
  $A_m,B_m,C_m,N_m$ of \S\ref{sec:decomp}. Evidence class $\VERIFIED$.
\item \textbf{(DEPTH).} The linear system of \S\ref{sec:depth} cutting out the
  depth-conditioned family in which $\wfive$ is chosen is consistent, with the ranks recorded
  in Theorem~\ref{thm:depthcons}. Evidence class $\CERT$: an exact rank computation reproduced
  at two auxiliary primes with identical pivot sets. It is what Proposition~\ref{prop:LIFT}
  and hence the depth bound (DEPTH-gen) rest on, and it is inherited by every statement below
  that quotes a depth bound.
\end{itemize}

Neither is $\PROVED$: see \S\ref{sec:certs} for exactly what has
been checked for (T1-top), what a creative-telescoping certificate would have to contain, and a
structural obstruction (Theorem~\ref{thm:degfit}) which shows why the cheap route that works at
weight $3$ has no weight-$5$ analogue. Section~\ref{sec:certs} is written so that a single
sentence flips when that certificate lands. We stress that the weight-$3$ decomposition
identity (Theorem~\ref{thm:B}) is \emph{not} among the hypotheses of
Theorem~\ref{thm:main}: nothing in the $P_n$ chain of \S\S\ref{sec:Qrow}--\ref{sec:nucleus}
consumes $\wthree$. Theorem~\ref{thm:B} is a third unproved input of this paper, used only for
the middle row (Theorem~\ref{thm:Eintro}).

Everything else in the chain is proved. The statements below record the main intermediate
results together with the hypotheses each one carries; of them only
Theorem~\ref{thm:Aintro} is unconditional.

\begin{theorem}[Bottom row]\label{thm:Aintro}
For every prime $p$, every $a\ge0$ and every $0\le r<p$,
\[
Q_{ap+r}\equiv Q_a\,Q_r \pmod p,
\]
hence $Q_n\equiv\prod_i Q_{n_i}\pmod p$ over the base-$p$ digits of $n$. Moreover
\[
Q_{ap+r}\equiv Q_a\bigl(Q_r+p\,a\,\Psi_a\bigr)\pmod{p^2}\qquad(p\ge5,\ 1\le a<p,\ p\nmid Q_a)
\]
with $\Psi_a$ the explicit level-$r$ weight-$1$ harmonic functional of
Corollary~\textup{\ref{cor:Qsuper}}.
\end{theorem}

\begin{theorem}[Base case]\label{thm:baseintro}
Assume \textup{(T1-top)} for the representative $w_5^{\mathrm I}$ of \S\ref{sec:decomp} and the
linear-algebra certificate \textup{(DEPTH)}. Then $\ord_p(P_n)\ge0$ for every prime $p\ge5$
and every $n<p$.
\end{theorem}

\begin{theorem}[One radical short, unconditionally in $n$]\label{thm:51intro}
Assume \textup{(T1-top)} and \textup{(DEPTH)}. For every prime $p\ge5$ and every $n\ge1$,
\[
\ord_p(P_n)\ \ge\ -5\lfloor\log_p n\rfloor-1,
\]
i.e.\ $d_n^5\cdot\!\!\prod_{5\le p\le3n}\!\!p\cdot P_n\in\ZZ[1/6]$; and $\ord_p(P_n)\ge0$
outright whenever $p>3n$.
\end{theorem}

Theorem~\ref{thm:51intro} does not use the base case or the off-regime descent, and it is
recorded separately because it is what the depth calculus alone buys.

\begin{theorem}[Middle row]\label{thm:Eintro}
Assume the weight-$3$ decomposition identity \textup{(Theorem B)} of \S\ref{sec:decomp}. Let
$p\ge5$, $1\le a<p$, $0\le r<p$, $n=ap+r$ and $p\nmid Q_a$. Then
\[
v_p\bigl(p^3\hat P_n-\hat P_a Q_r\bigr)\ \ge\ 1+\min\bigl(0,v_p(\hat P_a)\bigr).
\]
In particular $p^3\hat P_{ap+r}\equiv \hat P_aQ_r\pmod p$ if and only if $\hat P_a\in\Zp$.
\end{theorem}

Theorem~\ref{thm:Eintro} corrects a natural expectation: the na\"ive ``Frobenius
$\operatorname{diag}(1,p^3,p^5)$'' product congruence is \emph{false} in the middle slot, and
\S\ref{sec:middle} explains exactly why --- the letter $A_3(k)$ has upper argument $n+k$
reaching $2n$, which is the source of the $d_{2n}$ in \eqref{eq:BZincl}.

Finally, two obstruction theorems delimit the landscape; they are stated in
\S\ref{sec:wall} and are, we think, of independent interest.

\begin{theorem}[The $\theta$ wall, informally]\label{thm:wallintro}
The entire content of the base case is a linear functional $\theta_{p,n}$ on the space of
depth-conditioned weight-$5$ forms which \emph{factors through the value map}
$w\mapsto\sum_{k,l}T(n,k,l)w(n,k,l)$. Consequently no exact summation identity
$\sum_{k,l}T\cdot M=0$ --- no member of the ``Lemma $\Phi$ species'', including the new
weight-$2$ identities of Lemma~\textup{\ref{lem:Phi2}} --- can ever move $\theta$; and
$\theta\not\equiv0$ on that space. The depth calculus plus arbitrarily much summand-side
identity work therefore cannot prove the base case.
\end{theorem}

\subsection{What is new}

The proof of Theorem~\ref{thm:main} is a digit-scaled induction whose four ingredients are, we
believe, all new:
\begin{enumerate}[leftmargin=1.9em,label=\textup{(\arabic*)}]
\item a \emph{level-lifting} proposition (Proposition~\ref{prop:LIFT}) showing that a finite
  list of $42$ linear conditions on the coefficients of $\wfive$, computed at a single
  base-$p$ digit, controls the $p$-adic depth of $\wfive$ at \emph{every} digit level
  simultaneously and at every prime $p\ge5$ at once;
\item an exact block factorisation of $T$ over a base-$p$ fibre (Lemma~\ref{lem:B}) together
  with an exact residue identity (Lemma~\ref{lem:Phi}) that annihilates the entire first-order
  dependence of the fibre sum on the fibre coordinates, giving the sharp fibre congruence
  Lemma~\ref{lem:F} and its multi-digit weakening Lemma~\ref{lem:Fgen}; the only external
  input is Wolstenholme's congruence, used once;
\item a descent Kummer lemma (Lemma~\ref{lem:DK}) whose one-line mechanism is that
  \emph{off-regime means a carry in base-$p$ position $0$, the depth-pattern indicators live
  in position $L\ge1$, and Kummer's theorem counts both};
\item a base case that is not a cancellation argument at all above or below the midpoint, and
  at the midpoint reduces to three cells whose residues sum to $-14\sum_{j<p}j^{-2}$
  (Theorem~\ref{thm:nucleus}).
\end{enumerate}

Along the way we obtain: a $\bmod\ p^2$ supercongruence for the bottom row
(Corollary~\ref{cor:Qsuper}); an explicit order-$4$ left multiple $\tilde L$ of $\LBZ$ whose
leading coefficient is $2D(\nu+3)^5(2\nu+5)$, i.e.\ with the factor $a_0$ removed
(\S\ref{sec:desing}); a universal normalised midpoint row
$(11907,-334374,-19292)$ (\S\ref{sec:recstar}); new exact weight-$2$ residue identities on the
Brown--Zudilin summand (Lemma~\ref{lem:Phi2}); and two new harmonic closed forms for Ap\'ery's
$\zeta(3)$ numerators $b_n$ (\S\ref{sec:apery}), one of which is cell-wise $p$-integral.

\subsection{Conventions on evidence}\label{sec:evidence}

Because part of the chain is machine-assisted, we are explicit about evidence classes and use
them uniformly.

\begin{itemize}[leftmargin=1.6em]
\item $\PROVED$ --- a complete human proof is given here or in a named predecessor.
\item $\CERTIFIED$ --- a machine proof object (a creative-telescoping operator together with
  its certificate) has been produced \emph{and} independently re-checked by exact
  rational-function arithmetic in a session that does not trust the search.
\item $\CERT$ --- a linear-algebra certificate: an exact rank computation reproduced at two or
  three auxiliary primes with identical pivot sets.
\item $\VERIFIED\ r$ --- an exact finite check on the stated range $r$, $0$ failures. This is
  evidence, never proof. Ranges are always given.
\item $\LEAN$ --- formalised in Lean~4 with $0$ \texttt{sorry}s and no axioms beyond the three
  standard ones (\texttt{propext}, \texttt{Classical.choice} and \texttt{Quot.sound}).
\item $\OPEN$ --- not settled; stated precisely.
\end{itemize}

All numerical checks quoted below are exact (integer, \texttt{Fraction}, $\FF_q$, or tracked-precision
$p$-adic arithmetic); no floating point enters any arithmetic statement. Appendix~\ref{app:repro}
lists the scripts.

%%%%%%%%%%%%%%%%%%%%%%%%%%%%%%%%%%%%%%%%%%%%%%%%%%%%%%%%%%%%%%%%%%%%%%%%%%%%%%%%
\section{Structure of the proof}\label{sec:structure}
%%%%%%%%%%%%%%%%%%%%%%%%%%%%%%%%%%%%%%%%%%%%%%%%%%%%%%%%%%%%%%%%%%%%%%%%%%%%%%%%

This section is a narrative map. Nothing is proved here; every statement is proved in the
section named.

\subsection{The shape of the object}

The three rows of the ladder sit in a graded tower. The bottom row is integral and obeys an
exact Lucas congruence (\S\ref{sec:Qrow}). The upper rows are not integral, and their
denominators are exactly what we must control. The first move (\S\ref{sec:Hlayer}) is to strip
off the part of $P_n$ that is visibly responsible for the pole: writing
\[
W_n:=P_n-H^{(5)}_n Q_n,
\]
the ``$H$-layer'' $H^{(5)}_nQ_n$ contributes $-5L$ for free because $Q_n\in\ZZ$, and all
remaining content sits in $W_n$. The second move is to give $W_n$ a summand-level shape:
\[
W_n=\sum_{k,l}T(n,k,l)\,\vfive(n,k,l),\qquad \vfive:=\wfive-H^{(5)}_n,
\]
which is exactly the decomposition identity (T1-top). Everything downstream is a
\emph{cell-by-cell ledger}: a comparison, at each summand, between the $p$-adic \emph{depth}
of $\vfive$ (how negative its valuation can be) and the Kummer \emph{carries} of $T$ (how many
powers of $p$ the summand supplies).

\subsection{The three walls, and what they forced}

The route to Theorem~\ref{thm:main} was not the expected one, and three hard negatives shaped
it. We state them because each is a theorem, and each closes off a natural approach.

\begin{enumerate}[leftmargin=1.9em,label=\textup{(W\arabic*)}]
\item \textbf{Lemma $F$ is sharp} (\S\ref{sec:F}). The fibre congruence delivers relative
  precision $p^{2+\min(v_pT,2)}$ and that value is \emph{attained}; so no strengthening of the
  fibre analysis can pay for the weight-$5$ depth. The fix had to come from the $\wfive$ side,
  and it is a \emph{linear} fix: the depth conditions of \S\ref{sec:depth}.
\item \textbf{No $p$-independent representative makes the base case cell-wise}
  (\S\ref{sec:depth}, Theorem~\ref{thm:DEPTHplus}). The strengthened depth systems
  $(\mathrm{DEPTH}^+)$ and $(\mathrm{DEPTH}^{++})$ are \emph{inconsistent} with the fitting
  system, at two auxiliary primes, and the inconsistency localises to exactly the three pole
  patterns carrying the deficit. Enlarging the alphabet by Ap\'ery-type letters
  $R^{(a)}$ or by genuinely nested depth-$2$ letters $Y_{ab},V_{ab}$ does not help
  (Theorem~\ref{thm:rletter}).
\item \textbf{The $\theta$ wall} (\S\ref{sec:wall}, Theorem~\ref{thm:wall}). The obstruction
  functional factors through the value map, so no summand-side identity can move it. This
  \emph{explains} (W2) rather than merely recording it: those linear systems were trying to
  buy, with $p$-independent linear conditions, something that is not $p$-independent-linear.
\end{enumerate}

The walls all point the same way: the base case is a statement about the \emph{number} $P_n$,
so a proof must import a property of $P_n$. There are exactly two such properties available:
the summand decomposition, and the recurrence $\LBZ$. The proof in \S\ref{sec:nucleus} uses
both.

\subsection{The nucleus dissolves at the midpoint}

The base case is $\ord_p(P_n)\ge0$ for $n<p$. The cell-wise bound the depth calculus supplies
is $-1$, and the deficit is \emph{attained}, at $(n,k,l)=\bigl(\tfrac{p+1}{2},0,\tfrac{p-1}{2}\bigr)$
for every $p$. So a genuine cancellation is needed. The proof splits the range at the midpoint
and finds that the cancellation is a three-term identity:

\begin{itemize}[leftmargin=1.6em]
\item For $n\le(p-1)/2$ the region $\mathrm{III}=\{k,l\ge q:=p-n,\ p\le k+l<p+q\}$ where the
  chosen representative $w_5^{\mathrm I}$ can fail to be cell-wise integral is \emph{empty}: it
  requires $2n\ge p$. Half the base case is free.
\item At $n=(p+1)/2$ the region collapses to the \emph{three corner cells}
  $(n-1,n)$, $(n,n-1)$, $(n,n)$. There $T/p^2\equiv(2,2,24)$ by Wilson's theorem, Lucas'
  theorem and Kummer's theorem, and the relevant coefficients of the pole expansion are
  $K_3=(3-s_2/2,\,3-s_2/2,\,-1/2-s_2/2)$ where $s_2=\sum_{j<p}j^{-2}$; the corner sum is
  \[
  2\Bigl(3-\tfrac{s_2}{2}\Bigr)+2\Bigl(3-\tfrac{s_2}{2}\Bigr)+24\Bigl(-\tfrac12-\tfrac{s_2}{2}\Bigr)
  \ =\ -14\,s_2\ \equiv\ 0\pmod p ,
  \]
  because $\sum_{j<p}j^{-2}\equiv\sum_{j<p}j^{2}=\tfrac{(p-1)p(2p-1)}{6}\equiv0$ for $p\ge5$.
  \emph{One classical input; no Fermat quotient, no Bernoulli number.}
\item Above the midpoint the forward $\LBZ$ induction has only one kind of exceptional step,
  the roots of $a_0$ modulo $p$, and those are \emph{apparent singularities}: all three
  $2\times2$ minors of a certain $2\times3$ coefficient matrix are divisible by $a_0(\nu)$ in
  $\QQ[\nu]$. Equivalently there is an explicit order-$4$ left multiple $\tilde L$ of $\LBZ$
  with $a_0$ removed from the leading coefficient (\S\ref{sec:desing}).
\end{itemize}

The two halves dovetail with no residue: \emph{on the desingularised ladder the only singular
step below $p-3$ is exactly the midpoint, which is exactly the single level the region
computation supplies.} A congruence $(\mathrm{REC}\text{-}\star)$ that had been expected to be
the \emph{input} to the base case comes out as its \emph{output}, and its universal normalised
row $(11907,-334374,-19292)$ turns out to be the same row that governs the $a=1$ midpoint band
one digit up in the prior campaign of this program \cite{ProgA1MID}.

\subsection{The induction}

With the base case in hand, the induction (\S\ref{sec:IND}) is a per-digit ledger. Writing
$n=ap+r$ and grouping the cells of $W_n$ by $(b,c)=(\lfloor k/p\rfloor,\lfloor l/p\rfloor)$,
\[
p^5W_n-Q_rW_a=\underbrace{\sum_{k,l}T(n,k,l)\,\Ecal(n,k,l)}_{\text{(I) descent}}
+\underbrace{\sum_{b,c}\vfive(a,b,c)\bigl[\Tcal(b,c)-Q_rT(a,b,c)\bigr]}_{\text{(II) fibre}} ,
\]
$\Ecal=p^5\vfive(n,k,l)-\vfive(a,b,c)$ and $\Tcal(b,c)=\sum_{s,t}T(n,bp+s,cp+t)$. Term (II) is
paid for by the depth bound (DEPTH-gen) against the fibre congruence Lemma~\ref{lem:Fgen},
with slack \emph{exactly zero}. Term (I) splits by regime: in-regime the letter descent is
exact and costs nothing; off-regime it is Lemma~\ref{lem:DK}. Every line of the ledger is
tight; there is no spare power of $p$ anywhere in the chain.

Two remarks on that. First, the route the off-regime term was expected to take --- a
letter-wise carry ledger --- is \emph{refuted} (\S\ref{sec:trap}): at weight $5$ the mismatch
of the letter $A_m(k)$ costs $m\lambda$ against a total Kummer gain of $1+\lambda$, and at
$p=5$, $n=19$, $(k,l)=(6,0)$ the letter-wise ledger is short by exactly one power. The pole is
not in $\Ecal$ at all; it is annihilated inside $\vfive$ by the depth conditions at level $n$.
Second, the fact that the two sharpnesses --- Lemma $F$'s and (DEPTH)'s --- meet exactly is not
a coincidence of bookkeeping but the reason the theorem is true at the stated exponent.

%%%%%%%%%%%%%%%%%%%%%%%%%%%%%%%%%%%%%%%%%%%%%%%%%%%%%%%%%%%%%%%%%%%%%%%%%%%%%%%%
\section{The bottom row}\label{sec:Qrow}
%%%%%%%%%%%%%%%%%%%%%%%%%%%%%%%%%%%%%%%%%%%%%%%%%%%%%%%%%%%%%%%%%%%%%%%%%%%%%%%%

\subsection{The closed form is certified}

\begin{theorem}\label{thm:Qcert}
$\CERTIFIED$ The double sum $\sum_{k,l}T(n,k,l)$ is annihilated by $\LBZ$; together with the
natural boundaries and agreement at $n=0,1,2$ this proves \eqref{eq:BZQ}.
\end{theorem}

The certificate was produced by two-step creative telescoping \cite{Zei91,WZ92} using the
package of \cite{Ko10}. The single call
$\mathtt{CreativeTelescoping}[\mathrm{ann},\{S_k{-}1,S_l{-}1\},\{S_n\}]$ returns
\texttt{\$Failed} even here, where a telescoper provably exists; the two-step split is the fix:
\[
\begin{gathered}
\mathrm{ann}=\mathtt{Annihilator}[T,\{S_n,S_k,S_l\}]\ \leadsto\
\mathrm{ct}_1=\mathtt{CreativeTelescoping}[\mathrm{ann},S_k-1,\{S_n,S_l\}]\\
\leadsto\ \mathtt{OreGroebnerBasis}\ \leadsto\ \mathrm{ct}_2 .
\end{gathered}
\]
The resulting order-$3$ telescoper has coefficients \emph{identically} those of
\eqref{eq:LBZcoeffs}: the coefficient ratio is $\{1,1,1,1\}$ and
$\mathtt{Expand}[\text{coeffs}-\LBZ]=\{0,0,0,0\}$. Every check was redone by exact
rational-function arithmetic on the explicit binomial expression for $T$ --- each annihilator
generator satisfies $\mathtt{Together}[(L\cdot T)/T]=0$, and each $k$-step pair satisfies
$\mathrm{tel}\cdot T+(S_k-1)(\mathrm{cert}\cdot T)=0$. Separately, a \emph{single}-certificate
(WZ-pair) form
\[
\LBZ\cdot T=\Delta_k(\rho T)+\Delta_l(\sigma T)
\]
with explicit rational $\rho,\sigma$ has been produced and checked to exactly $0$ twice, the
second time in a kernel that never loaded the package at all; this is the object every
weight-lowering step of \S\ref{sec:certs} consumes.

\subsection{Lucas for the bottom row}

The following is the ``$T$-level'' input used everywhere below. Its proof is uniform in $p$
(no prime is excluded) and, importantly, uniform in $a$ --- it is already multi-digit.

\begin{lemma}[Lucas step]\label{lem:L1}
For $a,b\ge0$ and $0\le r,s<p$: $\binom{ap+r}{bp+s}\equiv\binom{a}{b}\binom{r}{s}\pmod p$.
\end{lemma}

\begin{proof}
Lucas' theorem applied twice: first to the last digit $s$ against $r$, then to the reassembled
higher digits.
\end{proof}

\begin{lemma}[Single carry annihilation]\label{lem:L2}
With $N=ap+r$, $k=bp+s$: if $r+s<p$ then $\binom{N+k}{N}\equiv\binom{a+b}{a}\binom{r+s}{r}$,
and if $r+s\ge p$ then $\binom{N+k}{N}\equiv0\pmod p$.
\end{lemma}

\begin{proof}
If $r+s<p$ then $N+k=(a+b)p+(r+s)$ and Lemma~\ref{lem:L1} applies. If $r+s\ge p$ then
$N+k=(a+b+1)p+(r+s-p)$ with $0\le r+s-p<p$, and Lemma~\ref{lem:L1} gives
$\binom{N+k}{N}\equiv\binom{a+b+1}{a}\binom{r+s-p}{r}$; but $s<p$ forces $r+s-p<r$, so
$\binom{r+s-p}{r}=0$.
\end{proof}

\begin{lemma}[Double carry annihilation]\label{lem:L3}
Assume $s\le r$, $t\le r$, $r+s<p$, $r+t<p$. Then $s+t<p$, and
$\binom{N+k+l}{N}\equiv\binom{a+b+c}{a}\binom{r+s+t}{r}\pmod p$ if $r+s+t<p$, while
$\binom{N+k+l}{N}\equiv0\pmod p$ otherwise.
\end{lemma}

\begin{proof}
First, $s+t<p$. Suppose $s+t\ge p$. From $s\le r$ and $t\le r$ we get $p\le s+t\le 2r$, so
$r\ge p/2$; from $r+s<p$ and $r+t<p$ we get $s,t\le p-1-r$, so $s+t\le2(p-1-r)<2(p-p/2)=p$, a
contradiction. Hence $s+t<p$, so $r+s+t<2p$ and $N+k+l=(a+b+c+\varepsilon)p+(r+s+t-\varepsilon p)$
with $\varepsilon\in\{0,1\}$. If $\varepsilon=0$, Lemma~\ref{lem:L1} gives the stated product.
If $\varepsilon=1$, the low digit is $\rho=r+s+t-p$, and $\rho<r$ precisely because $s+t<p$;
Lemma~\ref{lem:L1} then produces a factor $\binom{\rho}{r}=0$.
\end{proof}

The deduction ``$s+t<p$'' in Lemma~\ref{lem:L3} is the non-obvious step, and it is what
upgrades a previously verified ``carry-kill predicate'' to a proof.

\begin{lemma}[Summand factorisation]\label{lem:L4}
With $N=ap+r$, $k=bp+s$, $l=cp+t$ as above,
\[
T(N,bp+s,cp+t)\equiv\ii{r+s+t<p}\cdot T_{\mathrm{hi}}(a,b,c)\,T_{\mathrm{lo}}(r,s,t)\pmod p,
\]
\[
T_{\mathrm{hi}}(a,b,c)=\tbinom{a+b}{a}\tbinom{a}{b}^2\tbinom{a+c}{a}\tbinom{a}{c}^2\tbinom{a+b+c}{a},
\qquad
T_{\mathrm{lo}}(r,s,t)=\tbinom{r+s}{r}\tbinom{r}{s}^2\tbinom{r+t}{r}\tbinom{r}{t}^2\tbinom{r+s+t}{r}.
\]
\end{lemma}

\begin{proof}
$\binom{N}{k}^2\equiv\binom{a}{b}^2\binom{r}{s}^2$ and $\binom{N}{l}^2\equiv\binom{a}{c}^2\binom{r}{t}^2$
by Lemma~\ref{lem:L1}; these vanish unless $s\le r$ and $t\le r$. The factors
$\binom{N+k}{N}$, $\binom{N+l}{N}$ are handled by Lemma~\ref{lem:L2} (they vanish unless
$r+s<p$, $r+t<p$). On the surviving regime Lemma~\ref{lem:L3} applies to $\binom{N+k+l}{N}$.
Multiplying the five congruences, and noting that when any factor is $\equiv0$ both sides
vanish, gives the statement.
\end{proof}

\begin{proof}[Proof of Theorem~\textup{\ref{thm:Aintro}}, Lucas part]
Write every summation index in base $p$: $k=bp+s$, $l=cp+t$ with $0\le s,t<p$; since
$k,l\le N<(a+1)p$ we have $0\le b,c\le a$. Summing Lemma~\ref{lem:L4} over
$0\le b,c\le a$ and $0\le s,t<p$,
\[
Q_N\equiv\sum_{b,c=0}^{a}\ \sum_{s,t:\,r+s+t<p}T_{\mathrm{hi}}(a,b,c)\,T_{\mathrm{lo}}(r,s,t)\pmod p .
\]
Two boundary points. First, the true index region is
$\{(b,s):bp+s\le N\}\times\{(c,t):cp+t\le N\}$ intersected with the surviving set; completing
to the full box adds only terms with $b=a$, $s>r$ (for which $\binom{r}{s}=0$), and
identically for $(c,t)$, so the completed sum is congruent to the true one. Second, the
$(b,c)$-region and the $(s,t)$-region are independent, since $r+s+t<p$ constrains only $s,t$;
so the double sum factors. The high factor is $\sum_{b,c=0}^{a}T_{\mathrm{hi}}(a,b,c)=Q_a$, an
\emph{integer identity} (it is \eqref{eq:BZQ} at $n=a$). For the low factor,
$Q_r=\sum_{s,t=0}^{r}T_{\mathrm{lo}}(r,s,t)$, and every term with $r+s+t\ge p$ vanishes mod
$p$: if $r+s\ge p$ or $r+t\ge p$ this is the computation of Lemma~\ref{lem:L2}; otherwise
Lemma~\ref{lem:L3} applies and $\binom{r+s+t}{r}\equiv0$. Hence
$\sum_{s,t:\,r+s+t<p}T_{\mathrm{lo}}\equiv Q_r$ and $Q_N\equiv Q_aQ_r$. The multi-digit
product form follows by induction on the number of base-$p$ digits, since the two-digit
statement holds for \emph{all} $a\ge0$.
\end{proof}

\begin{evidence}\label{ev:A}
$\PROVED$, and $\LEAN$: the statement $Q_{ap+r}\equiv Q_aQ_r\pmod p$ and its full digit-product
form are formalised in Lean~4 with $0$ \texttt{sorry}s, no \texttt{native\_decide}, and axioms
exactly $\{\texttt{propext},\texttt{Classical.choice},\texttt{Quot.sound}\}$. The same
development formalises the Ap\'ery-row analogue $a_{ap+r}\equiv a_aa_r$, an abstract
$\chi$-twisted weight-$w$ Frobenius descent theorem, and its minimal-Ap\'ery
($p^3b_{ap+r}\equiv b_aa_r$) and Franel ($p^2B_{ap+r}\equiv B_aA_r$) instances. In the
minimal-Ap\'ery instance the formalised object is $b_{\min}(n)=\sum_k\binom nk^2\binom{n+k}k^2
(2H^{(3)}_n-H^{(3)}_k)$; the identity $b_{\min}=b$ is proved on paper but is not formalised, and
likewise for the Franel analogue. The $\bmod\ p^2$ supercongruence of
Corollary~\ref{cor:Qsuper} is not formalised either.
Independently, $\VERIFIED$ $0$ failures for $p\in\{5,\dots,31\}$, all single-digit cells and
iterated digits, $n\le360$.
\end{evidence}

%%%%%%%%%%%%%%%%%%%%%%%%%%%%%%%%%%%%%%%%%%%%%%%%%%%%%%%%%%%%%%%%%%%%%%%%%%%%%%%%
\section{The graded \texorpdfstring{$H$}{H}-layer}\label{sec:Hlayer}
%%%%%%%%%%%%%%%%%%%%%%%%%%%%%%%%%%%%%%%%%%%%%%%%%%%%%%%%%%%%%%%%%%%%%%%%%%%%%%%%

\begin{lemma}[Digit split of a harmonic sum]\label{lem:H}
For $x\ge0$ and $m\ge1$,
$H^{(m)}_x=U_m(x)+p^{-m}H^{(m)}_{\lfloor x/p\rfloor}$ with
$U_m(x):=\sum_{j\le x,\ p\nmid j}j^{-m}\in\Zp$.
\end{lemma}

\begin{proof}
Split the summation range according to whether $p\mid j$.
\end{proof}

\begin{theorem}[$H$-layer reduction]\label{thm:C}
Let $p\ge5$, $1\le a<p$, $0\le r<p$, $n=ap+r$ \textup{(}so $n<p^2$\textup{)}. Put
$W_n:=P_n-H^{(5)}_nQ_n$ and $\hat W_n:=\hat P_n-H^{(3)}_nQ_n$. Then
\[
p^5P_n-P_aQ_r\equiv p^5W_n-W_aQ_r\pmod p,\qquad
p^3\hat P_n-\hat P_aQ_r\equiv p^3\hat W_n-\hat W_aQ_r\pmod p,
\]
each congruence being the assertion that the difference of its two sides lies in $p\Zp$.
\end{theorem}

%% FIXED [m-5]: ``both sides being $p$-integral'' was asserted, not shown, and is not free at
%% FIXED this point in the exposition. The statement now asserts only what the proof gives.
\begin{remark}
The difference of the two sides is $p^5H^{(5)}_nQ_n-H^{(5)}_aQ_aQ_r$ (respectively
$p^3H^{(3)}_nQ_n-H^{(3)}_aQ_aQ_r$), and both of \emph{those} terms are $p$-integral for free,
since $v_p(H^{(m)}_N)\ge-m\lfloor\log_pN\rfloor$ gives $p^5H^{(5)}_n,\,p^3H^{(3)}_n\in\Zp$ for
$n<p^2$, while $H^{(m)}_a\in\Zp$ for $a<p$ and $Q_n,Q_a,Q_r\in\ZZ$. That is what the proof
below establishes, and it is unconditional. The $p$-integrality of the individual sides is
\emph{not} available here: $\ord_p(P_a)\ge0$ for $a<p$ is the base case
(Theorem~\ref{thm:baseintro}), $\ord_p(p^5P_n)\ge0$ for $n<p^2$ is Theorem~\ref{thm:main} at
$L=1$, and the middle-row analogues are Proposition~\ref{prop:Phatsize} and
Theorem~\ref{thm:Eintro} --- all of them later, and all of them conditional. Nothing below
needs them.
\end{remark}

\begin{proof}
For $n=ap+r<p^2$ the indices $m\le n$ divisible by $p$ are exactly $m=jp$ with
$1\le j\le\lfloor n/p\rfloor=a$, and $p\nmid j$ since $j\le a<p$. Hence
$H^{(5)}_n=S+p^{-5}H^{(5)}_a$ with $S:=\sum_{m\le n,\ p\nmid m}m^{-5}\in\Zp$, so
$p^5H^{(5)}_n=p^5S+H^{(5)}_a\equiv H^{(5)}_a\pmod{p^5}$. Since $Q_n\in\ZZ$, we get
$p^5H^{(5)}_nQ_n\equiv H^{(5)}_aQ_n\pmod{p^5}$; reducing mod $p$ and using
Theorem~\ref{thm:Aintro} together with $H^{(5)}_a\in\Zp$ (as $a<p$) gives
$p^5H^{(5)}_nQ_n\equiv H^{(5)}_aQ_aQ_r\pmod p$. On the other side
$P_aQ_r=H^{(5)}_aQ_aQ_r+W_aQ_r$. Subtract. The weight-$3$ statement is identical with
$5\mapsto3$.
\end{proof}

Theorem~\ref{thm:C} is unconditional and weight-agnostic: it needs nothing but
Theorem~\ref{thm:Aintro} and a two-line valuation computation, and it works verbatim at every
weight. Its consequence is that the entire ``$p^5$ pole'' of the one-digit product congruence
\[
(\mathrm{LB}_5)\qquad p^5P_{ap+r}\equiv P_a\,Q_r\pmod p\qquad(p\ge5,\ 1\le a<p,\ 0\le r<p)
\]
is free, and $(\mathrm{LB}_5)$ is \emph{equivalent} to the same statement for $W$:
\[
(\mathrm{W}5)\qquad p^5W_{ap+r}\equiv W_a\,Q_r\pmod p .
\]

\begin{evidence}\label{ev:W5}
$\VERIFIED$ $0$ failures: $(\mathrm{W}5)$ holds for every $p\in\{5,7,11,13,17,19,23,29,31\}$ and
every single-digit cell $n=ap+r\le360$, with
$\min v_p(p^5W_n-W_aQ_r)=1$ exactly at every one of those primes --- so $(\mathrm{W}5)$ is
sharp and carries the full content. Also $\min_{n<p^2}v_p(W_n)=-5$ at every one of those
primes. The prime $p=5$ is not exceptional.
\end{evidence}

\begin{remark}
The induction of \S\ref{sec:IND} does \emph{not} use $(\mathrm{LB}_5)$. As assembled, the last
step of $(\mathrm{LB}_5)$ needs $v_p((\Lambda-Q_r)W_a)\ge1$, i.e.\ it needs the base case; the
inequality-shaped induction on $v_p(W_n)$ avoids this entirely and needs only
Lemma~\ref{lem:Fgen}, which is unconditional. We record $(\mathrm{LB}_5)$ because it is the
sharp statement at $n<p^2$ and because the middle-row theorem is of that shape.
\end{remark}

%%%%%%%%%%%%%%%%%%%%%%%%%%%%%%%%%%%%%%%%%%%%%%%%%%%%%%%%%%%%%%%%%%%%%%%%%%%%%%%%
\section{Carry lemmas for the Brown--Zudilin summand}\label{sec:carry}
%%%%%%%%%%%%%%%%%%%%%%%%%%%%%%%%%%%%%%%%%%%%%%%%%%%%%%%%%%%%%%%%%%%%%%%%%%%%%%%%

In the $\zeta(3)$ template the summand $\binom nk^2\binom{n+k}{k}^2$ is a perfect square, so a
single Kummer carry already gives $v_p\ge2$. The Brown--Zudilin summand \eqref{eq:T} has
\emph{three unsquared} binomials, so a single carry gives only $v_p\ge1$. The replacement is
that one digit overflow forces carries in two \emph{different} unsquared binomials.

\begin{lemma}[Triple-carry slack]\label{lem:D}
Let $p$ be prime, $N=ap+r$ with $0\le a,r<p$, and $k=bp+s$, $l=cp+t$ with $0\le s,t<p$ and
$0\le b,c\le a$ \textup{(}automatic when $k,l\le N$\textup{)}. If $a+b\ge p$ then
$v_p(T(N,k,l))\ge2$.
\end{lemma}

\begin{proof}
By Kummer's theorem $v_p\binom{N+k}{N}$ is the number of carries in the base-$p$ addition
$N+k$. The position-$0$ carry is $\varepsilon_0=\ii{r+s\ge p}$; the position-$1$ sum is
$a+b+\varepsilon_0\ge a+b\ge p$, so there is a carry at position $1$. Hence
$v_p\binom{N+k}{N}\ge1$, with $\ge2$ if $r+s\ge p$. So assume, else we are done:
(i) $r+s<p$; (ii) $s\le r$ and $t\le r$ --- otherwise $s>r$ (resp.\ $t>r$) forces a borrow in
$N-k$ (resp.\ $N-l$), so $v_p\binom{N}{k}\ge1$ and the \emph{squared} factor already gives
$v_p\ge2$; (iii) $a+c<p$ --- otherwise the position-$1$ argument applied to $\binom{N+l}{N}$
gives another power; (iv) $r+t<p$ --- otherwise $\binom{N+l}{N}$ has a position-$0$ carry.
From (i),(ii),(iv), $s+t<p$: if $s+t\ge p$ then $p\le s+t\le 2r$ forces $r\ge p/2$, while
$r+s<p$, $r+t<p$ force $s,t\le p-1-r$, hence $s+t\le2(p-1-r)<p$, a contradiction. From (iii)
and $b\le a$, $b+c\le a+c<p$. Therefore $k+l=(b+c)p+(s+t)$ is the base-$p$ expansion, and the
position-$1$ sum in $N+(k+l)$ is $a+(b+c)+\ii{r+s+t\ge p}\ge a+b\ge p$: a carry. So
$v_p\binom{N+k+l}{N}\ge1$ as well, and $v_pT\ge2$.
\end{proof}

\begin{lemma}[Refined triple carry]\label{lem:Dplus}
With $N,k,l$ as above set $\beta:=\lfloor(N+k)/p\rfloor=a+b+\ii{r+s\ge p}$. If $\beta\ge p$
then $v_p(T(N,k,l))\ge2$ for every $l$.
\end{lemma}

\begin{proof}
If $a+b\ge p$ this is Lemma~\ref{lem:D}. Otherwise $\beta\ge p$ forces $\ii{r+s\ge p}=1$ and
$a+b=p-1$; then $N+k$ carries at position $0$ and at position $1$ ($a+b+1=p$), so
$v_p\binom{N+k}{N}\ge2$ by Kummer alone.
\end{proof}

\begin{lemma}[Boundary carry]\label{lem:Dpp}
Let $p$ be prime, $N=ap+r$, $k=bp+s$, $l=cp+t$ with $0\le a,r,s,t<p$, $0\le b,c\le a$,
$k,l\le N$, and put $\beta=\lfloor(N+k)/p\rfloor$. If $\beta\ge p$ \emph{and} $a+b<p$ ---
equivalently $a+b=p-1$ and $r+s\ge p$ --- then $v_p(T(N,k,l))\ge4$.
\end{lemma}

\begin{proof}
$a+b=p-1$ and $r+s\ge p$ make $N+k=p^2+(r+s-p)$: the addition carries at position $0$ and at
position $1$, so by Kummer
\begin{equation}\label{eq:Dppstar}
v_p\tbinom{N+k}{N}=2 .
\end{equation}
Two further powers are needed. Note $b\le a$ and $a+b=p-1$ force $2a+1\ge p$.

\emph{Case $s>r$.} The subtraction $N-k$ borrows at position $0$, so $v_p\binom{N}{k}\ge1$ and
the squared factor supplies $2$; with \eqref{eq:Dppstar}, $v_pT\ge4$.

\emph{Case $s\le r$, $s+t<p$.} Then $k+l=(b+c)p+(s+t)$ is the base-$p$ expansion and
$r+s+t\ge r+s\ge p$, so $N+(k+l)$ carries at position $0$; its position-$1$ sum is
$a+(b+c)+1\ge a+b+1=p$, so it carries there too. Hence $v_p\binom{N+k+l}{N}\ge2$.

\emph{Case $s\le r$, $s+t\ge p$.} From $s\le r$ and $s+t\ge p$ we get $t\ge p-s\ge p-r$, i.e.
\begin{equation}\label{eq:Dppstst}
r+t\ge p ,
\end{equation}
so $N+l$ carries at position $0$ and $v_p\binom{N+l}{N}\ge1$. Write $\beta':=b+c+1$, the
position-$1$ digit sum of $k+l$ (here $k+l=\beta'p+(s+t-p)$).
If $\beta'<p$, the position-$1$ sum of $N+(k+l)$ is
$a+\beta'+\ii{r+s+t\ge2p}=p+c+\ii{\cdot}\ge p$ (using $a+b=p-1$), a carry, so
$v_p\binom{N+k+l}{N}\ge1$; total $2+1+1=4$. If $\beta'\ge p$ then $c\ge p-1-b=a$, so $c=a$; by
\eqref{eq:Dppstst} the position-$0$ carry of $N+l$ exists and its position-$1$ sum is
$a+c+1=2a+1\ge p$, a second carry, so $v_p\binom{N+l}{N}\ge2$ and with \eqref{eq:Dppstar},
$v_pT\ge4$.
\end{proof}

\begin{evidence}\label{ev:D}
Lemma~\ref{lem:D}: $\VERIFIED$ $24\,741$ cases ($p\in\{5,7\}$, all $a,r$, all $k,l\le N$), $0$
failures. The three-carry strengthening is \emph{false}: ``$a+b\ge p$ and $a+c\ge p$'' does
\emph{not} give $v_p\ge3$ ($649$ counterexamples in the same range), so Lemma~\ref{lem:D} is
exactly the available slack. Lemma~\ref{lem:Dplus}: $\VERIFIED$ $28\,128$ cases, $0$ failures.
Lemma~\ref{lem:Dpp}: $\VERIFIED$ $0$ failures with $\min v_p(T)=4$ exactly (hence sharp) at
$p=5,7,11,13$ on $518/2\,869/28\,340/65\,918$ index triples, of which
$174/892/8\,175/18\,620$ lie in the last sub-case $s\le r$, $s+t\ge p$. By contrast,
$\beta\ge p$ \emph{with} $a+b\ge p$ gives $\min v_p(T)=2$ exactly over $1.3$M triples ---
Lemma~\ref{lem:D} is sharp there.
\end{evidence}

%%%%%%%%%%%%%%%%%%%%%%%%%%%%%%%%%%%%%%%%%%%%%%%%%%%%%%%%%%%%%%%%%%%%%%%%%%%%%%%%
\section{The decomposition identities}\label{sec:decomp}
%%%%%%%%%%%%%%%%%%%%%%%%%%%%%%%%%%%%%%%%%%%%%%%%%%%%%%%%%%%%%%%%%%%%%%%%%%%%%%%%

\subsection{The alphabet}

For $n\ge0$, $0\le k,l\le n$ and $1\le m\le5$ define the \emph{letters}
\begin{equation}\label{eq:letters}
A_m(x):=H^{(m)}_{n+x}-H^{(m)}_{x},\qquad
B_m(x):=H^{(m)}_{n-x}-H^{(m)}_{x},\qquad
C_m:=H^{(m)}_{n+k+l}-H^{(m)}_{k+l},\qquad
N_m:=H^{(m)}_{n},
\end{equation}
with $x\in\{k,l\}$; each has \emph{weight} $m$. These are exactly the parameter
log-derivatives of the $\Gamma$-quotient form of $T$, which is why they are the right
alphabet: with
\[
\begin{aligned}
F(n,k,l;\alpha_k,\alpha_l,\beta_k,\beta_l,g)
&=\frac{\Gamma(n+k+1+\alpha_k)}{\Gamma(n+1)\Gamma(k+1+\alpha_k)}
\Bigl(\frac{\Gamma(n+1)}{\Gamma(k+1+\beta_k)\Gamma(n-k+1-\beta_k)}\Bigr)^{2}\\
&\qquad\times(\text{the same with } k\mapsto l)\times
\frac{\Gamma(n+k+l+1+g)}{\Gamma(n+1)\Gamma(k+l+1+g)},
\end{aligned}
\]
one has $\partial_{\alpha_k}^m\log F|_0=(-1)^{m+1}(m-1)!\,A_m(k)$,
$\partial_{\beta_k}^m\log F|_0=2(-1)^{m+1}(m-1)!\,B_m(k)$ and
$\partial_g^m\log F|_0=(-1)^{m+1}(m-1)!\,C_m$.

\subsection{The middle row: Theorem B}

\begin{theorem}[Theorem B]\label{thm:B}
$\VERIFIED$
\[
\hat P_n=\sum_{k,l=0}^{n}T(n,k,l)\,\wthree(n,k,l),
\]
where
\begin{equation}\label{eq:w3}
\begin{aligned}
\wthree(n,k,l)={}& H^{(3)}_n+A_3(k)+A_3(l)
-\tfrac14\bigl[A_2(k)A_1(k)+A_2(l)A_1(l)\bigr]\\
&-\tfrac34\bigl[A_2(k)B_1(k)+A_2(l)B_1(l)\bigr]
-\tfrac38\bigl[A_2(k)+A_2(l)\bigr]C_1
-\tfrac18\bigl[A_2(k)A_1(l)+A_2(l)A_1(k)\bigr].
\end{aligned}
\end{equation}
\end{theorem}

This is a decomposition of the middle row over the \emph{same} Brown--Zudilin summand that
gives $Q_n$, with a graded harmonic weight and no extra nested corrections. Two structural
features are worth recording, since they are not assumptions but outputs: every non-constant
term carries an $A$-letter of the highest weight in its monomial, and every monomial whose top
letter is $B$, $C$ or $H^{(r)}_n$ has coefficient $0$ --- the fit was run in a basis that
allowed all of these.

The single most consequential feature of \eqref{eq:w3} is that the letter
$A_r(k)=H^{(r)}_{n+k}-H^{(r)}_k$ has upper argument up to $2n$. \emph{This one fact is the
source of the $d_{2n}$ in \eqref{eq:BZincl} and of the entire middle-row anomaly of
\S\ref{sec:middle}}: $A_3(k)$ acquires a second $p$-pole as soon as
$\lfloor(n+k)/p\rfloor\ge p$, which for $n=a<p$ happens iff $a+k\ge p$.

\subsection{The top row: the family, and (T1-top)}

At weight $5$ the situation is subtler, and the history is instructive. A first search, in a
basis of $149$ symmetric weight-$5$ monomials with at most three factors, was inconsistent, and
this was read as evidence that the depth-$2$ period $\zeta(5)+2\zeta(3)\zeta(2)$ forces
genuinely nested letters. That reading is \emph{false}. Removing only the factor-count cap ---
same alphabet, no nested letters --- makes the system consistent.

\begin{proposition}[Existence]\label{prop:w5exists}
$\VERIFIED$ In the alphabet $\{A_r,B_r\}$ (at $k$ and $l$) $\cup\{C_r\}\cup\{N_r\}$,
$r=1,\dots,5$, with at most $5$ factors in each single-variable slot and at most $2$ in the
coupling and constant slots, the fitting system
$P_n=\sum_{k,l}T(n,k,l)\,\wfive(n,k,l)$ has $448$ basis monomials and rank $313$; it is
\emph{consistent}, with all $687$ independent excess equations satisfied at $N=1000$. The
solution space is a $135$-dimensional affine family.
\end{proposition}

The ablation that pins this is decisive. At $N=1200$ equations, caps
$(\mathrm{mf}_k,\mathrm{mf}_c,\mathrm{mf}_n)$:
\[
\begin{gathered}
(3,2,2)\ \text{inconsistent},\quad(3,2,5)\ \text{inconsistent},\quad(3,5,2)\ \text{inconsistent},\\
(5,2,2)\ \text{CONSISTENT},\quad(5,5,5)\ \text{CONSISTENT}.
\end{gathered}
\]
\emph{The missing ingredient is monomials with four or five factors in one variable} ---
e.g.\ $A_1(k)^4A_1(l)$, $A_1(k)^3B_1(k)^2$, $A_1(k)^5$. Likewise the alphabet ablation
(at $N=1000$) shows $\{A_r,B_r,C_r,N_r\}$ is exactly right: dropping $B$ or dropping $C$
is inconsistent, and $H^{(r)}_{2n}$ is not needed. \emph{So the alphabet of $\wfive$ is
exactly the alphabet of $\wthree$; the sole difference between the weights is the number of
factors} --- $\wthree$ uses at most two, $\wfive$ needs up to five.

The family is then cut down by the depth conditions of \S\ref{sec:depth} to a
$124$-dimensional subfamily, and inside it we fix explicit representatives.

%% FIXED [m-4]: the stored coefficient files are now named, so that a reader does not pick up
%% FIXED the 155-term intermediate (denominators $\{2,3,71\}$, not $p$-integral at $p=71$).
\begin{definition}[The representatives]\label{def:reps}
\leavevmode
\begin{itemize}[leftmargin=1.6em]
\item $w_5^{\mathrm{allp}}$: $178$ terms, coefficient denominators supported on $\{2,3\}$,
  hence $p$-integral for \emph{every} $p\ge5$. Obtained by a CRT combination of two
  depth-minimal representatives with disjoint bad primes (Lemma~\ref{lem:CRT}).
\item $w_5^{\mathrm I}$: $207$ terms, denominators on $\{2,3\}$; in addition to the depth
  conditions it satisfies the full cell-wise cap on every pole pattern except one band, which
  is what makes the base case of \S\ref{sec:nucleus} work.
\end{itemize}
Both are the CRT outputs, not the intermediates they are built from; in the accompanying
archive they are the files \texttt{w5\_allp.json} ($178$ terms) and
\texttt{w5\_exIII\_allp.json} ($207$ terms). The intermediate from which the second is built,
\texttt{w5\_exIII.json} (equivalently \texttt{w5\_I.json}), has $155$ terms and denominators
$\{2,3,71\}$, so it is \emph{not} $p$-integral at $p=71$ and must not be substituted for
$w_5^{\mathrm I}$.
\end{definition}

\begin{lemma}[CRT combination]\label{lem:CRT}
$\PROVED$ Let $x_1,x_2$ lie in the depth-conditioned family; suppose $x_1$ is $p$-integral for
all $p\ge5$ except $P_1$ (with maximal denominator exponent $e_1$) and $x_2$ except $P_2$
(exponent $e_2$), $P_1\ne P_2$. Put $d=x_2-x_1$ and choose $t\in\ZZ$ with
$t\equiv1\pmod{P_1^{e_1+1}}$, $t\equiv0\pmod{P_2^{e_2+1}}$. Then $x=x_1+td$ lies in the family
and is $p$-integral for every $p\ge5$.
\end{lemma}

\begin{proof}
At $P_1$, $x=x_2-(1-t)d$ with $v_{P_1}((1-t)d)\ge(e_1+1)-e_1=1$; at $P_2$, $x=x_1+td$ with
$v_{P_2}(td)\ge(e_2+1)-e_2=1$; at every other $p\ge5$ both $x_1$ and $d$ are $p$-integral and
$t\in\ZZ$.
\end{proof}

\noindent We can now state the first of the two hypotheses of Theorem~\ref{thm:main} (the
second, (DEPTH), is stated in \S\ref{sec:depth}).

\begin{quote}
\textbf{(T1-top).} $P_n=\displaystyle\sum_{k,l=0}^{n}T(n,k,l)\,\wfive(n,k,l)$ for all $n\ge0$,
where $\wfive$ is $w_5^{\mathrm{allp}}$ or $w_5^{\mathrm I}$ as appropriate.
\end{quote}

\begin{remark}[which representative]\label{rem:onerep}
The induction of \S\ref{sec:IND} consumes $\wfive$ \emph{only} through the depth bound
(DEPTH-gen), so it works for any point of the depth-conditioned family; and the base case of
\S\ref{sec:nucleus} uses $w_5^{\mathrm I}$. Hence the whole $p\ge5$ theorem can be run
end-to-end on the \emph{single} representative $w_5^{\mathrm I}$, and its certificate target is
then the \emph{one} identity $P_n=\sum_{k,l}T\cdot w_5^{\mathrm I}$. Alternatively one may
certify $w_5^{\mathrm{allp}}$ and add the homogeneous identity
$\sum_{k,l}T\cdot(w_5^{\mathrm I}-w_5^{\mathrm{allp}})=0$; these are genuinely different
obligations (see Proposition~\ref{prop:ptrank}), and which is cheaper is a computational
question, not a mathematical one.
\end{remark}

\subsection{New closed forms for Ap\'ery's \texorpdfstring{$\zeta(3)$}{zeta(3)} numerators}\label{sec:apery}

The residue construction behind \eqref{eq:w3} also produces pure harmonic-monomial closed forms
one weight down, replacing the classical non-generalising weight
$\sum_{m\le k}(-1)^{m-1}/(2m^3\binom nm\binom{n+m}{m})$.

\begin{proposition}\label{prop:apery-b}
$\PROVED$ For Ap\'ery's $\zeta(3)$ numerators \cite{Ap79},
\[
b_n=\sum_{j=0}^{n}\binom nj^2\binom{n+j}{n}^2
\Bigl[H^{(3)}_j+\bigl(2H_j-H_{n+j}-H_{n-j}\bigr)H^{(2)}_j\Bigr].
\]
\end{proposition}

\begin{proof}
Put $F(t)=\prod_{m=1}^{n}(t-m)\big/\prod_{m=0}^{n}(t+m)$, so that
$F(t)^2=\sum_{j=0}^{n}\bigl[\beta_{2,j}(t+j)^{-2}+\beta_{1,j}(t+j)^{-1}\bigr]$ with
$\beta_{2,j}=\binom nj^2\binom{n+j}{n}^2$ and
$\beta_{1,j}=\beta_{2,j}\cdot\frac{d}{dt}\log[(t+j)^2F(t)^2]\big|_{t=-j}
=2\beta_{2,j}(2H_j-H_{n+j}-H_{n-j})$. Since $\deg\mathrm{num}-\deg\mathrm{den}=-2$, the
residues sum to zero: $\sum_j\beta_{1,j}=0$. Now
$-(F^2)'(t)=\sum_j\bigl[2\beta_{2,j}(t+j)^{-3}+\beta_{1,j}(t+j)^{-2}\bigr]$, so summing over
$t\ge1$ and using $\sum_{t\ge1}(t+j)^{-m}=\zeta(m)-H^{(m)}_j$,
\[
\sum_{t\ge1}-(F^2)'(t)=2\Bigl(\sum_j\beta_{2,j}\Bigr)\zeta(3)+\Bigl(\sum_j\beta_{1,j}\Bigr)\zeta(2)
-\Bigl[2\sum_j\beta_{2,j}H^{(3)}_j+\sum_j\beta_{1,j}H^{(2)}_j\Bigr].
\]
The $\zeta(2)$ coefficient vanishes, $\sum_j\beta_{2,j}=a_n$ is Ap\'ery's identity, and the
linear form produced is $2(a_n\zeta(3)-b_n)$.
\end{proof}

%% FIXED [CRITICAL C-1]: the left-hand side was $a_n$; the displayed right-hand side evaluates
%% FIXED to Ap\'ery's NUMERATOR $b_n$ (0, 6, 351/4, 62531/36, ...) at every $n\le14$ (exact
%% FIXED Fraction recomputation), and to $a_n$ at no $n>0$. Root cause: work/p1g/zeta3.py and
%% FIXED work/p1g/z3ex.py name the numerator sequence ``a_n''. The proposition is now stated for
%% FIXED $b_n$ --- a SECOND closed form for the same sequence --- and \S1.3's promise of a closed
%% FIXED form ``for $a_n$'' has been withdrawn. The cell-wise integrality claim is unaffected and
%% FIXED was re-confirmed (819 cells, 0 violations).
\begin{proposition}\label{prop:apery-a}
$\PROVED$ over $\QQ$ Writing $A_1(k)=H_{n+k}-H_k$, $B_1(k)=H_{n-k}-H_k$,
$B_2(k)=H^{(2)}_{n-k}-H^{(2)}_k$, the same numerators $b_n$ also have the closed form
\[
b_n=\sum_{k=0}^{n}\binom nk^2\binom{n+k}{k}^2\Bigl[H^{(3)}_n
+\tfrac12\bigl(A_1(k)^2B_1(k)+A_1(k)B_1(k)^2+A_1(k)B_2(k)\bigr)\Bigr],
\]
i.e.\ with the weight $w_3=H^{(3)}_n+\tfrac12A_1\bigl(B_1(A_1+B_1)+B_2\bigr)$, and this weight
is \emph{cell-wise $p$-integral}: $d_3(n,k)\le v_pT_A(n,k)$ with $0$ violations for
$p\in\{5,7,11,13,17,19,23\}$, all $n<p$ and all $k$.
\end{proposition}

\begin{remark}
Propositions~\ref{prop:apery-b} and \ref{prop:apery-a} are two closed forms for \emph{the same}
sequence $b_n$; they are not a pair of forms for $b_n$ and $a_n$. (The denominators $a_n$ need
no harmonic weight: $a_n=\sum_k\binom nk^2\binom{n+k}{k}^2$ is Ap\'ery's own closed form.) That
the two weights differ while their $T_A$-sums agree is an instance of the non-uniqueness that
Proposition~\ref{prop:ptrank} quantifies one weight up: the fitting system has a large kernel of
genuine summation identities, so a representative is never determined by the identity it
satisfies.
\end{remark}

Proposition~\ref{prop:apery-a} matters here for a negative reason: it shows that Ap\'ery's own
base case is cell-wise in the \emph{purely harmonic} alphabet, so the $\zeta(3)$ precedent does
\emph{not} argue for binomial-reciprocal letters at weight $5$. The mechanism it does name ---
spend the weight on the non-polar $B$ slot --- is exactly what the weight-$5$ alphabet turns
out not to have room for, by exactly one functional (\S\ref{sec:wall}).

\begin{evidence}\label{ev:apery}
Proposition~\ref{prop:apery-b}: $\VERIFIED$ exact agreement with the standard $b_n$ for
$n=0,\dots,5$: $0,\,6,\,351/4,\,62531/36,\,11424695/288,\,35441662103/36000$.
Proposition~\ref{prop:apery-a}: exact $\QQ$ solve over $14$ levels plus one $u^3$ condition
(rank $15$), then $\VERIFIED$ on $24$ levels and $7$ primes, $0$ mismatches; and its right-hand
side agrees with the standard $b_n$ exactly for $n=0,\dots,14$.
\end{evidence}

%%%%%%%%%%%%%%%%%%%%%%%%%%%%%%%%%%%%%%%%%%%%%%%%%%%%%%%%%%%%%%%%%%%%%%%%%%%%%%%%
\section{The pole calculus and the depth conditions}\label{sec:depth}
%%%%%%%%%%%%%%%%%%%%%%%%%%%%%%%%%%%%%%%%%%%%%%%%%%%%%%%%%%%%%%%%%%%%%%%%%%%%%%%%

Throughout this section fix $p\ge5$ and $n\ge1$, put $L=\lfloor\log_p n\rfloor$, $M:=L+1$,
$P:=p^{M}$ (so $n<P$), fix a cell $0\le k,l\le n$, and write $u:=p^{-1}$. Set
\[
\vfive(n,k,l):=\wfive(n,k,l)-H^{(5)}_n,\qquad
d_5(n,k,l):=\max\bigl(0,-v_p\vfive(n,k,l)\bigr).
\]

\subsection{Level expansion}

\begin{lemma}[Level expansion]\label{lem:U}
For all $N\ge0$ and $m\ge1$,
$H^{(m)}_N=\sum_{e\ge0}u^{em}\,\Sigma^{(m)}_e(N)$ where
$\Sigma^{(m)}_e(N):=\sum_{j\le\lfloor N/p^e\rfloor,\ p\nmid j}j^{-m}\in\Zp$, and
$\Sigma^{(m)}_e(N)=0$ for $p^e>N$.
\end{lemma}

\begin{proof}
Partition $\{1,\dots,N\}$ by $e=v_p(i)$: $i=p^ej$ with $p\nmid j$, $j\le N/p^e$, and
$i^{-m}=p^{-em}j^{-m}$.
\end{proof}

Two consequences are used constantly: $\Sigma^{(m)}_e(N)=\Sigma^{(m)}_{e-1}(\lfloor N/p\rfloor)$
for $e\ge1$, and $v_p(H^{(m)}_N)\ge-m\lfloor\log_pN\rfloor$.

Since $n+k,n+l<2P$, $k+l<2P$, $n+k+l<3P$ and $n-k,k,l,n<P$, and since $p\ge5$ gives
$3P<pP=p^{M+1}$, \emph{every} letter has an expansion $X_m=\sum_{e=0}^{M}u^{em}\mathsf X_{m,e}$
with $\mathsf X_{m,e}\in\Zp$, and a weight-$5$ monomial has $u$-order at most $5M$. Hence
\begin{equation}\label{eq:Kexp}
\vfive(n,k,l)=\sum_{j=0}^{5M}K_j\,u^{j},\qquad K_j\in\Zp,
\qquad\text{so}\qquad d_5\le\max\{j:K_j\ne0\}.
\end{equation}

\subsection{Top-level residues and the pattern census}

\begin{definition}\label{def:pattern}
Put
\[
\begin{gathered}
\alpha:=\ii{n+k\ge P},\qquad \gamma:=\ii{n+l\ge P},\qquad
\varepsilon:=\lfloor(k+l)/P\rfloor\in\{0,1\},\\
\kappa:=\ii{n+k+l\ge(\varepsilon+1)P},\qquad \theta:=\varepsilon+1,
\end{gathered}
\]
($\theta$ set to $1$ when $\kappa=0$), and $s:=\alpha+\gamma+\kappa$. The tuple
$\pi=(\alpha,\gamma,\kappa,\theta)$ is the \emph{pole pattern} of the cell.
\end{definition}

\begin{lemma}[Top-level residues]\label{lem:R}
The $e=M$ coefficient of a letter is: $\alpha$ for $A_m(k)$, $\gamma$ for $A_m(l)$, $0$ for
every $B_m$ and for $N_m$, and $\kappa\,\theta^{-m}$ for $C_m$.
\end{lemma}

\begin{proof}
$k,l,n,n-k<P$ gives $\Sigma_M(\cdot)=0$ for those arguments. $\Sigma^{(m)}_M(n+k)$ is a sum
over $j\le\lfloor(n+k)/P\rfloor\le1$ with $p\nmid j$, i.e.\ equals $\alpha$. For $C_m$: with
$\Theta:=\lfloor(n+k+l)/P\rfloor\le2$ and $\varepsilon=\lfloor(k+l)/P\rfloor$, every $j\le2<p$,
so the coefficient is $\sum_{j=\varepsilon+1}^{\Theta}j^{-m}$; and $n<P$ forces
$\varepsilon\le\Theta\le\varepsilon+1$, so $\Theta-\varepsilon=\kappa$ and the coefficient is
$\kappa\theta^{-m}$.
\end{proof}

Note $\theta^{-m}$ is a $p$-unit for every $p\ge5$, so the whole calculus is $p$-independent.

\begin{lemma}[Kummer floor]\label{lem:K}
$v_p\binom{n+k}{n}\ge\alpha$, $v_p\binom{n+l}{n}\ge\gamma$, $v_p\binom{n+k+l}{n}\ge\kappa$;
hence
\begin{equation}\label{eq:vTfloor}
vT:=v_pT(n,k,l)\ \ge\ \alpha+\gamma+\kappa=s .
\end{equation}
\end{lemma}

\begin{proof}
Kummer: $v_p\binom{n+k}{n}$ is the number of carries in $n+k$. Both $n,k<P=p^M$, so a carry out
of position $M-1$ occurs iff $n+k\ge P$, i.e.\ iff $\alpha=1$; same for $\gamma$. For $\kappa$:
write $k+l=\varepsilon P+\rho$, $0\le\rho<P$; the carry out of position $M-1$ in $n+(k+l)$
occurs iff $n+\rho\ge P$, i.e.\ iff $n+k+l\ge(\varepsilon+1)P$.
\end{proof}

\begin{lemma}[The census, at every level]\label{lem:census}
The pattern $\pi$ takes exactly the seven values
\[
(0,0,0,1),\ (0,0,1,1),\ (1,0,1,1),\ (0,1,1,1),\ (1,1,0,1),\ (1,1,1,1),\ (1,1,1,2),
\]
with $s=0,1,2,2,2,3,3$ respectively.
\end{lemma}

\begin{proof}
(i) $\kappa=0\Rightarrow\alpha=\gamma$. Suppose $\alpha=1$, $\kappa=0$. If $\varepsilon=0$ then
$\kappa=\ii{n+k+l\ge P}=0$ gives $n+k+l<P$, contradicting $n+k\ge P$. So $\varepsilon=1$, i.e.
$k+l\ge P$; then $n+l\ge k+l\ge P$ (using $k\le n$), so $\gamma=1$. Symmetrically for
$\gamma=1$, $\kappa=0$. So the $\kappa=0$ patterns are $(0,0,0)$ and $(1,1,0)$.
(ii) $\theta=2$ (i.e.\ $k+l\ge P$) forces $n+k\ge k+l\ge P$ (using $l\le n$) and
$n+l\ge k+l\ge P$, i.e.\ $\alpha=\gamma=1$.
(iii) With $\kappa=1$, $\varepsilon=0$ all four $(\alpha,\gamma)$ occur. Total $2+4+1=7$.
\end{proof}

\begin{remark}
Lemma~\ref{lem:census} is level-free and does \emph{not} go through Lemma~\ref{lem:D}: at
$M=1$ it \emph{reproves} ``$a+b\ge p\Rightarrow vT\ge2$'' as a corollary rather than using it.
The absence of the patterns $(1,0,0,\cdot)$ and $(0,1,0,\cdot)$ is precisely the census form of
that statement.
\end{remark}

\subsection{The depth conditions and their consistency}

The Lemma-$F$ budget of \S\ref{sec:F} supplies precision $p^{2+\min(vT,2)}$, and the ledger
needs $p^{1+d_5}$; so the pattern-determined cap we must enforce is
\begin{equation}\label{eq:cap}
J(\pi):=1+\min(s,2)\qquad\text{(and $J=0$ at the trivial pattern $\pi=(0,0,0,1)$)}.
\end{equation}

\begin{definition}[(DEPTH)]\label{def:depth}
For each reachable pattern $\pi$ and each $j>J(\pi)$, require $K_j=0$ \emph{identically in the
$\Zp$-symbols} $\mathsf X_{m,e}$.
\end{definition}

Treating the symbols as independent indeterminates is a strengthening; it is $\QQ$-linear and
$p$-independent in the coefficients of $\wfive$, and --- as the next theorem shows --- it costs
nothing.

\begin{theorem}[Consistency of (DEPTH)]\label{thm:depthcons}
$\CERT$ Over two auxiliary primes $q=33554393$ and $q=33554467$, with $N=600$ fitting
equations:
\[
\begin{array}{ll}
\text{raw condition rows}&68\\
\rank(\text{conditions alone})&42\\
\rank(\text{fitting system alone})&313\\
\rank(\text{joint system})&324\\
\rank(\text{augmented joint system})&324\\
\end{array}
\]
so the joint system is \emph{consistent} at both primes, and the depth-conditioned family has
dimension $448-324=124$.
\end{theorem}

Thus of the $42$ independent depth conditions, $31$ are already implied by the decomposition
identity itself and only $11$ are new.

\begin{corollary}[Depth-minimal family]\label{cor:W5DEPTH}
$\PROVED$ modulo the certificate of Theorem~\textup{\ref{thm:depthcons}}. The decomposition
family contains a $124$-dimensional affine subfamily on which, for \emph{every} prime $p\ge5$
and every cell $0\le b,c\le a<p$,
\[
d_5(a,b,c)\ \le\ 1+\min\bigl(v_pT(a,b,c),2\bigr),
\]
provided the chosen representative's coefficients are $p$-integral --- which
Lemma~\textup{\ref{lem:CRT}} arranges for every $p\ge5$ at once.
\end{corollary}

\begin{evidence}\label{ev:depth}
$\VERIFIED$ $0$ failures. Three independent exact tests on each saved representative:
(V1) the exact-$\QQ$ ladder identity $P_n=\sum T\wfive$; (V2) the depth sweep
$d_5\le1+\min(vT,2)$ over \emph{all} cells; (V3) the cell-by-cell test
$v_p(\Tcal(b,c)-(Q_n/Q_a)T(a,b,c))\ge1+d_5(b,c)$ computed mod $p^{10}$ from the true fibre
tables. For $w_5^{\mathrm{allp}}$: (V1) $n\le34$, $0$; (V2) $\max d_5=3$, $0$ violations, min
slack $0$, up to $p=31$; (V3) $0$ failures at $p=5,7,11,13$. The comparison is the point:
\[
\begin{array}{lcccc}
\text{representative} & p=5 & p=7 & p=11 & p=13\\\hline
\text{130-term (pre-depth)} & 1/270 & 18/707 & 3385/5379 & 1285/10634\\
\text{106-term sparsest canonical} & 34/270 & 91/707 & 3503/5379 & \text{---}\\
\textbf{depth-minimal} & \mathbf 0 & \mathbf 0 & \mathbf 0 & \mathbf 0
\end{array}
\]
i.e.\ $0/16\,990$ cells fail, where the sparsest representatives failed thousands.
\emph{Sharpness:} min slack is $0$ at every prime --- the value $d_5=1+\min(vT,2)$ is attained.
So (DEPTH) sits exactly at the edge of what Lemma~\ref{lem:F} supplies.
\end{evidence}

\begin{remark}[sparsest is not depth-minimal]
Canonicalisation by sparsity and canonicalisation by depth are different, and it is depth that
is load-bearing. The sparsest representative found ($106$ terms) has $\max d_5=5$ and violates
the cap on the $vT\ge2$ cells; a $130$-term one has $\max d_5=4$. Neither works. Depth is not
visible in any sparsity or weight statistic --- and neither is the second load-bearing
criterion, $p$-integrality of the coefficients: a stray prime in a denominator destroys the
bound at that one prime and nowhere else.
\end{remark}

\subsection{Level-lifting: the conditions are level-independent}

This is the technical heart of the induction. Fix a pattern $\pi$. For a letter-multiset $S$
(letters counted with multiplicity) write $\wt(S)=\sum_{X\in S}\text{weight}(X)$ and
\begin{equation}\label{eq:Phifunctional}
\Phi_\pi(S):=\sum_{\Mcal\supseteq S}c_\Mcal\prod_{X\in\Mcal\setminus S}\rho(X;\pi),
\end{equation}
where $c_\Mcal$ is the $\wfive$-coefficient of the monomial $\Mcal$ and
$\rho(X;\pi)\in\{\alpha,\gamma,\kappa\theta^{-m},0\}$ is the top-level residue of
Lemma~\ref{lem:R}. These are $\QQ$-linear functionals of the $448$ coefficients, independent of
$p$, $n$ and $M$.

\begin{proposition}[LIFT]\label{prop:LIFT}
$\PROVED$ Let $J\ge0$ and suppose
\begin{equation}\label{eq:DEPTHpiJ}
\Phi_\pi(S)=0\quad\text{for every letter-multiset $S$ with }\wt(S)\le4-J .
\end{equation}
Then for \emph{every} $M\ge1$ and every cell $(n,k,l)$ with $\lfloor\log_pn\rfloor=M-1$ and
pattern $\pi$,
\[
d_5(n,k,l)\ \le\ 5(M-1)+J\ =\ 5L+J .
\]
Moreover at $M=1$ the conditions \eqref{eq:DEPTHpiJ} are \emph{exactly} the conditions
``$K_j=0$ identically in the symbols for $j>J$'' of Definition~\textup{\ref{def:depth}}.
\end{proposition}

\begin{proof}
Expand each letter by Lemma~\ref{lem:U} and collect. Writing $j=5M-\delta$ and
$f_i:=M-e_i\ge0$ (so $\sum_if_im_i=\delta$),
\[
K_{5M-\delta}=\sum_{\Mcal}c_\Mcal\sum_{\substack{f:\ \sum f_im_i=\delta\\ f_i\le M}}
\Bigl(\prod_{i:\,f_i\ge1}\mathsf X_{i,M-f_i}\Bigr)
\Bigl(\prod_{i:\,f_i=0}\rho(X_i;\pi)\Bigr).
\]
Regard the $\mathsf X_{i,e}$ ($0\le e\le M-1$, one indeterminate per (letter, level)) as
independent symbols --- this is the same strengthening already made at $M=1$. Two contributions
carry the same symbol monomial iff they have the same multiset
$\{(X_i,M-f_i):f_i\ge1\}$; that multiset determines the dropped multiset $S$ and the drop
profile $f|_S$. Summing over $\Mcal$ at fixed $(S,f|_S)$ produces $c(S,f)\,\Phi_\pi(S)$, where
$c(S,f)$ is the multinomial count of ways to distribute the copies of a repeated letter among
the prescribed levels --- a \emph{nonzero rational independent of $\Mcal$}: for a letter
occurring $N\ge q$ times in $\Mcal$ the count is $\binom Nq q!/\prod_uq_u!$, and its ratio to
the $M=1$ count $\binom Nq$ is $q!/\prod_uq_u!$, free of $N$. Hence
\[
\begin{gathered}
K_{5M-\delta}=0\ \text{identically in the symbols}\\
\iff\ \Phi_\pi(S)=0\ \text{for every $S$ realisable by an admissible $f$ with }\textstyle\sum_if_im_i=\delta .
\end{gathered}
\]
The minimal $\delta$ realising a given support $S$ is $\delta=\wt(S)$ (all $f_i=1$), admissible
for every $M\ge1$. Therefore
$\{K_j=0\ \forall j>5(M-1)+J\}\Longleftarrow\{\Phi_\pi(S)=0\ \forall S,\ \wt(S)\le4-J\}$,
because $j>5(M-1)+J\iff\delta=5M-j<5-J\iff\wt(S)\le\delta\le4-J$. With \eqref{eq:Kexp} this
gives $d_5\le5(M-1)+J$. At $M=1$, $\delta=5-(u\text{-order})$, so $u>J\iff\wt(S)\le4-J$: the two
condition sets coincide row for row.
\end{proof}

\begin{corollary}[(DEPTH-gen)]\label{cor:depthgen}
$\PROVED$ With $\wfive$ any point of the depth-conditioned $124$-dimensional family, for
\emph{every} prime $p\ge5$, \emph{every} $n\ge1$ and \emph{every} cell,
\[
d_5(n,k,l)\ \le\ 5L+1+\min(\alpha+\gamma+\kappa,2)\ \le\ 5L+1+\min\bigl(v_pT(n,k,l),2\bigr).
\]
\end{corollary}

\begin{proof}
The caps are $J(\pi)=1+\min(s,2)$ for $s\ge1$; for $\pi=(0,0,0,1)$ every top residue vanishes,
so $\Phi_\pi(S)=0$ automatically whenever $\Mcal\setminus S\ne\emptyset$, i.e.\ whenever
$\wt(S)\le4$, and the corollary holds there with $J=0$. Apply Proposition~\ref{prop:LIFT}. The
second inequality is Lemma~\ref{lem:K}.
\end{proof}

\begin{evidence}\label{ev:depthgen}
$\VERIFIED$ $0$ violations in $150\,955$ cells: $p\in\{5,7,11,13\}$, $n\le40/55/45/45$ (digit
levels $L=0,1,2$), all $(k,l)$. In the same sweep $v_pT\ge\alpha+\gamma+\kappa$ has $0$
counterexamples, and the bound is \emph{sharp} (min slack $0$ at every prime, attained already
at $L=0$).
\end{evidence}

\begin{corollary}[Theorem~\ref{thm:51intro}]\label{cor:51}
$\PROVED$ Assume \textup{(T1-top)} and \textup{(DEPTH)}. Then
$\ord_p(P_n)\ge-5\lfloor\log_pn\rfloor-1$ for every prime $p\ge5$ and every $n\ge1$;
equivalently $d_n^5\cdot\prod_{5\le p\le3n}p\cdot P_n\in\ZZ[1/6]$. For $p>3n$ all letter
arguments are $<p$, so $\ord_p(P_n)\ge0$ outright.
\end{corollary}

\begin{proof}
Cell by cell, $v_p(T\vfive)\ge vT-d_5\ge vT-5L-1-\min(vT,2)\ge-5L-1$ by
Corollary~\ref{cor:depthgen}; and $v_p(H^{(5)}_nQ_n)\ge-5L$ because $Q_n\in\ZZ$ and by
Lemma~\ref{lem:U}.
\end{proof}

\subsection{A hard negative: the deficit is irremovable}\label{sec:hardneg}

The cell-wise bound at $L=0$ is $v_p(T\vfive)\ge-1$, and the value $-1$ is \emph{attained}, at
$(n,k,l)=\bigl(\tfrac{p+1}{2},0,\tfrac{p-1}{2}\bigr)$ for every $p$ (so $p=5:(3,0,2)$,
$p=7:(4,0,3)$, $p=11:(6,0,5)$, $p=13:(7,0,6)$, \dots), where $vT=2$, $d_5=3$ and the pattern is
$(0,1,1,1)$. The cell-wise bound would follow from $d_5\le v_pT$, i.e.\ from strengthened
conditions. It does not exist.

\begin{theorem}\label{thm:DEPTHplus}
$\CERT$ at two auxiliary primes $q=33554393$, $q=33554467$, $N=600$:
\[
\small
\begin{array}{lccccc}
\text{system} & \text{caps }J(\pi) & \text{rows} & \rank(\text{cond}) & \rank(\text{joint}) & \text{consistent?}\\\hline
\text{(DEPTH)} & 1+\min(s,2) & 68 & 42 & 324 & \text{yes, nullity }124\\
(\mathrm{DEPTH}^{+}) & s & 239 & 123 & 342 & \textbf{NO}\\
(\mathrm{DEPTH}^{++}) & 2\text{ on the three }s=2\text{ patterns} & 149 & 81 & 342 & \textbf{NO}
\end{array}
\]
\end{theorem}

The third row is the sharp localisation: it is \emph{exactly the three $s=2$ patterns --- the
ones carrying the $-1$ deficit --- that cannot be tightened}; the $s\le1$ and $s=3$ patterns
are already at the cell-wise optimum. So within the $p$-independent linear framework the
deficit is irreducible, and the base case requires a genuine cancellation \emph{between} cells.
It is moreover global in both indices: $\VERIFIED$ ($p\le13$) the partial sums
$\sum_cT(a,b,c)\vfive(a,b,c)$ at fixed $b$ still have $v_p=-1$; only the full double sum is
integral. No row identity suffices.

\begin{theorem}[Alphabet enlargement does not help]\label{thm:rletter}
$\VERIFIED$ Let $R^{(a)}$ be the Ap\'ery-type letters and $Y_{a,b},V_{a,b}$ the nested interval
letters,
\[
\begin{gathered}
R^{(a)}(n,k)=\sum_{m\le k}\frac{(-1)^{m-1}}{m^a\binom nm\binom{n+m}{m}},\\
Y_{a,b}(n,k)=\!\!\sum_{k<m_2<m_1\le n+k}\!\! m_1^{-a}m_2^{-b},\qquad
V_{a,b}(n,k+l)=\!\!\sum_{k+l<m_2<m_1\le n+k+l}\!\! m_1^{-a}m_2^{-b}.
\end{gathered}
\]
Then $R^{(a)}$ has pole order $\le1$ at every weight $a=1,\dots,5$ with indicator exactly
$\alpha$, and $Y_{a,b}$ has pole order exactly $\max(a,b)<a+b$ with indicator $\alpha$
\textup{(}$V_{a,b}$ likewise with indicator $\kappa$\textup{)} --- i.e.\ both families supply
precisely the missing order. Nevertheless, with the full Ap\'ery alphabet $R^{(1..5)}$ at both
slots ($1210$ coefficients, $\rank(\mathrm{fit})$ $313\to960$, depth-conditioned value space
$\dim U$ $261\to641$) the strengthened system is \emph{still inconsistent}, at both the
$(\mathrm{DEPTH}^{++})$ and the $(\mathrm{DEPTH}^{+})$ cap, at $N=1300$ (row-saturated). The
same for $Y/V$.
\end{theorem}

The structural reason is worth recording, since it applies to \emph{any} enlargement.

\begin{proposition}[Block diagonality]\label{prop:blockdiag}
$\PROVED$ Let the alphabet be enlarged by new letters whose $u^0$ part is a \emph{new} $\Zp$
symbol (true for $R,D,Y,V,Z$). Split the basis into \emph{old} (monomials using no new letter)
and \emph{new}. Then the depth-condition matrix is block diagonal,
$C=\operatorname{diag}(C_{\mathrm{old}},C_{\mathrm{new}})$, because every $u$-expansion term of
a new monomial carries at least one new symbol in its symbol key and no old term does. Hence
$\rank C=\rank C_{\mathrm{old}}+\rank C_{\mathrm{new}}$, and with
$W_{\mathrm{old}}:=M_{\mathrm{old}}\ker C_{\mathrm{old}}$,
$W_{\mathrm{new}}:=M_{\mathrm{new}}\ker C_{\mathrm{new}}$ in value-sequence space,
\[
\text{the enlarged system is consistent}\iff P\in W_{\mathrm{old}}+W_{\mathrm{new}} .
\]
\end{proposition}

So enlargement does not relax the old conditions at all: the new letters must supply the
missing value-direction from scratch. And the obstruction is \emph{global}: the smallest prefix
of levels $n=1,\dots,\Lambda$ on which the strengthened system is already inconsistent is
$\Lambda=262$, one more than $\dim U=261$, exactly the generic threshold. There is no short
window of levels and hence no small-$n$ congruence carrying the obstruction.

\begin{remark}[the honest caveat]\label{rem:honest}
The depth conditions used throughout are the \emph{symbol-independent} form: for every
reachable pattern $\pi$ and every $u$-power $j>\mathrm{cap}(\pi)$, $K_j$ must vanish
identically in the $\Zp$ symbols. This is \emph{sufficient} for the cell-wise depth bound at
every prime simultaneously, and it is what makes the computation possible at all. It is
\emph{not necessary}: at a fixed prime the symbols range over the finitely many values the
cells supply. So Theorem~\ref{thm:rletter} refutes the $p$-independent symbol-independent form
of the enlargement route; the prime-by-prime form is a different, non-linear question which we
have not settled. Two remarks bound its plausibility: the symbol-independent conditions are
\emph{sharp} at the base cap (min slack $0$), so nothing is being wasted; and as $p\to\infty$
the cells fill out and the symbols become generic, so the two forms agree asymptotically. Any
gap lives at small primes.
\end{remark}

%%%%%%%%%%%%%%%%%%%%%%%%%%%%%%%%%%%%%%%%%%%%%%%%%%%%%%%%%%%%%%%%%%%%%%%%%%%%%%%%
\section{The fibre congruence}\label{sec:F}
%%%%%%%%%%%%%%%%%%%%%%%%%%%%%%%%%%%%%%%%%%%%%%%%%%%%%%%%%%%%%%%%%%%%%%%%%%%%%%%%

Throughout this section $p\ge5$, $n=ap+r$, $0\le r<p$, $k=bp+s$, $l=cp+t$ with $0\le b,c\le a$,
and
\[
\Tcal(b,c):=\sum_{s,t=0}^{p-1}T(n,\,bp+s,\,cp+t)\in\ZZ .
\]

\subsection{The block factorisation}

Write $T$ as a \emph{balanced factorial ratio} --- this is the form that makes everything work:
\begin{equation}\label{eq:Tfact}
T(n,k,l)=\frac{(n+k)!\,(n+l)!\,(n+k+l)!\,n!}
{k!^3\,l!^3\,(n-k)!^2\,(n-l)!^2\,(k+l)!},
\end{equation}
which is checked by expanding the five binomials; the argument sums balance,
$(n+k)+(n+l)+(n+k+l)+n=3k+3l+2(n-k)+2(n-l)+(k+l)=4n+2k+2l$. Call the $15$ factorial slots
$(n+k),(n+l),(n+k+l),n$ (numerator) and $k,k,k,l,l,l,(n-k),(n-k),(n-l),(n-l),(k+l)$
(denominator). For $m\ge0$ write $m=m_1p+m_0$ with $0\le m_0<p$, and
$(m!)_p:=\prod_{j\le m,\ p\nmid j}j$, so that $m!=p^{m_1}m_1!\,(m!)_p$.

\begin{lemma}[Wolstenholme block]\label{lem:W}
$\PROVED$ For $p\ge5$ and every $i\ge0$,
\[
B_i:=\prod_{u=1}^{p-1}(ip+u)\ \equiv\ (p-1)!\pmod{p^3}.
\]
\end{lemma}

\begin{proof}
$B_i=\sum_{j\ge0}(ip)^je_{p-1-j}(1,\dots,p-1)$ with $e_{p-1}=(p-1)!$,
$e_{p-2}=(p-1)!\,H_{p-1}$ and $e_{p-3}=(p-1)!\,(H_{p-1}^2-H^{(2)}_{p-1})/2$. Wolstenholme's
congruence ($p\ge5$) gives $H_{p-1}\equiv0\pmod{p^2}$ and $H^{(2)}_{p-1}\equiv0\pmod p$. Hence
the $j=1$ term lies in $p^3\Zp$, the $j=2$ term is
$p^2i^2(p-1)!(H_{p-1}^2-H^{(2)}_{p-1})/2\in p^3\Zp$, and the terms $j\ge3$ lie in $p^3\Zp$.
\end{proof}

Consequently, writing
\begin{equation}\label{eq:G}
G(m_1,m_0):=\Bigl[\prod_{i<m_1}B_i\Big/((p-1)!)^{m_1}\Bigr]\prod_{u=1}^{m_0}\Bigl(1+\frac{m_1p}{u}\Bigr),
\end{equation}
one has the \emph{exact} identity $(m!)_p=((p-1)!)^{m_1}\,m_0!\,G(m_1,m_0)$ together with
\[
G(m_1,m_0)=1+p\,m_1H_{m_0}+p^2m_1^2\frac{H_{m_0}^2-H^{(2)}_{m_0}}{2}+O(p^3),
\qquad G\in1+p\Zp
\]
(all $H_{m_0}$ here are $p$-integral because $m_0<p$).

\begin{lemma}[Exact fibre block factorisation]\label{lem:B}
$\PROVED$ Let $p\ge5$, $n=ap+r$ with $0\le r<p$, $k=bp+s$, $l=cp+t$ with $0\le b,c\le a$ and
$0\le s\le r$, $0\le t\le r$. Put
\[
e_1=\ii{r+s\ge p},\quad e_2=\ii{r+t\ge p},\quad e_3=\lfloor(r+s+t)/p\rfloor\in\{0,1,2\},\quad
e_4=\ii{s+t\ge p}.
\]
Then, \emph{exactly},
\begin{equation}\label{eq:LemB}
T(n,k,l)=T(a,b,c)\,T(r,s,t)\,\Pi\,\hat G,\qquad
\Pi=\frac{\binom{a+b+e_1}{e_1}\binom{a+c+e_2}{e_2}\binom{a+b+c+e_3}{e_3}}{\binom{b+c+e_4}{e_4}},
\qquad \hat G\in1+p\Zp,
\end{equation}
where $\hat G=\prod_{\mathrm{slots}}\bigl[G(m_1,m_0)/G(e_{\mathrm{slot}},m_0)\bigr]^{\pm1}$,
with $(m_1,m_0)$ the true base-$p$ digits of each slot and $e_{\mathrm{slot}}\in\{e_1,e_2,e_3,e_4\}$
on the four shifted slots, $0$ on the other eleven. The hypothesis $a<p$ is \emph{not} used.
\end{lemma}

\begin{proof}
Under $s\le r$, $t\le r$ the true digit pairs of the $15$ slots are
\[
\begin{gathered}
(n+k)\to(a+b+e_1,\,r+s-e_1p),\qquad (n+l)\to(a+c+e_2,\,r+t-e_2p),\\
(n+k+l)\to(a+b+c+e_3,\,r+s+t-e_3p),
\end{gathered}
\]
$n\to(a,r)$, $k\to(b,s)$, $l\to(c,t)$, $(n-k)\to(a-b,r-s)$, $(n-l)\to(a-c,r-t)$,
$(k+l)\to(b+c+e_4,\,s+t-e_4p)$. Substituting
$m!=p^{m_1}m_1!\,((p-1)!)^{m_1}m_0!\,G(m_1,m_0)$ in all $15$ slots gives
\[
T(n,k,l)=p^{E}((p-1)!)^{E}\Bigl[\prod(m_1!)^{\pm1}\Bigr]\Bigl[\prod(m_0!)^{\pm1}\Bigr]
\Bigl[\prod G(m_1,m_0)^{\pm1}\Bigr],
\]
with $E=\sum\pm m_1=e_1+e_2+e_3-e_4$.
For each shifted slot write $m_1!=(m_1-e)!\cdot m_1!/(m_1-e)!$ and
$m_0!=(m_0+ep)!\big/\bigl(p^ee!\,((p-1)!)^e\,G(e,m_0)\bigr)$ (the second by the same
substitution applied to $m_0+ep$, whose digits are $(e,m_0)$). The unshifted level-$a$ digits
are exactly $a+b,\,a+c,\,a+b+c,\,a;\,b,b,b,\,c,c,c,\,a-b,a-b,\,a-c,a-c,\,b+c$, i.e.
$\prod(m_1-e)!^{\pm1}=T(a,b,c)$; and the unshifted level-$r$ digits are
$r+s,\,r+t,\,r+s+t,\,r;\,s,s,s,\,t,t,t,\,r-s,r-s,\,r-t,r-t,\,s+t$, i.e.
$\prod(m_0+ep)!^{\pm1}=T(r,s,t)$. Collecting, the powers $p^{\pm e}$ and $((p-1)!)^{\pm e}$
cancel $p^E((p-1)!)^E$ exactly, $\prod[m_1!/((m_1-e)!e!)]^{\pm1}=\Pi$, and the leftover $G$'s
are $\hat G$. Finally $\hat G\in1+p\Zp$ because each $G(m_1,m_0)\in1+p\Zp$ by
Lemma~\ref{lem:W}, for every $m_1\ge0$.
\end{proof}

The admissible carry patterns involve only $r,s,t$ and are
\begin{equation}\label{eq:epatterns}
(e_1,e_2,e_3,e_4)\in\{(0,0,0,0),(0,0,1,0),(1,0,1,0),(0,1,1,0),(1,1,1,0),(1,1,1,1),(1,1,2,1)\}
\end{equation}
($e_4=1\Rightarrow e_1=e_2=1$ since $t\le r$ gives $r+s\ge s+t\ge p$; $e_1=1\Rightarrow e_3\ge1$;
$e_3=2\Rightarrow e_4=1$).

\begin{evidence}\label{ev:LemB}
$\VERIFIED$ $8\,247\,294$ in-regime cells, $0$ failures ($p\in\{5,7,11\}$, all $a\le3p+1$, all
$r$, all $0\le b,c\le a$, all $s,t\le r$): $\hat G$ recomputed as
$T(n,k,l)/(T(a,b,c)T(r,s,t)\Pi)$ over $\QQ$ and tested for $\hat G-1\in p\Zp$.
\end{evidence}

\subsection{The residue identity}

\begin{lemma}[$\Phi$]\label{lem:Phi}
$\PROVED$ For all integers $r\ge0$, $t\ge0$,
\[
\sum_{s=0}^{r}T(r,s,t)\,\Phi_b(s,t)=0,\qquad
\Phi_b(s,t):=H_{r+s}+H_{r+s+t}-3H_s+2H_{r-s}-H_{s+t}=A_1(s)+2B_1(s)+C_1
\]
in the level-$r$ letters. By the $k\leftrightarrow l$ symmetry of $T$, also
$\sum_{t=0}^{r}T(r,s,t)\Phi_c(s,t)=0$ for each $s$, with $\Phi_c(s,t)=A_1(t)+2B_1(t)+C_1$.
\end{lemma}

\begin{proof}
$\Phi_b=\partial_s\log T(r,s,t)$ in the $\Gamma$-continuation, and the identity is a residue
sum. Concretely, put
\[
G(x):=\prod_{i=1}^{r}(x+i)\prod_{i=1}^{r}(x+t+i)\Big/\prod_{j=0}^{r}(x-j)^2 .
\]
$G$ is rational with $\deg(\mathrm{den})-\deg(\mathrm{num})=(2r+2)-2r=2$, so $G(x)=O(x^{-2})$
and the sum of all its residues vanishes. Its poles are exactly the double poles at
$x=0,1,\dots,r$ (the numerator zeros are at $x=-1,\dots,-r$ and $x=-t-1,\dots,-t-r$, all
negative, so no cancellation). At $x=s$ write $G=N/((x-s)^2E)$ with $E(x)=\prod_{j\ne s}(x-j)^2$;
then
\[
\begin{aligned}
\Res_{x=s}G&=\frac NE(s)\Bigl[\frac{N'}{N}(s)-\frac{E'}{E}(s)\Bigr]\\
&=\frac NE(s)\bigl[(H_{r+s}-H_s)+(H_{r+s+t}-H_{s+t})-2H_s+2H_{r-s}\bigr]
=\frac NE(s)\,\Phi_b(s,t),
\end{aligned}
\]
and $\frac NE(s)=\frac{(r+s)!(r+s+t)!}{s!^3(s+t)!(r-s)!^2}=T(r,s,t)/K_t$ with
$K_t=\frac{(r+t)!\,r!}{t!^3(r-t)!^2}$ independent of $s$. Summing the residues gives the claim.
\end{proof}

\begin{remark}
$\Phi_b,\Phi_c$ are exactly the level-$r$ weight-$1$ letter combinations $A_1+2B_1+C_1$ ---
the same alphabet as $\wthree$ and $\wfive$. Lemma~\ref{lem:Phi} is the weight-$1$ shadow of
the same residue mechanism that produced the closed forms of \S\ref{sec:apery}; the weight-$2$
level is Lemma~\ref{lem:Phi2}. $\VERIFIED$ exact, $0$ exceptions, $r\le11$, all $t$, plus the
aggregated $\sum_{s,t}$ version for $r\le10$ over $\QQ$.
\end{remark}

\subsection{The two carry lemmas and Lemma \texorpdfstring{$F$}{F}}

Write $vT:=v_pT(a,b,c)=\alpha+\gamma+\kappa$ with $\alpha=\ii{a+b\ge p}$,
$\gamma=\ii{a+c\ge p}$, $\kappa=v_p\binom{a+b+c}{a}\in\{0,1\}$, each a single Kummer carry
because $a,b,c<p$; so $vT\le3$.

\begin{lemma}[Off-fibre-regime vanishing]\label{lem:F2}
$\PROVED$ If $s>r$ or $t>r$ then $v_p\bigl(T(n,bp+s,cp+t)\bigr)\ge2+\min(vT,2)$.
\end{lemma}

\begin{proof}
By $k\leftrightarrow l$ symmetry assume $s>r$. If $b=a$ then $k=ap+s>ap+r=n$ and $T=0$; so
$b<a$. The subtraction $n-k$ borrows at position $0$, so $v_p\binom nk\ge1$ and the
\emph{squared} factor contributes $2$. It remains to produce $\min(vT,2)$ further powers.
If $\alpha=1$, the position-$1$ digit sum of $n+k$ is $a+b+\ii{r+s\ge p}\ge a+b\ge p$, a carry,
so $v_p\binom{n+k}{n}\ge1$; identically if $\gamma=1$. If $\kappa=1$, write $k+l=\beta p+\sigma$
with $\beta=b+c+e_4$, $\sigma=s+t-e_4p$: if $b+c\ge p$ then $\kappa=1$ means $a+b+c\ge2p$ and
the position-$1$ sum of $n+(k+l)$ is $a+(\beta-p)+\text{carry}\ge a+b+c-p\ge p$, a carry; if
$b+c<p$ and $\beta<p$, the position-$1$ sum is $a+\beta+\text{carry}\ge a+b+c\ge p$, a carry.
The only gap is $\beta=p$, i.e.\ $b+c=p-1$ and $e_4=1$.
Now count. If $vT\le1$, at most one of $\alpha,\gamma,\kappa$ is $1$ and the above supplies it
(in the gap case $\kappa=1$, $\alpha=\gamma=0$ one gets $b+c=p-1$, $b,c<p-a$, hence
$p-1<2(p-a)$ and $p-1\le2a$, forcing $a=(p-1)/2$ and $b=c=a$, contradicting $b<a$). If
$vT\ge2$, at least two of $\alpha,\gamma,\kappa$ equal $1$; if they are $\alpha,\gamma$ we are
done. Otherwise the gap $b+c=p-1$, $e_4=1$ may occur: if $\alpha=1,\gamma=0$ then $\gamma=0$
gives $c\le p-a-1$, hence $b=p-1-c\ge a$, contradicting $b<a$ --- vacuous. If
$\gamma=1,\alpha=0$ then symmetrically $c=a$; if $t>r$ then $\binom nl^2$ gives $2$ more and we
are done, so $t\le r$, and $e_4=1$ gives $s\ge p-t\ge p-r$, i.e.\ $r+s\ge p$; then $n+k$
carries at position $0$ and (since $a+b+1=a+(p-1-a)+1=p$) at position $1$, so
$v_p\binom{n+k}{n}\ge2$ and the total is $\ge2+1+2=5\ge4$. If $\alpha=\gamma=1$ then
$v_p\binom{n+k}{n}+v_p\binom{n+l}{n}\ge2$.
\end{proof}

\begin{lemma}[First-order fibre expansion]\label{lem:F1}
$\PROVED$ For $0\le s\le r$, $0\le t\le r$,
\[
T(n,bp+s,cp+t)\equiv T(a,b,c)T(r,s,t)\bigl(1+p[a\Phi_a+b\Phi_b+c\Phi_c]\bigr)\pmod{p^{2+vT}},
\]
\[
\Phi_a(s,t)=H_{r+s}+H_{r+t}+H_{r+s+t}+H_r-2H_{r-s}-2H_{r-t},
\]
except in the single carry pattern $(e_1,e_2,e_3,e_4)=(1,1,1,1)$ with $\alpha=\gamma=1$, where
the congruence holds mod $p^{1+vT}$ (and then $vT=3$, so $1+vT=4=2+\min(vT,2)$). In all cases
the congruence holds mod $p^{2+\min(vT,2)}$.
\end{lemma}

\begin{proof}
Set $\Delta_1:=a\Phi_a+b\Phi_b+c\Phi_c$; expanding the definitions,
$\Delta_1=\sum_{\mathrm{slots}}\pm\mu_1H_{\mu_0}$ where $(\mu_1,\mu_0)$ are the \emph{na\"ive}
(level-$a$, level-$r$) digits, i.e.\ $\mu_0=m_0+e_{\mathrm{slot}}p$. Each $H_{\mu_0}$ has
$\mu_0<3p$, hence $H_{\mu_0}=(\Zp)+p^{-1}H_{\lfloor\mu_0/p\rfloor}$ with
$\lfloor\mu_0/p\rfloor\le2<p$, so $p\Delta_1\in\Zp$; moreover
$H_{\mu_0}-H_{m_0}=p^{-1}H_{e_{\mathrm{slot}}}+(\Zp)$, whence
$1+p\Delta_1=1+pX+\sum_{\text{shifted}}\pm\mu_1H_{e_{\mathrm{slot}}}+p(\Zp)$ with
$X=\sum_{\mathrm{slots}}\pm\mu_1H_{m_0}$. By Lemma~\ref{lem:B},
$T(n,k,l)/(T(a,b,c)T(r,s,t))=\Pi\hat G=\Pi(1+pX+O(p^2))$, so the difference equals
$T(a,b,c)T(r,s,t)\bigl[\Pi-1-\sum\pm\mu_1H_e+O(p)\bigr]$. Let $J:=v_pT(r,s,t)$ and $K$ the
valuation of the bracket; the conclusion is $v_p(\mathrm{diff})\ge vT+J+K$. Running the seven
patterns \eqref{eq:epatterns}:
\[
\begin{array}{lcccc}
(e_1,e_2,e_3,e_4) & \Pi & 1+\sum\pm\mu_1H_e & K\ge & J\ge\\\hline
(0,0,0,0) & 1 & 1 & 2 & 0\\
(0,0,1,0) & a{+}b{+}c{+}1 & 1+(a{+}b{+}c)\ \text{---\ equal} & 1 & 1\\
(1,0,1,0),(0,1,1,0) & (a{+}b{+}1)(a{+}b{+}c{+}1) & 1+(a{+}b)+(a{+}b{+}c) & 0 & 2\\
(1,1,1,0) & \text{3 factors} & \text{---} & 0 & 3\\
(1,1,1,1) & \frac{(a+b+1)(a+c+1)(a+b+c+1)}{b+c+1} & \text{---} & -1 & 2\\
(1,1,2,1) & \text{with }\binom{a+b+c+2}{2} & \text{---} & -1 & 3
\end{array}
\]
$K\ge0$ whenever $e_4=0$ (then $\Pi\in\ZZ$), and $K\ge-1$ always since $b+c+1\le2a+1<2p$. So
$J+K\ge2$, i.e.\ $v_p(\mathrm{diff})\ge vT+2$, in every pattern except $(1,1,1,1)$, and there
$K=-1$ forces $v_p\binom{b+c+1}{1}=1$, i.e.\ $b+c=p-1$, with $v_p(a+c+1)=v_p(a+b+1)=0$. But
$b+c=p-1$ gives $\kappa=1$, and: $\alpha=\gamma=0$ forces (as in Lemma~\ref{lem:F2})
$a=(p-1)/2$, $b=c=a$, so $a+b+1=p$, a contradiction; exactly one of $\alpha,\gamma$ equal to
$1$ forces $b=a$ resp.\ $c=a$, hence $a+c+1=p$ resp.\ $a+b+1=p$, again a contradiction. So
$K=-1$ implies $\alpha=\gamma=1$, hence $vT=2+\kappa=3$ and
$v_p(\mathrm{diff})\ge3+2-1=4=2+\min(vT,2)$.
\end{proof}

\begin{theorem}[Lemma $F$]\label{lem:F}
$\PROVED$ Let $p\ge5$, $n=ap+r$, $1\le a<p$, $0\le r<p$, $p\nmid Q_a$ and $0\le b,c\le a$. Then
\[
\Tcal(b,c)\ \equiv\ \frac{Q_n}{Q_a}\,T(a,b,c)\pmod{p^{\,2+\min(v_pT(a,b,c),\,2)}} .
\]
\end{theorem}

\begin{proof}
Split the fibre at $s\le r$, $t\le r$. By Lemma~\ref{lem:F2} the complementary terms are
$\equiv0$ mod $p^{2+\min(vT,2)}$, and there $T(r,s,t)=0$. By Lemma~\ref{lem:F1}, for
$s,t\le r$,
$T(n,bp+s,cp+t)\equiv T(a,b,c)T(r,s,t)(1+p[a\Phi_a+b\Phi_b+c\Phi_c])$ mod $p^{2+\min(vT,2)}$.
Every product $T(r,s,t)\Phi_x(s,t)$ is $p$-\emph{integral}: a pole of $\Phi_x$ needs $r+s\ge p$
(so $\binom{r+s}{r}$ carries), $r+t\ge p$, $s+t\ge p$ (so $r+s\ge s+t\ge p$), or $r+s+t\ge p$
with $r+s,r+t<p$ (so $s+t<p$ by Lemma~\ref{lem:L3}, hence $\binom{r+s+t}{r}$ carries); in each
case $v_pT(r,s,t)\ge1$, and every pole has order exactly $1$ because all arguments are $<3p$.
Summing over $0\le s,t\le p-1$ and using $T(r,s,t)=0$ for $s>r$ or $t>r$,
\[
\Tcal(b,c)\equiv T(a,b,c)\Bigl[\sum_{s,t=0}^{r}T(r,s,t)+p\,a\,\Psi_a
+p\,b\!\sum_{s,t}\!T\Phi_b+p\,c\!\sum_{s,t}\!T\Phi_c\Bigr]\pmod{p^{2+\min(vT,2)}},
\]
with $\Psi_a:=\sum_{s,t=0}^{r}T(r,s,t)\Phi_a(s,t)\in\Zp$. Now
$\sum_{s,t=0}^{r}T(r,s,t)=Q_r$ by \eqref{eq:BZQ} at $n=r$, and by Lemma~\ref{lem:Phi} the last
two sums vanish. \emph{The whole $(b,c)$-dependence of the first-order term therefore cancels},
leaving $\Tcal(b,c)\equiv\mu\,T(a,b,c)$ mod $p^{2+\min(vT(b,c),2)}$ with
$\mu:=Q_r+p\,a\,\Psi_a\in\Zp$ independent of $(b,c)$. Summing over $0\le b,c\le a$ and using
$\sum_{b,c}\Tcal(b,c)=Q_n$, $\sum_{b,c}T(a,b,c)=Q_a$ and $2+\min(vT,2)\ge2$ gives
$Q_n\equiv\mu Q_a\pmod{p^2}$, hence $\Lambda:=Q_n/Q_a\equiv\mu\pmod{p^2}$ as $p\nmid Q_a$.
Finally $\Tcal(b,c)-\Lambda T(a,b,c)=T(a,b,c)(\mu-\Lambda)+[\Tcal(b,c)-\mu T(a,b,c)]$, whose
two terms have $v_p\ge vT+2$ and $\ge2+\min(vT,2)$ respectively.
\end{proof}

\begin{corollary}[A $\bmod\ p^2$ supercongruence for the bottom row]\label{cor:Qsuper}
$\PROVED$ For $p\ge5$, $1\le a<p$, $0\le r<p$ with $p\nmid Q_a$,
\[
Q_{ap+r}\equiv Q_a\bigl(Q_r+p\,a\,\Psi_a\bigr)\pmod{p^2},\qquad
\Psi_a=\sum_{s,t=0}^{r}T(r,s,t)\,\Phi_a(s,t),
\]
with $\Phi_a=A_1(s)+A_1(t)+C_1-2B_1(s)-2B_1(t)-H_s-H_t+H_{s+t}+H_r$ in the level-$r$ letters.
\end{corollary}

This is a genuine $\bmod\ p^2$ refinement of Theorem~\ref{thm:Aintro}, and $\Psi_a$ is the
level-$r$ weight-$1$ harmonic functional of the Brown--Zudilin summand. $\VERIFIED$ $0$
mismatches at $p=5,7,11,13$, all $a,r$ with $p\nmid Q_a$, $n\le360$.

\begin{evidence}[sharpness of Lemma $F$]\label{ev:F}
$\VERIFIED$ $0$ failures, and \emph{sharp}: over $20\,998$ cells at $p=5,7,11,13$ the pairs
$(vT,v_p(\mathrm{diff}))$ realised are exactly $(0,2),(1,3),(2,4),(3,4)$, i.e.\ the bound
$2+\min(vT,2)$ is attained in every stratum (min slack $0$). Consequently ``prove Lemma $F$ one
order deeper'' is \emph{impossible}: the congruence mod $p^{3+\min(vT,2)}$ is false. Lemmas
\ref{lem:F2} and \ref{lem:F1} are each $\VERIFIED$ at $p=5,7$ over all $(a,b,c,r,s,t)$,
stratified by carry pattern ($26$ strata; min slack $\ge0$ in every one and $=0$ in seven).
\end{evidence}

\begin{remark}[what Lemma $F$ did \emph{not} need]
Not needed: Jacobsthal's $\binom{ap}{bp}\equiv\binom ab\pmod{p^3}$, and Granville's general
prime-power binomial congruences. The only external input is Wolstenholme, used once, in
Lemma~\ref{lem:W} --- and that is exactly why $p\ge5$ is the right hypothesis. The real content
is (i) the exact block factorisation, where $T$'s being a \emph{balanced} factorial ratio is
used, and (ii) Lemma~\ref{lem:Phi}. The proof is weight-agnostic: Lemma~\ref{lem:F} is a
statement about $T$ alone.
\end{remark}

\subsection{Two carry inequalities}

The next three sections use two elementary base-$p$ facts repeatedly; we isolate them once.
Write $\car(x,y,\epsilon)$ for the number of carries in the base-$p$ addition $x+y+\epsilon$,
with $\epsilon\in\{0,1\}$ the carry into position $0$; write $\mathrm s(N)$ for the base-$p$
digit sum of $N$, and $z(N)$ for the number of trailing digits of $N$ equal to $p-1$.

\begin{lemma}[Carry inequalities]\label{lem:carry}
$\PROVED$ With $n=ap+r$, $k=bp+s$, $l=cp+t$ and $e_1=\ii{r+s\ge p}$, $e_2=\ii{r+t\ge p}$ as in
Lemma~\textup{\ref{lem:B}}:
\begin{enumerate}[leftmargin=2.4em,label=\textup{(C\arabic*)}]
\item $\car(x,y,1)=\car(x,y,0)+z(x+y)\ \ge\ \car(x,y,0)$ for all $x,y\ge0$;
\item $v_p\binom{n+k}{n}=e_1+\car(a,b,e_1)\ge v_p\binom{a+b}{a}$, and likewise for
  $(n+l,n)$ with $e_2$; and
  $v_p\binom nk=\ii{s>r}+\car(b,\,a-b-\ii{s>r},\,\ii{s>r})\ge v_p\binom ab+\ii{s>r}$,
  and likewise for $(n,l)$.
\end{enumerate}
\end{lemma}

\begin{proof}
(C1): $\car(x,y,\epsilon)=(\mathrm s(x)+\mathrm s(y)+\epsilon-\mathrm s(x+y+\epsilon))/(p-1)$,
and $\mathrm s(N+1)=\mathrm s(N)+1-(p-1)z(N)$.

(C2): by Kummer's theorem $v_p\binom{n+k}{n}$ counts the carries of $n+k$ in base $p$; splitting
that addition at position $0$ contributes $e_1$ there and $\car(a,b,e_1)$ above, and (C1) gives
$\car(a,b,e_1)\ge\car(a,b,0)=v_p\binom{a+b}{a}$. For the second, split the subtraction
$n-k$ at position $0$: it borrows exactly when $s>r$, the higher-position carries of
$k+(n-k)$ sum to $a$, and $\mathrm s(a-b-1)+1\ge\mathrm s(a-b)$.
\end{proof}

\subsection{The multi-digit weakening}

The induction consumes one order less precision, and at that precision the statement holds for
\emph{all} $a$, with a shorter proof and no auxiliary hypotheses at all.

\begin{lemma}[Lemma $F$-gen]\label{lem:Fgen}
$\PROVED$ Let $p\ge5$, $a\ge1$, $0\le r<p$, $n=ap+r$, $0\le b,c\le a$. Let $\alpha,\gamma,\kappa$
be the level-$a$ indicators of Definition~\textup{\ref{def:pattern}} evaluated at $(a,b,c)$ with
$P_a=p^{\lfloor\log_pa\rfloor+1}$, and $s_a:=\alpha+\gamma+\kappa$. Then
\[
v_p\bigl(\Tcal(b,c)-Q_r\,T(a,b,c)\bigr)\ \ge\ 1+\min(s_a,2)\ =:\ G\ (\le3).
\]
\end{lemma}

Three features distinguish this from Theorem~\ref{lem:F}: it holds for \emph{all} $a$, not only
$a<p$; the comparison constant is $Q_r$ itself, with no $\mu=Q_r+pa\Psi_a$ correction and no
$\Lambda=Q_n/Q_a$, hence \emph{no $p\nmid Q_a$ hypothesis} and no exceptional-prime
compensation; and it is one order weaker, which is what makes the general-$a$ proof short. It
is exactly the precision the induction consumes.

%% FIXED [m-6]: (C1) and (C2) were stated and proved inside this proof body and then invoked
%% FIXED twice more, two sections later. They are now the displayed Lemma \ref{lem:carry}.
\begin{proof}
Both parts use the carry inequalities of Lemma~\ref{lem:carry}.

\emph{(A) Off-regime terms $s>r$ or $t>r$.} By symmetry assume $s>r$. If $b=a$ then $k>n$ and
$T=0$. So $b<a$; the subtraction $n-k$ borrows at position $0$, so $v_p\binom nk\ge1$ and the
squared factor gives $v_pT\ge2$. That settles $G\le2$, i.e.\ $s_a\le1$. If $s_a\ge2$ then at
least two of $\alpha,\gamma,\kappa$ are $1$; $\alpha=\gamma=0$ would force $s_a=\kappa\le1$, so
$\alpha=1$ or $\gamma=1$, and then Lemma~\ref{lem:carry}(C2) gives one further power: $v_pT\ge3=G$.

\emph{(B) In-regime terms $s\le r$, $t\le r$.} Since $Q_r=\sum_{s,t\le r}T(r,s,t)$ exactly,
\eqref{eq:LemB} gives
$\sum_{s,t\le r}T(n,k,l)-Q_rT(a,b,c)=T(a,b,c)\sum_{s,t\le r}T(r,s,t)[\Pi\hat G-1]$, so it
suffices termwise that $vT+J+v_p(\Pi\hat G-1)\ge G$ with $J:=v_pT(r,s,t)$. Note
$v_p(\Pi\hat G-1)\ge\min(v_p(\Pi-1),\,v_p(\Pi)+1)\ge\min(v_p(\Pi),0)$, and $G\le3$,
$vT\ge s_a$. Running the seven patterns \eqref{eq:epatterns}:
\[
\begin{array}{lcccl}
(e_1,e_2,e_3,e_4) & \Pi & J\ge & v_p(\Pi\hat G-1)\ge & \text{total}\ge\\\hline
(0,0,0,0) & 1 & 0 & 1 & s_a+1\ \checkmark\\
(0,0,1,0) & a{+}b{+}c{+}1\in\ZZ & 1 & 0 & s_a+1\ \checkmark\\
(1,0,1,0),(0,1,1,0) & \in\ZZ & 2 & 0 & s_a+2\ \checkmark\\
(1,1,1,0) & \in\ZZ & 3 & 0 & 3\ \checkmark\\
(1,1,1,1) & \frac{(a+b+1)(a+c+1)(a+b+c+1)}{b+c+1} & 2 & \text{see below} & \checkmark\\
(1,1,2,1) & \frac{(a+b+1)(a+c+1)\binom{a+b+c+2}{2}}{b+c+1} & 3 & \text{see below} & \checkmark
\end{array}
\]
(the $J$ column: $e_1=1\Rightarrow\binom{r+s}{r}$ carries; $e_2=1\Rightarrow\binom{r+t}{r}$
carries; if $e_4=0$ and $e_3\ge1$ then $\binom{r+s+t}{r}$ carries at position $0$; if $e_4=1$
and $e_3=2$ then $r+\sigma\ge p$ with $\sigma=s+t-p$, so $\binom{r+s+t}{r}$ carries).

For the two $e_4=1$ patterns put $m:=v_p(b+c+1)$ and $\tau:=v_p(a)$. Then $a+b+c+1=a+(b+c+1)$
has $v_p\ge\min(\tau,m)$, so $v_p(\Pi)\ge\min(\tau,m)-m$ and therefore
$v_p(\Pi\hat G-1)\ge\min(\tau,m)-m$ (which is $\le0$, so the bound is valid whether or not
$v_p(\Pi)\ge0$).

\emph{Case $\tau\ge m$.} Then $v_p(\Pi\hat G-1)\ge0$ and the totals are $vT+2\ge2$ resp.
$vT+3\ge3$; since $G\le3$ and, for $(1,1,1,1)$, $vT+2\ge s_a+2\ge\min(s_a,2)+1=G$.

\emph{Case $\tau<m$ (so $m\ge1$).} Then $b+c\equiv-1\pmod{p^m}$, i.e.\ the digits of $b+c$ in
positions $0,\dots,m-1$ are all $p-1$, while $a$ has digits $0$ in positions $0,\dots,\tau-1$
and $a_\tau\ge1$. In the addition $a+(b+c)$: no carry at positions $0,\dots,\tau-1$
($0+(p-1)=p-1$), a carry at position $\tau$ ($a_\tau+(p-1)\ge p$), and a carry at each of
$\tau+1,\dots,m-1$ ($a_i+(p-1)+1\ge p$). Hence
\[
E:=v_p\tbinom{a+b+c}{a}\ge m-\tau,\qquad\text{so}\qquad vT\ge(m-\tau)+A+A'+2D+2D',
\]
with $A=v_p\binom{a+b}{a}\ge\alpha$, $A'=v_p\binom{a+c}{a}\ge\gamma$, $D=v_p\binom ab$,
$D'=v_p\binom ac$. For $(1,1,1,1)$ the total is $\ge vT+2+(\tau-m)\ge2+A+A'+2D+2D'$: if
$s_a\le1$ then $G\le2$; if $s_a\ge2$ then as in (A) $\alpha=1$ or $\gamma=1$, so $A+A'\ge1$ and
the total is $\ge3=G$. For $(1,1,2,1)$ the total is $\ge vT+3+(\tau-m)\ge3\ge G$.
\end{proof}

\begin{remark}
Lemma~\ref{lem:Phi} is \emph{not} used. The single-digit Lemma~\ref{lem:F} needs it because it
aims at relative precision $p^2$, so its first-order term must be $(b,c)$-independent; at
relative precision $p^1$ the first-order term is simply absorbed, and the whole first-order
analysis --- Lemmas \ref{lem:F1}, \ref{lem:F2}, \ref{lem:Phi} and the $\Psi_a$
supercongruence --- disappears. That is the structural reason the multi-digit statement is
\emph{easier} than the single-digit one.
\end{remark}

\begin{evidence}\label{ev:Fgen}
$\VERIFIED$ $0$ failures in $111\,963$ cells: $p=5,7,11,13$, multi-digit $a\le30/22/15/13$, all
$r$, all $(b,c)$; sharp (min slack $0$).
\end{evidence}

%%%%%%%%%%%%%%%%%%%%%%%%%%%%%%%%%%%%%%%%%%%%%%%%%%%%%%%%%%%%%%%%%%%%%%%%%%%%%%%%
\section{The digit-scaled induction}\label{sec:IND}
%%%%%%%%%%%%%%%%%%%%%%%%%%%%%%%%%%%%%%%%%%%%%%%%%%%%%%%%%%%%%%%%%%%%%%%%%%%%%%%%

\subsection{The inductive statement}

\begin{definition}[(IND)]\label{def:IND}
Let $p\ge5$. For $L\ge0$ let $\mathcal I(L)$ be the assertion
\[
\text{for every }m\ge1\text{ with }\lfloor\log_pm\rfloor\le L:\quad
v_p(W_m)\ \ge\ -5\lfloor\log_pm\rfloor,\qquad W_m=P_m-H^{(5)}_mQ_m .
\]
\end{definition}

Note $v_p(H^{(5)}_mQ_m)\ge-5\lfloor\log_pm\rfloor$ always ($Q_m\in\ZZ$, Lemma~\ref{lem:U}), so
$\mathcal I(L)$ is \emph{equivalent} to $v_p(P_m)\ge-5\lfloor\log_pm\rfloor$ on the same range.

\begin{remark}[why this form]
The na\"ive per-digit ratio statement $p^5P_n/Q_n\equiv P_a/Q_a\pmod p$ fails multi-digit: the
depth of the single-step ratio is $5-2s$, negative for $s\ge3$. (IND) is scaled: it never
divides by $Q$, never renormalises, and its per-level target is an inequality on $v_p(W)$, not
a congruence.
\end{remark}

\subsection{The step and its ledger}

Let $L\ge1$, $\lfloor\log_pn\rfloor=L$, $n=ap+r$ with $0\le r<p$, $a=\lfloor n/p\rfloor$, so
$\lfloor\log_pa\rfloor=L-1$. Group the cells of $W_n$ by
$(b,c)=(\lfloor k/p\rfloor,\lfloor l/p\rfloor)$, $0\le b,c\le a$, and split
\begin{equation}\label{eq:split}
p^5W_n-Q_rW_a=(\mathrm I)+(\mathrm{II}),\qquad
\begin{aligned}
(\mathrm{II})&:=\sum_{b,c}\vfive(a,b,c)\bigl[\Tcal(b,c)-Q_rT(a,b,c)\bigr],\\
(\mathrm I)&:=\sum_{k,l}T(n,k,l)\,\Ecal(n,k,l),\ \
\Ecal:=p^5\vfive(n,k,l)-\vfive(a,b,c).
\end{aligned}
\end{equation}
If $v_p(\mathrm I),v_p(\mathrm{II})\ge-5(L-1)$ then, using $\mathcal I(L-1)$ for
$v_p(Q_rW_a)\ge v_p(W_a)\ge-5(L-1)$,
\[
v_p(p^5W_n)\ge-5(L-1)\quad\Longrightarrow\quad v_p(W_n)\ge-5-5(L-1)=-5L,
\]
which is $\mathcal I(L)$.

\begin{proposition}[Ledger for the fibre term]\label{prop:ledgerII}
$\PROVED$ Per cell, per digit level:
\[
\begin{array}{lll}
\text{provided: } v_p\vfive(a,b,c) & \ge-5(L-1)-1-\min(s_a,2) & \text{Corollary~\ref{cor:depthgen}}\\
\text{provided: } v_p(\Tcal-Q_rT) & \ge1+\min(s_a,2) & \text{Lemma~\ref{lem:Fgen}}\\
\text{consumed: target} & \ge-5(L-1) & \text{(IND) step}\\
\textbf{slack} & \textbf{exactly }0 & \text{both inputs sharp}
\end{array}
\]
\end{proposition}

So the two lemmas match to the digit: nothing is wasted and nothing is missing in (II). Note
what is \emph{not} needed: $\Lambda=Q_n/Q_a$, the hypothesis $p\nmid Q_a$, the supercongruence
of Corollary~\ref{cor:Qsuper}, and Lemma~\ref{lem:Phi}. The exceptional primes (where
$v_p(Q_n)$ can reach $7$ at $p=7$) are irrelevant to (IND).

\subsection{The descent term, in-regime}

By Lemma~\ref{lem:U}, $p^mH^{(m)}_N=H^{(m)}_{\lfloor N/p\rfloor}+p^m\Sigma^{(m)}_0(N)$. Hence,
letter by letter,
\begin{equation}\label{eq:letterdescent}
\begin{aligned}
p^mA_m^{(n)}(k)&=A_m^{(a)}(b)+e_1(a+b+1)^{-m}+p^m\sigma^A, &\quad \sigma^A&:=\Sigma_0(n+k)-\Sigma_0(k)\in\Zp,\\
p^mB_m^{(n)}(k)&=B_m^{(a)}(b)-\ii{s>r}(a-b)^{-m}+p^m\sigma^B,\\
p^mC_m^{(n)}&=C_m^{(a)}(b,c)+[e_3\text{-terms}]-[e_4\text{-terms}]+p^m\sigma^C,\\
p^mN_m^{(n)}&=N_m^{(a)}+p^m\sigma^N .
\end{aligned}
\end{equation}

\begin{definition}[regimes]\label{def:regime}
A cell is \emph{in-regime} if $s\le r$, $t\le r$ and $e_1=e_2=e_3=e_4=0$; otherwise
\emph{off-regime}.
\end{definition}

In-regime all the mismatch terms vanish and $p^mX^{(n)}_m=X^{(a)}_m+p^m\sigma^X$ exactly.
Expanding the product,
\[
\Ecal=\sum_{S\ne\emptyset}p^{\wt(S)}\sigma_S\Psi_S,\qquad
\sigma_S=\prod_{X\in S}\sigma^X\in\Zp,\qquad
\Psi_S:=\sum_{\Mcal\supseteq S}c_\Mcal\prod_{X\in\Mcal\setminus S}X^{(a)},
\]
a level-$a$ form of weight $5-\wt(S)$.

\begin{proposition}[In-regime descent costs nothing]\label{prop:inregime}
$\PROVED$ For every in-regime cell and every digit level, $v_p(T\Ecal)\ge-5(L-1)$.
\end{proposition}

\begin{proof}
By the argument of Proposition~\ref{prop:LIFT} applied to $\Psi_S$ --- its $u$-order is
$(5-\wt S)M_a-\delta'$ and its top non-vanishing $\delta'$ obeys $\wt(S)+\delta'\ge5-J(\pi)$
because $\Phi_\pi(S\cup S')=0$ for $\wt(S\cup S')\le4-J(\pi)$ --- one gets
\[
v_p\bigl(p^{\wt S}\sigma_S\Psi_S\bigr)\ \ge\ 5-J(\pi)-(5-\wt S)M_a
\]
whenever $\wt(S)\le5-J(\pi)$; if $\wt(S)>5-J(\pi)$ then $\wt(S)\ge6-J(\pi)\ge3$, so
$\wt(S)M_a\ge3\ge J(\pi)$ and the requirement below holds trivially. In the fully in-regime
case Lemma~\ref{lem:carry}(C2) is an equality at every slot, so $v_pT(n,k,l)=v_pT(a,b,c)=vT$. The requirement
$v_p(T\Ecal)\ge-5(M_a-1)$ becomes $vT\ge J(\pi)-\wt(S)M_a$, and with $J(\pi)=1+\min(s_a,2)\le3$,
$\wt(S)\ge1$, $vT\ge s_a$:
\[
\begin{array}{cccl}
M_a & \wt(S) & \text{needed }vT\ge & \\\hline
1 & 1 & \min(s_a,2) & \text{have }vT\ge s_a\ \checkmark\\
1 & 2 & \min(s_a,2)-1 & \checkmark\\
1 & \ge3 & \le0 & \checkmark\\
\ge2 & \ge1 & \le1,\ =1\text{ only when }J=3\text{ i.e. }s_a\ge2 & \text{have }vT\ge2\ \checkmark
\end{array}
\]
\end{proof}

%%%%%%%%%%%%%%%%%%%%%%%%%%%%%%%%%%%%%%%%%%%%%%%%%%%%%%%%%%%%%%%%%%%%%%%%%%%%%%%%
\section{The descent term off-regime}\label{sec:gapdesc}
%%%%%%%%%%%%%%%%%%%%%%%%%%%%%%%%%%%%%%%%%%%%%%%%%%%%%%%%%%%%%%%%%%%%%%%%%%%%%%%%

Fix $p\ge5$ and $n\ge p$; put $L=\lfloor\log_pn\rfloor\ge1$, $M=L+1$, $P=p^M$,
$a=\lfloor n/p\rfloor$, $r=n-ap$, $P_a=p^L$ (so $p^{L-1}\le a<P_a$ and
$\lfloor\log_pa\rfloor=L-1$). For a cell write $b=\lfloor k/p\rfloor$, $s=k-bp$,
$c=\lfloor l/p\rfloor$, $t=l-cp$, and set $\zeta:=e_3-e_4$. Let $(\alpha,\gamma,\kappa,\theta)$
be the level-$n$ pattern (with $P$) and $(\alpha_a,\gamma_a,\kappa_a,\theta_a)$ the level-$a$
pattern (with $P_a$); $s_n:=\alpha+\gamma+\kappa$, $s_a:=\alpha_a+\gamma_a+\kappa_a$.

\subsection{The digit dictionary}

\begin{lemma}[Digit dictionary]\label{lem:D1}
$\PROVED$
\begin{enumerate}[leftmargin=1.9em,label=\textup{(\arabic*)}]
\item $\lfloor(n+k)/p\rfloor=a+b+e_1$, $\lfloor(n+l)/p\rfloor=a+c+e_2$,
  $\lfloor(n+k+l)/p\rfloor=a+b+c+e_3$, $\lfloor(k+l)/p\rfloor=b+c+e_4$,
  $\lfloor(n-k)/p\rfloor=a-b-\ii{s>r}$, $\lfloor n/p\rfloor=a$.
\item $0\le\zeta=e_3-e_4\le1$, and $\zeta=1$ iff the addition $n+(k+l)$ carries out of
  position $0$.
\item $\alpha=\ii{a+b+e_1\ge P_a}\ge\alpha_a$ and $\gamma=\ii{a+c+e_2\ge P_a}\ge\gamma_a$.
\item $\varepsilon:=\lfloor(k+l)/P\rfloor=\lfloor(b+c+e_4)/P_a\rfloor
  =\varepsilon_a+e_4\ii{b+c=P_a-1}$, and $\kappa=\ii{a+b+c+e_3\ge(\varepsilon+1)P_a}$.
\item $\kappa\ge\kappa_a$ \emph{unless} $e_4=1$ and $b+c=P_a-1$; hence in all cases
  $s_n\ge s_a-\ii{e_4=1}$.
\item In-regime, $s_n=s_a$ and $v_pT(n,k,l)=v_pT(a,b,c)$.
\end{enumerate}
\end{lemma}

\begin{proof}
(1) is the definition of the four indicators ($n=ap+r$, $k=bp+s$, $l=cp+t$, and
$n-k=(a-b-\ii{s>r})p+((r-s)\bmod p)$).
(2) $e_4\le e_3$ because $r+s+t\ge s+t$, and $e_3\le e_4+1$ because $r<p$. The digit of $k+l$
in position $0$ is $s+t-e_4p$, so $n+(k+l)$ carries out of position $0$ iff
$r+s+t-e_4p\ge p$ iff $e_3\ge e_4+1$.
(3) $n+k\ge P=p\,P_a$ iff $\lfloor(n+k)/p\rfloor\ge P_a$ iff $a+b+e_1\ge P_a$, and
$a+b+e_1\ge a+b$.
(4) $b+c\le2a\le2P_a-2$, so the interval $[1,b+c+1]$ shifted appropriately contains exactly one
multiple of $P_a$, namely $P_a$ itself; hence
$\lfloor(b+c+1)/P_a\rfloor-\lfloor(b+c)/P_a\rfloor=\ii{b+c=P_a-1}$. The formula for $\kappa$ is
$n+k+l\ge(\varepsilon+1)P\iff\lfloor(n+k+l)/p\rfloor\ge(\varepsilon+1)P_a$.
(5) If $\varepsilon=\varepsilon_a$ then
$\kappa=\ii{a+b+c+e_3\ge(\varepsilon_a+1)P_a}\ge\ii{a+b+c\ge(\varepsilon_a+1)P_a}=\kappa_a$.
Otherwise $\varepsilon$ exceeds $\varepsilon_a$, which by (4) forces $e_4=1$ and $b+c=P_a-1$;
and always $\kappa\ge\kappa_a-1$ since both are $0/1$. Adding (3) gives the display.
(6) In-regime all four $e_i$ vanish, so (3),(4) give $\alpha=\alpha_a$, $\gamma=\gamma_a$,
$\varepsilon=\varepsilon_a$, $\kappa=\kappa_a$.
\end{proof}

Item (5) is sharp: $e_4=1$ with $b+c=P_a-1$ really does drop $\kappa$ from $1$ to $0$. Example:
$p=5$, $n=101$, $(k,l)=(63,63)$, so $L=2$, $a=20$, $r=1$, $b=c=12$, $s=t=3$, $b+c=24=P_a-1$,
$(e_3,e_4)=(1,1)$, level-$a$ pattern $(1,1,1,1)$ with $s_a=3$, level-$n$ pattern $(1,1,0,1)$
with $s_n=2$; there $B=4$ and $v_pT=10\ge1+\max(s_n,s_a)=4$.

\subsection{The descent Kummer lemma}

Define the \emph{bottom-carry count}
\begin{equation}\label{eq:Bcount}
B:=e_1+e_2+2\ii{s>r}+2\ii{t>r}+\zeta .
\end{equation}

\begin{lemma}[Descent Kummer]\label{lem:DK}
$\PROVED$ Let $p\ge5$, $n\ge p$, $L=\lfloor\log_pn\rfloor\ge1$, and let $(k,l)$ be off-regime.
Then
\[
\begin{gathered}
\text{(DK1)}\ \ B\ge1,\ \text{and if } e_4=1 \text{ then } B\ge2;\\
\text{(DK2)}\ \ v_pT(n,k,l)\ge s_n+B;\qquad
\text{(DK3)}\ \ v_pT(n,k,l)\ge1+\max(s_n,s_a).
\end{gathered}
\]
\end{lemma}

\begin{proof}
\emph{(DK1).} Three exhaustive cases. If $s>r$ or $t>r$ then $B\ge2$ outright, from the term
$2\ii{s>r}$ resp.\ $2\ii{t>r}$. If $s\le r$, $t\le r$ and $e_4=1$ then $r+s\ge s+t\ge p$ and
$r+t\ge s+t\ge p$, so $e_1=e_2=1$ and again $B\ge2$. If $s\le r$, $t\le r$ and $e_4=0$ then
$\zeta=e_3$, and off-regime forces $e_1+e_2+e_3\ge1$; each summand contributes $1$ to $B$ (for
$e_3$ through $\zeta=e_3$), so $B\ge1$. Since $e_4=1$ occurs only in the first two cases,
$e_4=1\Rightarrow B\ge2$.

\emph{(DK2).} Write $T=\binom{n+k}{n}\binom nk^2\binom{n+l}{n}\binom nl^2\binom{n+k+l}{n}$ and
let $V_1,\dots,V_5$ be the $v_p$ of the five binomials, so $v_pT=V_1+V_2+2V_3+2V_4+V_5$. By
Kummer each $V_i$ is a count of carries in one addition. In each addition we locate the
\emph{bottom} carry (out of base-$p$ position $0$) and, where the level-$n$ pattern demands
one, the \emph{top} carry (out of position $M-1=L$); the two positions are distinct precisely
because $L\ge1$ --- this is the one and only place the hypothesis is used.
\begin{itemize}[leftmargin=1.6em]
\item $V_1$ counts the carries of $n+k$. Position $0$: carries iff $r+s\ge p$, i.e.\ iff
  $e_1=1$. Position $L$: since $n,k<P$, it carries iff $n+k\ge P$, i.e.\ iff $\alpha=1$
  (Lemma~\ref{lem:K}). Hence $V_1\ge\alpha+e_1$; likewise $V_2\ge\gamma+e_2$.
\item $V_3$ counts the carries of $k+(n-k)$. In position $0$ the digits are $s$ and
  $(r-s)\bmod p$, summing to $r+p\ii{s>r}$, so it carries iff $s>r$. Hence $V_3\ge\ii{s>r}$,
  likewise $V_4\ge\ii{t>r}$.
\item $V_5$ counts the carries of $n+(k+l)$. Position $0$: carries iff $\zeta=1$
  (Lemma~\ref{lem:D1}(2)). Position $L$: writing $k+l=\varepsilon P+\rho$ with $0\le\rho<P$,
  the low $M$ digits of the addition sum to $n+\rho$, so it carries out of position $M-1=L$ iff
  $n+\rho\ge P$, i.e.\ iff $\kappa=1$. Hence $V_5\ge\kappa+\zeta$.
\end{itemize}
Summing with multiplicities $(1,1,2,2,1)$,
$v_pT\ge(\alpha+\gamma+\kappa)+(e_1+e_2+2\ii{s>r}+2\ii{t>r}+\zeta)=s_n+B$.

\emph{(DK3).} By (DK2) and (DK1), $v_pT\ge s_n+1$. For the level-$a$ half,
Lemma~\ref{lem:D1}(5) gives $s_n\ge s_a-\ii{e_4=1}$ and (DK1) gives $B\ge1+\ii{e_4=1}$; hence
$v_pT\ge s_n+B\ge(s_a-\ii{e_4=1})+(1+\ii{e_4=1})=s_a+1$.
\end{proof}

The mechanism is one line: \emph{off-regime means a carry in base-$p$ position $0$, the
depth-pattern indicators live in position $L\ge1$, and Kummer counts both.}

\begin{theorem}[Off-regime descent]\label{thm:gapdesc}
$\PROVED$ Let $p\ge5$, $n\ge p$, $L=\lfloor\log_pn\rfloor\ge1$, $a=\lfloor n/p\rfloor$, and let
$(k,l)$ with $0\le k,l\le n$ be off-regime. Then, with
$\Ecal=p^5\vfive(n,k,l)-\vfive(a,b,c)$,
\[
v_p\bigl(T(n,k,l)\,\Ecal(n,k,l)\bigr)\ \ge\ -5(L-1).
\]
\end{theorem}

\begin{proof}
Corollary~\ref{cor:depthgen} at level $n$ ($\lfloor\log_pn\rfloor=L$, pattern sum $s_n$) gives
$v_p\vfive(n,k,l)\ge-5L-1-\min(s_n,2)$, hence
$v_p(p^5\vfive(n,k,l))\ge-5(L-1)-1-\min(s_n,2)$. The same corollary at level $a$
($\lfloor\log_pa\rfloor=L-1$, pattern sum $s_a$) gives
$v_p\vfive(a,b,c)\ge-5(L-1)-1-\min(s_a,2)$. Therefore
\[
v_p(\Ecal)\ \ge\ -5(L-1)-1-\max\bigl(\min(s_n,2),\min(s_a,2)\bigr)\ \ge\ -5(L-1)-1-\max(s_n,s_a),
\]
while Lemma~\ref{lem:DK} gives $v_pT(n,k,l)\ge1+\max(s_n,s_a)$. Add.
\end{proof}

\begin{corollary}[The descent term]\label{cor:termI}
$\PROVED$ For every $p\ge5$ and every $n$ with $L=\lfloor\log_pn\rfloor\ge1$,
$v_p((\mathrm I))\ge-5(L-1)$, cell by cell: off-regime by Theorem~\ref{thm:gapdesc}, in-regime
by Proposition~\ref{prop:inregime}.
\end{corollary}

\begin{corollary}[The induction step]\label{cor:step}
$\PROVED$ $\mathcal I(L-1)\Rightarrow\mathcal I(L)$ for every $L\ge1$.
\end{corollary}

\begin{proof}
In \eqref{eq:split}, term (II) is bounded by Proposition~\ref{prop:ledgerII} and term (I) by
Corollary~\ref{cor:termI}.
\end{proof}

\begin{proof}[Proof of Theorem~\textup{\ref{thm:main}}]
$\mathcal I(0)$ is Theorem~\ref{thm:baseintro}, proved in \S\ref{sec:nucleus}; the step is
Corollary~\ref{cor:step}. Hence $v_p(W_m)\ge-5\lfloor\log_pm\rfloor$ for all $m\ge1$, and
therefore $\ord_p(P_m)\ge-5\lfloor\log_pm\rfloor$.
\end{proof}

\begin{remark}[scope]
Theorem~\ref{thm:gapdesc} is uniform in the level: nothing in \S\ref{sec:gapdesc} distinguishes
$a<p$ from $a\ge p$, and the hypothesis $L\ge1$ enters exactly once, in (DK2), to know that
digit positions $0$ and $L$ are distinct. At $L=0$ there is no descent to make. It uses no
$p\nmid Q_a$, no $\Lambda$, no supercongruence, no Lemma~\ref{lem:Phi}, no Lemma~\ref{lem:Dpp},
and --- in itself --- no Wolstenholme: $p\ge5$ is inherited only from
Corollary~\ref{cor:depthgen}. And it consumes $\wfive$ only through
Corollary~\ref{cor:depthgen}, so it holds verbatim for any point of the depth-conditioned
family, in particular for $w_5^{\mathrm I}$.
\end{remark}

\subsection{Why the expected route fails}\label{sec:trap}

It is worth recording that the natural letter-wise route to Theorem~\ref{thm:gapdesc} does not
exist, since the failure is instructive. That route reads off from \eqref{eq:letterdescent} the
mismatch $e_1(a+b+1)^{-m}$ and observes that a mismatch \emph{forces carries}:
\begin{equation}\label{eq:trap}
v_p\tbinom{n+k}{n}=v_p\tbinom{a+b}{a}+e_1(1+\lambda_1),\qquad \lambda_1:=v_p(a+b+1)
\end{equation}
(correct, and a consequence of Lemma~\ref{lem:carry}(C1) with $z(a+b)=\lambda_1$). \emph{But at weight $5$ the trade
loses.} Expanding $\Ecal$ letter-wise as $\sum_{S\ne\emptyset}(\prod_{X\in S}D_X)\Psi_S$ with
$D_X$ the per-letter mismatch, the single term $S=\{A_5(k)\}$ contributes
$c\,(a+b+1)^{-5}$ of valuation $-5\lambda_1$, where $c$ is the $\wfive$-coefficient of the
symmetrised $A_5(k)+A_5(l)$ monomial ($=3830046703/48$ for $w_5^{\mathrm{allp}}$, a $p$-unit),
against a total Kummer gain of only $1+\lambda_1$ in the whole of $T$. Cell-wise:
\[
\begin{array}{cccccccc}
p & n & (k,l) & (a,b) & \lambda_1 & v_pT & \text{letter-wise } vT-5\lambda_1 & \text{target }-5(L-1)\\\hline
5 & 19 & (6,0) & (3,1) & 1 & 4 & \mathbf{-1} & 0\\
5 & 19 & (6,19) & (3,1) & 1 & 4 & \mathbf{-1} & 0\\
5 & 99 & (26,30) & (19,5) & 2 & 8 & -2 & -5\\
7 & 244 & (100,50) & (34,14) & 2 & 8 & -2 & -5
\end{array}
\]
At $p=5$, $n=19$, $(k,l)=(6,0)$ the letter-wise ledger is short by exactly one power, and no
refinement of the carry bookkeeping can repair it because \eqref{eq:trap} is an
\emph{equality}. (The true value there is $v_p(T\Ecal)=1$, comfortably above the target.) What
actually happens is that the pole is not present in $\Ecal$ at all: the mismatch poles of the
individual letters cancel among themselves \emph{inside} $\vfive$, because the (DEPTH)
conditions annihilate its top $u$-coefficients at level $n$. Hence the correct instrument is
Corollary~\ref{cor:depthgen}, not a letter-wise ledger --- and the only thing left to check is
that $T$ pays for the \emph{two} depth caps, which is Lemma~\ref{lem:DK}. This is the weight-$5$
shadow of the fact that the $a<p$ analysis could afford the trade only up to weight $3$:
Lemma~\ref{lem:Dpp} supplies $v_pT\ge4$, which covers $A_3$ but not $A_5$.

\begin{evidence}\label{ev:DK}
All exact; $p$-adic valuations computed with tracked precision, the evaluator cross-checked
against exact \texttt{Fraction} evaluation of the $178$-term $w_5^{\mathrm{allp}}$ on $1290$
cells with $0$ mismatches of valuation or unit.
\[
\small
\begin{array}{llrr}
\text{sweep} & \text{range} & \text{cells} & \text{failures}\\\hline
\text{Lem.~\ref{lem:DK} (DK1)--(DK3), Lem.~\ref{lem:D1}(2),(5)}
 & 931\text{ levels},\ p\le31,\ L=1,2,3 & 142\,820\,576 & \mathbf 0\\
\text{the same, at } L=4\ (\text{and }L=3) & 13\text{ levels},\ p\le11,\ n\le3124 & 45\,533\,157 & \mathbf 0\\
\text{slot-wise inequalities behind (DK2)} & p\le31,\ L=1,2,3 & 11\,096\,075 & \mathbf 0\\
\text{Thm~\ref{thm:gapdesc}, exact }p\text{-adic }\Ecal & p\le13,\ L=1 & 174\,527 & \mathbf 0\\
\text{Thm~\ref{thm:gapdesc}, exact }p\text{-adic }\Ecal & p\le13,\ L=2 & 530\,900 & \mathbf 0\\
\text{Thm~\ref{thm:gapdesc}, exact }p\text{-adic }\Ecal & p\le7,\ L=3 & 934\,656 & \mathbf 0\\
\text{whole term (I), both regimes, }+\ v_p(\sum T\Ecal) & p\le13,\ L=2,3 & 15\,553+345\,334 & \mathbf 0\\
\text{Cor.~\ref{cor:depthgen} and Thm~\ref{thm:gapdesc} for } w_5^{\mathrm I} & p\le13,\ L\le2 & 15\,228+43\,213 & \mathbf 0
\end{array}
\]
The first two rows are decisive: together they check the criterion on \emph{every} cell of $944$
different levels --- $188\,353\,733$ off-regime cells, $0$ failures, no sampling inside a level.
\emph{Control:} in-regime, $B=0$ and $s_n=s_a$ at every one of the $13\,082\,671$ in-regime
cells of the sweep (as Lemma~\ref{lem:D1}(6) predicts) and the criterion
$v_pT\ge\max(J(\pi_n),J(\pi_a))$ \emph{fails} at $134\,964$ of them --- by design. So
off-regime is not a convenience of the proof: it is exactly the hypothesis that buys the extra
power of $p$. \emph{Sharpness:} the slack is $0$ at $41\,284$ cells, so neither
Lemma~\ref{lem:DK} nor Theorem~\ref{thm:gapdesc} can be weakened; and the true quantity
$v_p(T\Ecal)+5(L-1)$ attains $0$ (e.g.\ $p=13$, $n=153$, $(k,l)=(1,41)$: $v_pT=3$,
$v_p\Ecal=-3$, target $0$).
\end{evidence}

\begin{remark}[a formalisation target]
Lemmas~\ref{lem:D1} and \ref{lem:DK} are pure base-$p$ combinatorics over $\NN$ --- no harmonic
numbers, no representative, no analysis. They are the most formalisation-ready statements in
the chain after Theorem~\ref{thm:Aintro} (already $\LEAN$), and formalising them would make the
Kummer half of the induction step machine-checked, leaving Corollary~\ref{cor:depthgen} as the
analytic input.
\end{remark}

%%%%%%%%%%%%%%%%%%%%%%%%%%%%%%%%%%%%%%%%%%%%%%%%%%%%%%%%%%%%%%%%%%%%%%%%%%%%%%%%
\section{The nucleus: the base case}\label{sec:nucleus}
%%%%%%%%%%%%%%%%%%%%%%%%%%%%%%%%%%%%%%%%%%%%%%%%%%%%%%%%%%%%%%%%%%%%%%%%%%%%%%%%

By \S\ref{sec:hardneg} the base case is not cell-wise for any $p$-independent representative,
and by \S\ref{sec:wall} below no summand-side identity can supply it. What does work is a split
of the range $n<p$ at the midpoint. Throughout this section $\wfive=w_5^{\mathrm I}$ and
$\vfive=w_5^{\mathrm I}-H^{(5)}_n$.

\subsection{One band}

\begin{theorem}[Partial closure]\label{thm:oneband}
$\PROVED$ given \textup{(T1-top)} for $w_5^{\mathrm I}$ and the $\CERT$ linear-algebra
certificate. For every prime $p\ge5$ and every $n<p$, every cell outside the band
\begin{equation}\label{eq:III}
\mathrm{III}=\bigl\{(k,l):\ k,l\ge q:=p-n,\ \ p\le k+l<p+q\bigr\}
\end{equation}
contributes a $p$-integral summand to $P_n=\sum_{k,l}T(n,k,l)w_5^{\mathrm I}(n,k,l)$.
Consequently
\[
\text{\textup{(BASE)}}\quad\Longleftrightarrow\quad
v_p\Bigl(\sum_{(k,l)\in\mathrm{III}}T(n,k,l)\,\vfive(n,k,l)\Bigr)\ \ge\ 0 .
\]
\end{theorem}

The conditions cutting out $w_5^{\mathrm I}$ are the full cell-wise cap $d_5\le v_pT$ on every
pole pattern \emph{except} the one giving $\mathrm{III}$, where the cap of \eqref{eq:cap} is
kept. They are $p$-independent, so their sufficiency is the pole calculus of \S\ref{sec:depth};
the joint system has $\rank(\mathrm{cond})=111$, $\rank(\mathrm{joint})=342$ and is consistent,
against $\rank(\mathrm{cond})=123$, $\rank(\mathrm{joint})=342$ and \emph{inconsistent} with no
exemption. So \emph{the entire residue of the base case is one order of pole on one pattern
group}, and the harmonic alphabet already supplies everything else.

This is a strict sharpening: the previously known reduction was $2\Sigma_{\mathrm I}+\Sigma_{\mathrm{III}}\equiv0$
over two regions, and the attained deficit cell $\bigl(\tfrac{p+1}{2},0,\tfrac{p-1}{2}\bigr)$
has $k=0<q$, so it lies in region $\mathrm I$: \emph{the representative $w_5^{\mathrm I}$
genuinely moves the deficit off the cell where it was found.}

\begin{evidence}\label{ev:oneband}
$\VERIFIED$, exact \texttt{Fraction} arithmetic: (V1) $P_n=\sum T\,w_5^{\mathrm I}$ exactly over
$\QQ$ for $n=1,\dots,20$ (and $n=1,\dots,16$ for the all-primes form), $0$ mismatches;
(V2) $v_p(T\vfive)\ge0$ at every cell outside $\mathrm{III}$ over
$p=5,7,11,13,17,19,23$: $10\,092$ cells ($1\,025$ exempt), $0$ violations; and $5\,769$ cells,
$0$ violations, for the all-primes form; (V3) the reduction itself, $0$ failures.
\end{evidence}

\subsection{Below the midpoint the band is empty}

\begin{lemma}\label{lem:below}
$\PROVED$ If $n\le(p-1)/2$ then $\mathrm{III}=\emptyset$, hence $v_p(P_n)\ge0$.
\end{lemma}

\begin{proof}
$\mathrm{III}$ requires $k+l\ge p$ with $k,l\le n$, so $2n\ge p$; but $2n\le p-1$.
\end{proof}

So half of the base case is free, and the deficit can first appear only at $n=(p+1)/2$ ---
exactly where \S\ref{sec:hardneg} located the attained deficit cell. The two localisations
agree for a structural reason: the midpoint singularity of $\LBZ$ ($2n+5\equiv0$) and the
midpoint appearance of $\mathrm{III}$ ($2n\ge p$) are the same phenomenon seen from the
recurrence side and from the summand side.

\subsection{At the midpoint the band is three cells}

Let $n=(p+1)/2$, $q=p-n=(p-1)/2$.

\begin{lemma}\label{lem:threecells}
$\PROVED$ $\mathrm{III}=\{(q,q+1),(q+1,q),(q+1,q+1)\}=\{(n-1,n),(n,n-1),(n,n)\}$.
\end{lemma}

\begin{proof}
$k,l\ge q$ and $k,l\le n=q+1$ force $k,l\in\{q,q+1\}$. Then $k+l\in\{p-1,p,p+1\}$; the
condition $k+l\ge p$ excludes $(q,q)$, and $k+l<p+q$ holds for the other three since
$p+1<p+(p-1)/2$ for $p\ge5$.
\end{proof}

\begin{lemma}\label{lem:Tcorner}
$\PROVED$ At $n=(p+1)/2$, $p\ge5$: $v_p(T)=2$ at all three cells and
\[
T(n,n-1,n)/p^2\equiv T(n,n,n-1)/p^2\equiv2,\qquad T(n,n,n)/p^2\equiv24 \pmod p .
\]
\end{lemma}

\begin{proof}
Write $\epsilon:=(-1)^{(p+1)/2}$, $q=(p-1)/2$; Wilson's theorem gives $(q!)^2\equiv\epsilon$.

\emph{Cell $B=(n,n)$.} $T=\binom{p+1}{n}^2\binom{3n}{n}$. Kummer: $n+n=p+1$ has exactly one
carry in base $p$, so $v_p\binom{p+1}{n}=1$; $n+(p+1)$ has none, so $v_p\binom{3n}{n}=0$; thus
$v_p(T)=2$. As $p+1-n=n$, $\binom{p+1}{n}/p=(p+1)(p-1)!/(n!)^2\equiv-1/(n!)^2$; with
$n!=(q+1)q!$ and $q+1\equiv1/2$, $(n!)^2\equiv\epsilon/4$, so $\binom{p+1}{n}/p\equiv-4\epsilon$
and $(\binom{p+1}{n}/p)^2\equiv16$. Lucas on $3n=p+(p+3)/2$, $n=(p+1)/2$ gives
$\binom{3n}{n}\equiv\binom{(p+3)/2}{(p+1)/2}=(p+3)/2\equiv3/2$. Hence
$T/p^2\equiv16\cdot(3/2)=24$.

\emph{Cell $A=(n-1,n)=(q,q+1)$.} $\binom{n+k}{n}=\binom pn$ with $v_p=1$ and
$\binom pn/p=(p-1)!/(n!q!)\equiv-1/((q+1)(q!)^2)\equiv-2\epsilon$; $\binom nk=\binom{q+1}{q}=q+1\equiv1/2$;
$\binom{n+l}{n}=\binom{p+1}{n}$ with $\binom{p+1}{n}/p\equiv-4\epsilon$; $\binom nl=1$;
$\binom{n+k+l}{n}=\binom{(3p+1)/2}{(p+1)/2}\equiv\binom10\binom{(p+1)/2}{(p+1)/2}=1$ by Lucas.
So $v_p(T)=2$ and $T/p^2\equiv(-2\epsilon)(1/2)^2(-4\epsilon)=2\epsilon^2=2$.
\end{proof}

At $n=(p+1)/2$ every letter argument is $<2p$, so each letter contains at most one multiple of
$p$ and its $u$-expansion is exact and elementary. With $h_r:=H^{(r)}_q\bmod p$,
$s_r:=\sum_{j=1}^{p-1}j^{-r}\bmod p$, $(q+1)^{-1}\equiv2$, $(q+2)^{-1}\equiv2/3$:
\[
\begin{array}{lll}
\text{letter} & \text{cell } A=(q,q+1) & \text{cell } B=(q+1,q+1)\\\hline
A_r(k) & u^r+s_r-h_r & u^r+s_r+1-h_r-2^r\\
B_r(k) & 1-h_r & -h_r-2^r\\
A_r(l) & u^r+s_r+1-h_r-2^r & u^r+s_r+1-h_r-2^r\\
B_r(l) & -h_r-2^r & -h_r-2^r\\
C_r\ (\text{no pole}) & h_r+2^r & h_r+2^r+(2/3)^r-1\\
N_r & h_r+2^r & h_r+2^r
\end{array}
\]

\begin{lemma}\label{lem:K3}
$\PROVED$ (finite exact symbolic computation over $\QQ[h_1,\dots,h_5,s_1,\dots,s_5]$). For every
prime $p\ge5$, substituting the table above into the $207$ monomials of $w_5^{\mathrm I}$ and
extracting the $u^j$ coefficients gives
\[
K_4=K_5=0\ \text{ at both cells},\qquad K_3(A)=3-\tfrac{s_2}{2},\qquad K_3(B)=-\tfrac12-\tfrac{s_2}{2}.
\]
Both are \emph{free of $h_1,\dots,h_5$ and of $s_1,s_3,s_4,s_5$}: the entire harmonic content of
the $u^3$ coefficient cancels.
\end{lemma}

Only letters of weight $\le3$ can reach $K_3$ --- the pole part must have total weight $3$ and
the remainder weight $2$ --- which is why $s_4,s_5$ cannot appear and the lemma is valid down to
$p=5$. That $K_4=K_5=0$ is the Lemma-$F$ cap $d_5\le3$ being met.

\begin{theorem}[The nucleus, dissolved]\label{thm:nucleus}
$\PROVED$ For every prime $p\ge5$, at $n=(p+1)/2$,
\[
\sum_{(k,l)\in\mathrm{III}}\frac{T}{p^2}\,K_3
\ \equiv\ 2\Bigl(3-\frac{s_2}{2}\Bigr)+2\Bigl(3-\frac{s_2}{2}\Bigr)+24\Bigl(-\frac12-\frac{s_2}{2}\Bigr)
\ =\ -14\,s_2\ \equiv\ 0\pmod p,
\]
because $s_2=\sum_{j=1}^{p-1}j^{-2}\equiv\sum_{j=1}^{p-1}j^{2}=\frac{(p-1)p(2p-1)}{6}\equiv0$
for $p\ge5$. Since $K_0,K_1,K_2\in\Zp$ and $T/p^2\in\ZZ$, this gives
$v_p(\sum_{\mathrm{III}}T\vfive)\ge0$, hence by Theorem~\ref{thm:oneband}
$v_p(P_{(p+1)/2})\ge0$.
\end{theorem}

\emph{The whole nucleus consumes exactly one classical fact}, $\sum_{j<p}j^{-2}\equiv0$, i.e.\
the same $p\ge5$ hypothesis Wolstenholme imposes elsewhere. No Fermat quotient, no Bernoulli
number.

\begin{remark}[the statement is not vacuous]
The same extraction on $w_5^{\mathrm{allp}}$ gives an $h_1$-dependent $K_3$ with non-vanishing
corner combination. The content of Theorem~\ref{thm:nucleus} lives precisely in the extra depth
conditions defining $w_5^{\mathrm I}$ --- exactly as the wall theorem of \S\ref{sec:wall}
predicts: the value cannot be moved by summand-side identities, only by pinning the
representative. \emph{The wall is respected, and the door it points at is the representative.}
\end{remark}

\begin{evidence}\label{ev:corner}
$\VERIFIED$ $0$ failures: $v_p(T)=2$ and $T/p^2\bmod p=2,2,24$ at every prime
$p\in\{5,\dots,31\}$; and, by exact \texttt{Fraction} arithmetic through an independent
evaluator, at every prime $5\le p\le139$ the three corner residues $p\,T\vfive\bmod p$ are
exactly $(6,6,-12)\bmod p$ and sum to $0$ --- $32/32$ primes clean. Cell-wise integrality below
the midpoint: $1\,588$ cells over all levels $n<(p+1)/2$, $p\le23$, $0$ violations, plus the top
below-midpoint level $n=(p-1)/2$ at $p=29,31,37$ ($842$ cells, $|\mathrm{III}|=0$ as
Lemma~\ref{lem:below} predicts).
\end{evidence}

\subsection{Above the midpoint: the exceptional steps are apparent}\label{sec:desing}

For $n>(p+1)/2$ we use forward induction along $\LBZ$. To produce level $m\le p-1$ we use the
step $\nu=m-3\ge(p-3)/2$; on that range $2\nu+5\in[p+2,2p-3]$ is odd and $\ne p$, and
$\nu+3\le p-1$. \emph{Hence the only possible exceptional steps above the midpoint are roots of
$a_0$ modulo $p$.}

\begin{lemma}[Apparency]\label{lem:APP}
$\PROVED$ by exact polynomial division. Put $U(\nu):=(c_0,c_1,c_2)(\nu)$ and
$V(\nu):=(c_1,c_2,c_3)(\nu-1)$, both acting on $(Y_\nu,Y_{\nu+1},Y_{\nu+2})$. All three
$2\times2$ minors of $[U(\nu);V(\nu)]$ are divisible by $a_0(\nu)$ in $\QQ[\nu]$ (each minor has
degree $18$; the remainders on division by $a_0$ are exactly $0$). Hence if $p\mid a_0(\nu)$,
$\nu\ge1$, $V(\nu)\not\equiv0\pmod p$, and $Y$ solves $\LBZ$ with
$Y_{\nu-1},\dots,Y_{\nu+2}\in\Zp$, then $U(\nu)\cdot(Y_\nu,Y_{\nu+1},Y_{\nu+2})\equiv0
\pmod{p^{v_p(a_0(\nu))}}$, so $Y_{\nu+3}\in\Zp$.
\end{lemma}

\begin{proof}
Because $c_0(\nu-1)=\nu^5(\nu+1)a_0(\nu)$, the row at $\nu-1$ degenerates in the same place as
the row at $\nu$; it gives $V(\nu)\cdot(Y_\nu,Y_{\nu+1},Y_{\nu+2})=-\nu^5(\nu+1)a_0(\nu)Y_{\nu-1}
\equiv0\pmod{p^v}$ with $v:=v_p(a_0(\nu))$. Each minor is $a_0(\nu)\cdot(\text{integer})$, so
$\equiv0\pmod{p^v}$. Pick $i$ with $V_i$ a $p$-unit and set $\lambda:=U_i/V_i\in\Zp$; then
$U_j-\lambda V_j=(U_jV_i-U_iV_j)/V_i\equiv0\pmod{p^v}$, so
$U\cdot Y\equiv\lambda(V\cdot Y)\equiv0\pmod{p^v}$. On this range $v_p(c_3(\nu))=v$.
\end{proof}

Lemma~\ref{lem:APP} is the shadow of an explicit desingularisation. $a_0$ is irreducible over
$\QQ$, so $K:=\QQ[\nu]/(a_0)$ is a field and $U=\lambda V$ in $K$ with
$\lambda=c_0(\nu)/c_1(\nu-1)$ of degree $\le2$:
\[
\lambda(\nu)=\frac{392627556035671426586\,\nu^2+1282015597875460006266\,\nu
+1052781309790247665282}{D},
\]
\[
D=3641620092914355321=3\cdot7\cdot11\cdot29\cdot p_0,\qquad p_0:=543606522303979 .
\]
Every coefficient of $\mathrm{row}(\nu)-\lambda\,\mathrm{row}(\nu-1)$ --- an order-$4$ relation
on $Y_{\nu-1},\dots,Y_{\nu+3}$ --- is divisible by $a_0(\nu)$ in $\QQ[\nu]$; dividing through
and clearing $D$ leaves:

\begin{proposition}[Desingularisation]\label{prop:Ltilde}
$\PROVED$, $\VERIFIED$ There is an order-$4$ left multiple
$\tilde L:\ d_0Y_{\nu-1}+d_1Y_\nu+d_2Y_{\nu+1}+d_3Y_{\nu+2}+d_4Y_{\nu+3}=0$ of $\LBZ$ with
$d_i\in\ZZ[\nu]$, $\deg d_0,\dots,d_3=8$ and
\[
d_4=2D(\nu+3)^5(2\nu+5)=7283240185828710642\,(\nu+3)^5(2\nu+5):
\]
$a_0$ is gone from the leading coefficient. ($\tilde L$ annihilates $Q$, $P$ and $\hat P$
exactly over $\QQ$ for $\nu=1,\dots,40$, $0$ residuals. Note $d_1$ is not identically zero: the
$Y_\nu$ coefficient vanishes only modulo $a_0$, and
$d_1=(c_0(\nu)-\lambda c_1(\nu-1))/a_0(\nu)$ is a genuine degree-$8$ polynomial.)
\end{proposition}

Consequently, \emph{on the desingularised ladder the only singular step in $0\le n\le p-4$ is
the midpoint $n_0=(p-5)/2$} (the alternative $\nu+3\equiv0$ gives $\nu=p-3$, out of range), for
every $p\notin\{2,3,7,11,29,p_0\}$. The recurrence half and the
region half of the proof dovetail with no residue: $\tilde L$'s single singular step is
precisely the single level that Theorem~\ref{thm:nucleus} supplies.

%% FIXED [M-5]: (1) ``exactly three occurrences'' was false without a range hypothesis --- the
%% FIXED fourth pair is $(p,\nu)=(29,27)$, out of range at $p=29$ since $[(p-3)/2,p-4]=[13,25]$;
%% FIXED the lemma now carries the hypothesis $1\le\nu\le p-4$ and lists (29,27) explicitly.
%% FIXED (2) the uniqueness claim for $p_0$ was false: $\Res_\nu(a_0(\nu),a_0(\nu-1))
%% FIXED = -2^2\cdot11\cdot37^2\cdot557^3\cdot p_0$ and $p=11$ has the consecutive pair (5,6).
%% FIXED Both recomputed exactly (sympy resultants; root-by-root over 2,3,7,11,29,37,557,p_0).
\begin{lemma}[The fully degenerate steps]\label{lem:degen}
$\PROVED$ (explicit finite set). Let $p\ge5$ and let $\nu$ be a root of $a_0$ modulo $p$ with
$1\le\nu\le p-4$. Then $V(\nu)\equiv0$ forces $p$ to divide all
three resultants $\Res_\nu(a_0,c_i(\nu-1))$, whose gcd is
$2^6\cdot3^3\cdot7^3\cdot11\cdot29^3\cdot37^2\cdot557^3\cdot p_0$, and a root-by-root check over
every prime in that gcd gives exactly three occurrences:
\[
(p,\nu)=(7,2),\qquad (11,6),\qquad (p_0,\,416574044722681).
\]
In all three cases
$\mathrm{row}(\nu-1)/p$ is a valid $\ZZ$-functional with a unit $Y_{\nu+2}$ coefficient;
combined with the regular row at $\nu-2$ it gives two independent functionals, and the ranks
over $\FF_p$ are: $(11,6)$ and $(p_0,\cdot)$ give $\rank=2$ both with and without $U(\nu)$
\emph{(apparent)}; $(7,2)$ gives $2$ and $3$, so $p=7$ is a finite check (levels $5,6$;
$\VERIFIED$).
\end{lemma}

\begin{remark}
The range hypothesis $\nu\le p-4$ is not cosmetic. Dropping it admits exactly one further pair,
$(p,\nu)=(29,27)$: there $a_0(27)\equiv0$ and $c_1(26)\equiv c_2(26)\equiv c_3(26)\equiv0
\pmod{29}$, the last because $\nu+2=29$. It is harmless for us --- the forward induction of
\S\ref{sec:recstar} only ever uses steps $\nu\le p-4$, and at $p=29$ the above-midpoint range is
$[(p-3)/2,\,p-4]=[13,25]$ --- but it is a genuine fourth solution of the unrestricted system.
(At $p=2,3$ the pair $(p,1)$ is likewise fully degenerate; both are excluded by $p\ge5$.)

Two of the three listed cases arise from consecutive roots of $a_0$, and there are exactly two
such primes, not one: $\Res_\nu\bigl(a_0(\nu),a_0(\nu-1)\bigr)=-2^2\cdot11\cdot37^2\cdot557^3
\cdot p_0$, and of its prime divisors only $p=11$ (roots $5,6,9$; the pair $(5,6)$) and $p=p_0$
(roots $416574044722680$, $416574044722681$, $423023739121937$) actually carry a consecutive
pair --- $a_0$ has no root at all modulo $37$, a single root $49$ modulo $557$, and a single
root $1$ modulo $2$. The case $(11,6)$ is thus produced by the same mechanism as $(p_0,\cdot)$,
while $(7,2)$ and the out-of-range $(29,27)$ are not.
\end{remark}

\begin{evidence}\label{ev:above}
$\VERIFIED$ Sweep over all $428$ primes $5\le p<3000$: $296\,267$ steps above the midpoint,
$223$ exceptional, \emph{every one} an $a_0$-root, $0$ exceptions; $2$ fully degenerate, both
predicted by the resultant computation.
\end{evidence}

\subsection{Assembly, and a corollary that was expected to be an input}\label{sec:recstar}

\begin{proof}[Proof of Theorem~\textup{\ref{thm:baseintro}}]
For $n\le(p-1)/2$, Lemma~\ref{lem:below}. For $n=(p+1)/2$, Theorem~\ref{thm:nucleus}. For
$n>(p+1)/2$, forward induction from the four consecutive integral levels below: each step
$\nu\ge(p-3)/2$ has $p\nmid(\nu+3)(2\nu+5)$, so is exceptional only if $p\mid a_0(\nu)$, and
then it is apparent by Lemma~\ref{lem:APP} (Lemma~\ref{lem:degen} for the three fully
degenerate cases; $p=7$ by the finite check).
\end{proof}

\begin{proposition}[The universal midpoint row]\label{prop:univrow}
$\PROVED$ $n_0:=(p-5)/2$ is the unique $n$ with $2n+5=p$, so $n_0\equiv-5/2\pmod p$, and since
$\deg c_i=9$,
\[
2^9c_i(n_0)\equiv2^9c_i(-5/2)=28\,R_i\pmod p,\qquad (R_0,R_1,R_2,R_3)=(11907,-334374,-19292,0);
\]
equivalently, as an identity in $\ZZ[m]$, $128\,c_i(m)=7R_i+(2m+5)H_i(m)$. The content
$28=2^2\cdot7$ is a $p$-unit for every $p\ge5$ except $p=7$, where the whole row vanishes mod
$p$ and the recurrence obligation at $n_0$ is vacuous. The row $(R_0,R_1,R_2)$ is primitive
($11907=3^5\cdot7^2$, $334374=2\cdot3\cdot23\cdot2423$, $19292=2^2\cdot7\cdot13\cdot53$).
\end{proposition}

\begin{corollary}[$(\mathrm{REC}\text{-}\star)$]\label{cor:recstar}
$\PROVED$ For every prime $p\ge5$, with $n_0=(p-5)/2$,
\[
11907\,P_{n_0}-334374\,P_{n_0+1}-19292\,P_{n_0+2}\ \equiv\ 0\pmod p .
\]
\end{corollary}

\begin{proof}
$c_0(n_0)P_{n_0}+c_1(n_0)P_{n_0+1}+c_2(n_0)P_{n_0+2}=-c_3(n_0)P_{(p+1)/2}$ has
$v_p\ge v_p(c_3(n_0))\ge1$ by Theorem~\ref{thm:baseintro}; multiply by $2^9/28$ and use
Proposition~\ref{prop:univrow}. At $p=7$ use the finite check.
\end{proof}

\emph{The order of implication is inverted relative to expectation:}
$(\mathrm{REC}\text{-}\star)$ was the anticipated \emph{input} to the base case; it is its
\emph{output}. And its universal normalised row is \emph{identical} to the row governing the
$a=1$ midpoint band one digit up in the earlier campaign of this program \cite{ProgA1MID} ---
the base level and its $a=1$ lift are governed by the same universal row.

\begin{evidence}\label{ev:recstar}
$\VERIFIED$ $0$ failures, exact \texttt{Fraction} ladders, all $44$ primes $5\le p\le199$:
$v_p\bigl(11907P_{n_0}-334374P_{n_0+1}-19292P_{n_0+2}\bigr)=1$ \emph{exactly} at every one, so
the congruence is tight everywhere. The corresponding quantity for $Q$ is $1$ (except $p=29,131$
where it is $2$) --- automatic, since $Q\in\ZZ$ solves $\LBZ$ --- and for $\hat P$ it is $0$ at
$42$ of the $44$ primes, which is the control showing the congruence is real content.
A previously reported anomaly at $p=13$ was an artefact of the \emph{un-normalised} row:
writing $c_i(n_0)=uR_i+pe_i$ with $u=28\cdot2^{-9}$, the raw combination is
$u\varphi_R(P)+p\sum e_iP_{n_0+i}$, whose valuation is $\ge1$ always and jumps to $2$ exactly
when a second-order term vanishes mod $p$, as it happens to at $p=13$. In the normalised row
$p=13$ is tight like every other prime.
\end{evidence}

\begin{remark}[two harmless primes]
$v_p(c_3(n_0))=1+v_p(a_0(n_0))$ (as $n_0+3=(p+1)/2\in[1,p-1]$ and $2$ are units), and
$8a_0(-5/2)=-241144=-2^3\cdot43\cdot701$, so $p\mid a_0(n_0)$ exactly for $p\in\{43,701\}$,
both with multiplicity $1$. There the recurrence route would need $v_p\ge2$; the proof above
never uses the midpoint recurrence step, so $\{43,701\}$ cost nothing.
\end{remark}

%%%%%%%%%%%%%%%%%%%%%%%%%%%%%%%%%%%%%%%%%%%%%%%%%%%%%%%%%%%%%%%%%%%%%%%%%%%%%%%%
\section{Two structural theorems}\label{sec:wall}
%%%%%%%%%%%%%%%%%%%%%%%%%%%%%%%%%%%%%%%%%%%%%%%%%%%%%%%%%%%%%%%%%%%%%%%%%%%%%%%%

This section is not needed for Theorem~\ref{thm:main}; it explains why the proof has the shape
it has, and we believe the statements are of independent interest.

\subsection{The base case is one identity}

Let $n<p$. At $L=0$ one has $v_pT=\alpha+\gamma+\kappa$ \emph{exactly} (by Kummer:
$v_p\binom{n+k}{n}=\alpha$, $v_p\binom nl=v_p\binom nk=0$, $v_p\binom{n+k+l}{n}=\kappa$, all
arguments being $<2p$ and $n,k,l<p$). Write $\vfive=\sum_jK_ju^j$ as in \eqref{eq:Kexp} and
$S_j:=\sum_{k,l}T(n,k,l)K_j(n,k,l)$.

\begin{proposition}\label{prop:V2vac}
$\PROVED$ For $\wfive=w_5^{\mathrm{allp}}$, $K_2=0$ identically on every pattern with $s\le1$ ---
strictly stronger than the cap $J=1+\min(s,2)$, which permits $K_2\ne0$ at $s=1$. Hence at
$L=0$ every cell with $K_2\ne0$ has $v_pT=s\ge2$, so $(T/p)K_2\equiv0\pmod p$ cell by cell.
\end{proposition}

\begin{proposition}\label{prop:V3}
$\PROVED$ For $p\ge5$ and $n<p$, the base case is \emph{equivalent} to the single congruence
\[
(\mathrm V3)_0\qquad \sum_{k,l=0}^{n}\frac{T(n,k,l)}{p^2}\,K_3(n,k,l)\ \equiv\ 0\pmod p,
\]
a \emph{full} double sum: $K_3=0$ off the five patterns with $s\ge2$, and on those $v_pT=s\ge2$,
so every summand is $p$-integral.
\end{proposition}

\begin{proof}
$v_p(W_n)\ge0\iff v_p(P_n)\ge0$ ($n<p$ makes $H^{(5)}_nQ_n$ integral). $W_n=\sum_jp^{-j}S_j$
with $S_0\in\Zp$; $K_1\ne0\Rightarrow s\ge1\Rightarrow v_pT\ge1$, so $v_p(S_1)\ge1$;
$K_2\ne0\Rightarrow s\ge2$ by Proposition~\ref{prop:V2vac}, so $v_p(S_2)\ge2$; and
$K_4=K_5=0$. Hence $v_p(W_n)\ge0\iff v_p(S_3)\ge3$.
\end{proof}

With $q=p-n$, $\tilde n=2n-p$, the $s=2$ locus is exactly
\[
\begin{gathered}
\mathrm I:\ (\alpha,\gamma,\kappa)=(0,1,1),\ k<q,\ l\ge q;\qquad
\mathrm{II}:\ (1,0,1)\ \text{(the mirror)};\\
\mathrm{III}:\ (1,1,0),\ k,l\ge q,\ p\le k+l<p+q,
\end{gathered}
\]
and $\Sigma_{\mathrm I}=\Sigma_{\mathrm{II}}$, so $(\mathrm V3)_0\iff
2\Sigma_{\mathrm I}+\Sigma_{\mathrm{III}}\equiv0$. Theorem~\ref{thm:oneband} improves this to
one region.

\begin{evidence}\label{ev:V3}
$\VERIFIED$ $0$ failures. The residue density $R(k,l):=p\,T\vfive\bmod p$ is supported
\emph{exactly} on $\mathrm I\cup\mathrm{II}\cup\mathrm{III}$, is $k\leftrightarrow l$ symmetric,
and $\sum_{k,l}R=0$ for all $66$ levels $n<p$, $p\in\{5,7,11,13,17,19\}$; row sums are
\emph{not} $0$. Graded checks at $L=0$ (reconstruction of $\vfive$ from the $K_j$, integrality
of each $K_j$, the caps, and $v_p(S_j)\ge j$): $1\,516$ cells, $p\le13$, $0$ failures, with
$v_p(S_3)=3$ attained at most $n$. At $p=199$ the identity holds at all $99$ levels
$n=100,\dots,198$ and at all $15$ levels tested at $p=31$ --- by far the strongest single
verification, in exact $\FF_p$ arithmetic. Graded $L=1$ checks at $p=5$, $n\le14$: $0$
failures.
\end{evidence}

\begin{remark}[it is irreducibly a congruence]
$(\mathrm V3)_0$ can be completely de-$p$-adicised. Using Wilson's theorem, the reflections
$H_{p-1-j}\equiv H_j$ and $H^{(2)}_{p-1-j}\equiv-H^{(2)}_j$ ($p\ge5$), and the carry binomial
$\binom{p+m}{n}/p\equiv(-1)^{n-m+1}/(n\binom{n-1}{m})$ for $0\le m<n<p$, every ingredient
becomes an explicit rational function of $(n,q)$ alone, and $(\mathrm V3)_0$ becomes
$F(n,p-n)\equiv0\pmod p$ for an explicit finite double sum $F$ in which $p$ does not occur.
$\VERIFIED$ for all $24$ pairs $(n,q)$ with $2\le n\le12$, $q=p-n$, $p=n+q$ prime $\ge5$.
\emph{But $F(n,q)\ne0$ as a rational number} for every tested pair: the reduction consumes
Wilson and the reflections, so $(\mathrm V3)_0$ is genuinely a mod-$p$ congruence, not the
specialisation of a rational identity. In particular creative telescoping has nothing to
telescope here.
\end{remark}

\subsection{The wall}

\begin{theorem}[The $\theta$ wall]\label{thm:wall}
$\PROVED$ Let $\Dcal\subset\QQ^{448}$ be the space of weight-$5$ forms satisfying the
\textup{(DEPTH)} conditions ($\dim\Dcal=406$), and for $w\in\Dcal$ put
$V_w(n):=\sum_{k,l}T(n,k,l)w(n,k,l)$. \textup{(DEPTH)} gives $v_p(V_w(n))\ge-1$ for $n<p$, and
the entire content of the cancellation is the linear functional
\[
\theta_{p,n}:\Dcal\to\FF_p,\qquad \theta_{p,n}(w):=\bigl(p\,V_w(n)\bigr)\bmod p .
\]
Then $\theta_{p,n}(w)$ depends on $w$ only through the single rational number $V_w(n)$.
Consequently:
\begin{enumerate}[leftmargin=1.9em,label=\textup{(\arabic*)}]
\item every exact identity $\sum_{k,l}T\cdot M=0$ lies in $\ker V$ and therefore satisfies
  $\theta_{p,n}(M)=0$ trivially; adding any such $M$ to $\wfive$ changes neither $V$ nor
  $\theta$;
\item $(\mathrm V3)_0$ for $\wfive$ is \emph{literally} the assertion $v_p(P_n)\ge0$ and
  carries no information beyond it.
\end{enumerate}
\end{theorem}

\begin{proof}
Immediate from the definition: $V$ is linear and $\sum T\,M=0$ means $M\in\ker V$.
\end{proof}

\begin{proposition}\label{prop:thetanonzero}
$\VERIFIED$ $\theta\not\equiv0$ on $\Dcal$: three independent random $w\in\Dcal$ (each with
$\approx420$ nonzero coefficients, drawn by rref-completion of the $\QQ$ condition system) give
$\min_{n<p}v_p(\sum_{k,l}T(n,k,l)w(n,k,l))=-1$ at $p=5,7,11,13$ --- all $12$ (trial, prime)
pairs --- with cell-wise minimum also $-1$.
\end{proposition}

\begin{corollary}[the wall]\label{cor:wall}
The \textup{(DEPTH)} calculus \emph{plus arbitrarily much summand-side identity work} cannot
prove the base case. The property ``$v_p(V_w(n))\ge0$ for all $p\ge5$, $n<p$'' is a
$p$-\emph{dependent} linear condition on $w$ for each $(p,n)$, is false at a generic point of
$\Dcal$, and can only be enforced by conditions that pin the value sequence $V_w$ --- i.e.\ by
the fitting system itself.
\end{corollary}

This is the structural explanation of Theorem~\ref{thm:DEPTHplus}: those systems were trying to
buy, with $p$-independent linear conditions, something that is not $p$-independent-linear. It
also explains why the proof of \S\ref{sec:nucleus} had to import a property of the
\emph{number} $P_n$ --- which it does twice, through the decomposition identity and through
$\LBZ$.

\begin{remark}[a reflection route, refuted]
One might hope to pair the deficit cells by an involution. This fails concretely. The attained
cell $\bigl(\tfrac{p+1}{2},0,\tfrac{p-1}{2}\bigr)$ is the \emph{corner} of region $\mathrm I$,
not a reflection centre. In coordinates $(u,v)=(\lambda,k+\lambda)$ on $\mathrm I$ and
$(u,v)=(\tilde n-\lambda',\kappa')$ on $\mathrm{III}$, both regions are width-$q$ bands
$0\le v-u\le q-1$ and $\mathrm{III}=\mathrm I\cap\{v\le\tilde n\}$; so the index sets \emph{are}
related by $\lambda'\mapsto\tilde n-\lambda'$, but the summands are not --- the two $T/p^2$
weights carry $\binom n{v-u}^2\binom n{\tilde n-u}^2/\binom{n-1}{u}$ and
$-\binom n{\tilde n-v}^2\binom nu^2/\binom{n-1}{\tilde n-u}$, which no substitution matches.
Numerically, $|\mathrm I\cup\mathrm{II}|\ne|\mathrm{III}|$ ($12$ vs $3$ at $p=7$, $n=4$; $16$
vs $5$ at $p=7$, $n=5$) and the residue multisets differ ($\{1,2,3,5,1,6\}\times2$ vs
$\{1,1,4\}$ at $p=7$, $n=4$). And by Theorem~\ref{thm:wall} even a successful pairing would
have been a kernel identity and could not have proved the base case.
\end{remark}

\subsection{The weight-2 residue identities}

Lemma~\ref{lem:Phi} is the weight-$1$ member of a family. Here is the weight-$2$ level; it is
new, and by Theorem~\ref{thm:wall} it cannot prove the base case --- but it is exactly the kind
of kernel element by which the $\wfive$ family is translated, and it is reusable.

\begin{lemma}[$\Phi_2$]\label{lem:Phi2}
$\PROVED$ Fix $n\ge0$ and $0\le l\le n$; write $\Phi:=A_1(k)+2B_1(k)+C_1$ in the level-$n$
letters. Then
\[
\begin{aligned}
\text{(P1)}&& \sum_{k=0}^{n}T(n,k,l)\bigl[\Phi\,A_1(k)-A_2(k)\bigr]&=0,\\
\text{(P2)}&& \sum_{k=0}^{n}T(n,k,l)\bigl[\Phi\,C_1-C_2\bigr]&=0,\\
\text{(P3)}&& \sum_{k=0}^{n}T(n,k,l)\bigl[\Phi\,B_1(k)-\tfrac12(\Phi^2-A_2(k)-C_2)\bigr]&=0 .
\end{aligned}
\]
$\text{(P1)}+2\text{(P3)}+\text{(P2)}$ is the trivial $\Phi^2-\Phi^2=0$, so exactly two of the
three are independent; with $\text{(P0)}:\sum_kT\Phi=0$ they span a $3$-dimensional space of
exact $k$-row identities. The $k\leftrightarrow l$ mirrors hold by symmetry of $T$.
\end{lemma}

\begin{proof}
Let $G(x):=\prod_{i=1}^n(x+i)\prod_{i=1}^n(x+l+i)\big/\prod_{j=0}^n(x-j)^2$, so
$\deg\mathrm{den}-\deg\mathrm{num}=2$ and $\Res_{x=k}G=(T(n,k,l)/K_l)\Phi(k,l)$ with $K_l$
independent of $k$ (Lemma~\ref{lem:Phi}).

\emph{(P1).} For $1\le i\le n$ put $G_i:=G(x)/(x+i)$. Since $(x+i)$ divides the numerator of
$G$, $G_i$ has no pole at $x=-i$; its only poles are the double poles at $x=0,\dots,n$, and
$\deg\mathrm{den}-\deg\mathrm{num}=3\ge2$, so all residues sum to $0$. Writing $(x-k)^2G=N/E$
with $(N/E)(k)=T(n,k,l)/K_l$ and $(N/E)'(k)=(T/K_l)\Phi(k,l)$, we get
$\Res_{x=k}G_i=(T/K_l)[\Phi/(k+i)-(k+i)^{-2}]$. Sum over $k$, then over $i=1,\dots,n$, using
$\sum_i(k+i)^{-1}=A_1(k)$ and $\sum_i(k+i)^{-2}=A_2(k)$.

\emph{(P2).} Identical with $G(x)/(x+l+i)$ (again a numerator factor), using
$\sum_i(k+l+i)^{-1}=C_1$ and $\sum_i(k+l+i)^{-2}=C_2$.

\emph{(P3).} Take $G_j:=G(x)/(x-j)$, $0\le j\le n$; now $x=j$ is a \emph{triple} pole and
$\deg\mathrm{den}-\deg\mathrm{num}=3$. For $k\ne j$,
$\Res_{x=k}G_j=(T_k/K_l)[\Phi(k)/(k-j)-(k-j)^{-2}]$. At $x=j$, with $\Lambda:=\log(N/E)$,
$\Res=\tfrac12(N/E)''(j)=\tfrac12(T_j/K_l)[\Phi(j)^2+\Lambda''(j)]$ and
$\Lambda''(j)=-A_2(j)-C_2+2(H^{(2)}_j+H^{(2)}_{n-j})$. Sum over $j=0,\dots,n$ and use
$\sum_{j\ne k}(k-j)^{-1}=-B_1(k)$, $\sum_{j\ne k}(k-j)^{-2}=H^{(2)}_k+H^{(2)}_{n-k}$. The two
occurrences of the \emph{non-alphabet} symbol $H^{(2)}_k+H^{(2)}_{n-k}$ cancel exactly ($-1$
from the double poles, $+1$ from the triple pole), leaving (P3).
\end{proof}

The cancellation of $H^{(2)}_j+H^{(2)}_{n-j}$ in (P3) is why a \emph{closed} weight-$2$ identity
exists inside the $\{A,B,C,N\}$ alphabet at all; the same construction with
$G(x)/((x+i)(x+i'))$ produces the weight-$3$ level. $\VERIFIED$ exact over $\QQ$, $0$ failures:
all $120$ pairs $(n,l)$ with $0\le l\le n\le14$, each of (P0)--(P3) exactly $0$.

%%%%%%%%%%%%%%%%%%%%%%%%%%%%%%%%%%%%%%%%%%%%%%%%%%%%%%%%%%%%%%%%%%%%%%%%%%%%%%%%
\section{The middle row}\label{sec:middle}
%%%%%%%%%%%%%%%%%%%%%%%%%%%%%%%%%%%%%%%%%%%%%%%%%%%%%%%%%%%%%%%%%%%%%%%%%%%%%%%%

The weight-$3$ companion is proved by the same machinery, and it is instructive because it
\emph{fails} in a way the top row does not. Throughout, $n=ap+r$ with $p\ge5$, $1\le a<p$,
$0\le r<p$, and
\[
v(n,k,l):=\wthree(n,k,l)-H^{(3)}_n,\qquad
\hat W_n=\hat P_n-H^{(3)}_nQ_n=\sum_{k,l}T(n,k,l)\,v(n,k,l)
\]
by Theorem~\ref{thm:B}; the constant letter is exactly the $H$-layer removed by
Theorem~\ref{thm:C}, and $v$ has no constant letter. Level-$a$ letters are written
$A^{(a)}_m(b)=H^{(m)}_{a+b}-H^{(m)}_b$, $B^{(a)}_m(b)=H^{(m)}_{a-b}-H^{(m)}_b$,
$C^{(a)}_m(b,c)=H^{(m)}_{a+b+c}-H^{(m)}_{b+c}$.

\begin{lemma}[Letter descent]\label{lem:G}
$\PROVED$ In-regime (Definition~\ref{def:regime}) one has
$\lfloor(n+k)/p\rfloor=a+b$, $\lfloor(n-k)/p\rfloor=a-b$, $\lfloor(k+l)/p\rfloor=b+c$,
$\lfloor(n+k+l)/p\rfloor=a+b+c$, and hence by Lemma~\ref{lem:H} applied twice
\[
p^mA_m(k)=A^{(a)}_m(b)+p^m\Zp,\qquad p^mB_m(k)=B^{(a)}_m(b)+p^m\Zp,\qquad
p^mC_m=C^{(a)}_m(b,c)+p^m\Zp .
\]
Consequently, for every in-regime $(k,l)$,
$p^3\,T(n,k,l)\,v(n,k,l)\equiv T(n,k,l)\,v(a,b,c)\pmod p$. Off-regime the same congruence holds,
using Lemma~\textup{\ref{lem:Dpp}}.
\end{lemma}

\begin{proof}
For the displayed identities, e.g.\
$A_m(k)=[U_m(n+k)-U_m(k)]+p^{-m}[H^{(m)}_{a+b}-H^{(m)}_b]$ by Lemma~\ref{lem:H}, and the
bracket is $A^{(a)}_m(b)$; multiply by $p^m$. (For $B$ one needs $s\le r$ to get
$\lfloor(n-k)/p\rfloor=a-b$; for $C$ one needs $s+t<p$ and $r+s+t<p$.)
Now let $M=L_1L_2$ (or $L_1$) be one of the five monomial types of $v$, of total weight $3$;
then $p^3M=\prod(L_i^{(a)}+p^{m_i}z_i)$ with $z_i\in\Zp$, so
$p^3M-M^{(a)}=\sum_{\emptyset\ne S}\prod_{i\in S}p^{m_i}z_i\prod_{i\notin S}L_i^{(a)}$. The
level-$a$ letters obey $B^{(a)}_m\in\Zp$ (both arguments $<p$), $C^{(a)}_m\in p^{-m}\Zp$, and
$A^{(a)}_m(b)\in\Zp$ unless $a+b\ge p$, in which case $A^{(a)}_m(b)\in p^{-m}\Zp$. Checking the
five types: for $A_3(k)$ the error is $p^3z$; for $A_2(k)X_1$ with
$X_1\in\{A_1(k),B_1(k),C_1,A_1(l)\}$ the errors are $p^2z_1X_1^{(a)}$ (and
$X_1^{(a)}\in p^{-1}\Zp$, so this is in $p\Zp$), $p\,z_2A^{(a)}_2(b)$, and $p^3z_1z_2$. The only
dangerous term is $p\,z_2A^{(a)}_2(b)$, which lies in $p^{-1}\Zp$ \emph{only if} $a+b\ge p$ ---
and then Lemma~\ref{lem:D} gives $v_p(T(n,k,l))\ge2$, so after multiplying by $T$ it lies in
$p\Zp$. Off-regime the digit identities acquire a $+1$ and the mismatch of $p^mA_m(k)$ against
$A^{(a)}_m(b)$ is $\varepsilon_1(a+b+1)^{-m}$ with $\varepsilon_1=\ii{r+s\ge p}$, a $p$-unit
except when $a+b+1=p$; there Lemma~\ref{lem:Dpp} supplies $v_p(T)\ge4$, which the $A_3$
monomial needs, and $\ge3$ suffices for the $A_2X_1$ monomials.
\end{proof}

\begin{proof}[Proof of Theorem~\textup{\ref{thm:Eintro}}]
By Theorem~\ref{thm:C}, $p^3\hat P_n-\hat P_aQ_r\equiv p^3\hat W_n-\hat W_aQ_r\pmod p$. By
Theorem~\ref{thm:B} and Lemma~\ref{lem:G},
\[
p^3\hat W_n=\sum_{k,l}p^3T(n,k,l)v(n,k,l)\equiv\sum_{k,l}T(n,k,l)v(a,b,c)
=\sum_{b,c}v(a,b,c)\,\Tcal(b,c)\pmod p .
\]
Put $d(b,c):=\max(0,-v_p(v(a,b,c)))\in\{0,1,2,3\}$. Every monomial of $v$ contains an $A_2$ or
an $A_3$ letter, so $d(b,c)\ge2$ forces $a+b\ge p$ or $a+c\ge p$; and then
$v_p(T(a,b,c))\ge2$, whence $d(b,c)\le1+\min(v_pT(a,b,c),2)$. Therefore Theorem~\ref{lem:F}
applies and $v(a,b,c)\bigl[\Tcal(b,c)-(Q_n/Q_a)T(a,b,c)\bigr]$ has
$v_p\ge-d(b,c)+1+d(b,c)=1$, so
\[
p^3\hat W_n\equiv\frac{Q_n}{Q_a}\sum_{b,c}v(a,b,c)T(a,b,c)=\frac{Q_n}{Q_a}\hat W_a\pmod p .
\]
Hence $p^3\hat W_n-\hat W_aQ_r\equiv\lambda\hat W_a$ with $\lambda:=Q_n/Q_a-Q_r$ and
$v_p(\lambda)\ge1$ by Theorem~\ref{thm:Aintro}. Finally
$\hat W_a=\hat P_a-H^{(3)}_aQ_a$ with $H^{(3)}_a\in\Zp$ (as $a<p$) and $Q_a\in\ZZ$, so
$v_p(\hat W_a)\ge\min(0,v_p(\hat P_a))$ and $v_p(\lambda\hat W_a)\ge1+\min(0,v_p(\hat P_a))$.
\end{proof}

\begin{corollary}\label{cor:Eprime}
$\PROVED$ --- no Lemma~\ref{lem:F}, no off-regime descent, no integrality hypothesis. For
$1\le a<p/3$, every level-$a$ letter is $p$-integral ($a+b\le2a<p$ kills the $A$-poles,
$a+b+c\le3a<p$ kills the $C$-poles, $B$ never has one), hence $v(a,b,c)\in\Zp$,
$d(b,c)=0$, and Lemma~\ref{lem:F} degenerates to Lemma~\ref{lem:L4}. Therefore
\[
p^3\hat P_{ap+r}\equiv\hat P_a\,Q_r\pmod p\qquad\text{for all } 1\le a<p/3,\ 0\le r<p,
\]
and $\hat P_a\in\Zp$ automatically in that range.
\end{corollary}

\subsection{Why the product form fails, exactly}

\begin{proposition}[Size of the middle row]\label{prop:Phatsize}
$\PROVED$, conditional on Theorem~\ref{thm:B}. For $n<p^2$ and $p\ge5$,
$v_p(\hat P_n)\ge-4$; and for $n<p$, $v_p(\hat P_n)\ge-1$.
\end{proposition}

\begin{proof}[Proof sketch]
Inspect $\wthree$'s letters. $H^{(3)}_n=(\Zp)+p^{-3}H^{(3)}_a$ with $H^{(3)}_a\in\Zp$: pole
order $3$. $A_3(k)=(\Zp)+p^{-3}(H^{(3)}_\beta-H^{(3)}_b)$ with $\beta=\lfloor(n+k)/p\rfloor$ and
$b=\lfloor k/p\rfloor<p$; since $\beta\le2a+1<2p$, $H^{(3)}_\beta=(\Zp)+p^{-3}\ii{\beta\ge p}$,
so $A_3(k)=p^{-6}\ii{\beta\ge p}+p^{-3}(\Zp)+(\Zp)$ --- only the indicator term reaches depth
$6$, and by Lemma~\ref{lem:Dplus} every $(k,l)$ contributing to it has $v_pT\ge2$: pole order
$\le4$ after weighting. For $A_2(k)X_1$ the deepest term is $A_2(k)C_1$, which reaches $p^{-4}$
already with $\beta<p$, whenever $\lfloor(n+k+l)/p\rfloor\ge p>\lfloor(k+l)/p\rfloor$; that case
also gives $-4$. Taking the minimum gives $-4$. For $n=a<p$ the only offending letter is
$A_3(k)$ with $a+k\ge p$ (no $p$-divisible $m\le a$ exists, so $H^{(3)}_a$, $B_r$ and $C_r$ are
all $p$-integral), giving $-3+2=-1$ by Lemma~\ref{lem:D}.
\end{proof}

Proposition~\ref{prop:Phatsize} explains the anomaly: the na\"ive product congruence
$p^3\hat P_{ap+r}\equiv\hat P_aQ_r\pmod p$ is \emph{false}, and it fails precisely by the term
$\lambda\hat W_a$ of the proof above --- $\lambda$ supplies one power of $p$,
and $\hat W_a$ eats it back when $v_p(\hat P_a)=-1$. Equivalently, the middle Frobenius slot
needs $p^4$, not $p^3$, if one insists on an unnormalised product congruence.

\emph{The top row has no such defect}: $\VERIFIED$, $0$ exceptions for
$p\in\{5,7,\dots,31\}$, $v_p(P_a)\ge0$ \emph{and} $v_p(W_a)\ge0$ for every $1\le a<p$, whereas
$v_p(\hat P_a)=-1$ for most $a\in(p/2,p)$. That is why $(\mathrm{LB}_5)$ and $(\mathrm W5)$ are
clean and need no integrality hypothesis.

\begin{evidence}\label{ev:middle}
$\VERIFIED$ $0$ failures, all $1\le a<p$, $0\le r<p$, $n=ap+r\le360$,
$p\in\{5,7,11,13,17,19\}$: $v_p(p^3\hat P_n-\hat P_aQ_r)\ge1+\min(0,v_p(\hat P_a))$ with both
floors attained --- the minimum is $1$ on the cells with $v_p(\hat P_a)=0$ and $0$ on those with
$v_p(\hat P_a)=-1$. The plain product form fails on $3/5/24/55/84/128$ of the
$5/7/33/65/119/152$ cells with $v_p(\hat P_a)=-1$ at those primes and on \emph{none} of the
cells with $\hat P_a\in\Zp$. Also $\VERIFIED$: $v_p(\hat P_a)\ge0\Rightarrow v_p(\hat P_{ap+r})\ge-3$
for every $r$, and every cell with $v_p(\hat P_n)=-4$ has $v_p(\hat P_a)=-1$; and
$\min_{n<p^2}v_p(\hat P_n)=-4$ for $p\in\{5,\dots,23\}$, $\min_{a<p}v_p(\hat P_a)=-1$ for every
$p\in\{5,\dots,31\}$ --- both bounds of Proposition~\ref{prop:Phatsize} attained. The
\emph{master} form $p^3\hat P_nQ_a\equiv\hat P_aQ_n\pmod p$ holds on all cells; it does not
imply the product form when $p\mid Q_a$ or $\hat P_a\notin\Zp$.
\end{evidence}

%%%%%%%%%%%%%%%%%%%%%%%%%%%%%%%%%%%%%%%%%%%%%%%%%%%%%%%%%%%%%%%%%%%%%%%%%%%%%%%%
\section{Certification status of the decomposition identities}\label{sec:certs}
%%%%%%%%%%%%%%%%%%%%%%%%%%%%%%%%%%%%%%%%%%%%%%%%%%%%%%%%%%%%%%%%%%%%%%%%%%%%%%%%

This section states, as of the time of writing, exactly what is and is not certified. It is
written so that a single sentence flips when a certificate lands.

\subsection{Status}

The table records the state at the date of this preprint. Both $\VERIFIED$ entries are the
subject of ongoing machine computation; if either is upgraded, the only change to this paper is
the corresponding row and the sentence in Remark~\ref{rem:flip}.

\begin{center}
\begin{tabular}{@{}lll@{}}
\toprule
object & class & basis \\
\midrule
$Q_n=\sum_{k,l}T(n,k,l)$ &$\CERTIFIED$& two-step CT; telescoper $=\LBZ$ exactly; re-verified\\
$\LBZ\cdot T=\Delta_k(\rho T)+\Delta_l(\sigma T)$ &$\CERTIFIED$& explicit $\rho,\sigma$; checked to $0$ twice\\
(DEPTH) linear system &$\CERT$& $\rank$(joint)$\,=\,\rank$(aug)$\,=324$, two primes\\
$w_5^{\mathrm I}$ cell-wise conditions &$\CERT$& three primes, identical pivot sets\\
Theorem~\ref{thm:B} \ $\hat P_n=\sum T\wthree$ &$\VERIFIED$& see \S\ref{sec:certsB}\\
(T1-top) \ $P_n=\sum T\wfive$ &$\VERIFIED$& see \S\ref{sec:certsT1}\\
\bottomrule
\end{tabular}
\end{center}

\subsection{(T1-top)}\label{sec:certsT1}

\begin{evidence}\label{ev:T1}
$\VERIFIED$ For $\wfive=w_5^{\mathrm{allp}}$, evaluating $\sum_{k,l}T(n,k,l)\wfive(n,k,l)$
directly from the saved representative --- a code path different from the linear algebra that
produced it --- modulo $q=33554393$ and $q=33554467$ for $n=0,\dots,750$:
\[
\begin{array}{ll}
\text{against the exact ladder } P_n,\ \text{every } n\le360 & \textbf{0 mismatches, both primes}\\
\LBZ\ \text{residual},\ n=0,\dots,747 & \textbf{0 at all 748 values, both primes}\\
\text{minimal recurrence (search over order}\le12,\ \deg\le30) & \textbf{(order 3, degree 9), nullity 1}
\end{array}
\]
The fitting system imposed equations only for $n\le600$, so roughly $147$ of these are genuine
excess checks; and the minimal recurrence being exactly $\LBZ$ pins what a certificate must
produce. Separately, exact over $\QQ$: $n\le34$ for $w_5^{\mathrm{allp}}$, $n\le20$ for
%% FIXED [M-4]: ``287 of 687'' was a misreading of the source shorthand and, as printed, asserted
%% FIXED that 400 excess checks had FAILED. The two numbers are two separate all-satisfied counts:
%% FIXED 600-313=287 at N=600 (PHASE2_ENDGAME R2.1) and 1000-313=687 at N=1000
%% FIXED (PROOF_LB5_CLOSEOUT). Both now stated as complete successes, with their N named.
$w_5^{\mathrm I}$, $0$ mismatches, by standalone re-implementations reading only the stored
coefficients and the ladder; plus \emph{all} $287$ excess equations satisfied modulo two primes
at $N=600$ (that is, $600$ equations against rank $313$), and, in the earlier run at $N=1000$,
all $687$.
\end{evidence}

\begin{proposition}[No cheap representative]\label{prop:ptrank}
$\PROVED$ The $448$ basis monomials are linearly independent \emph{as functions} on the cells:
the cell-level matrix (rows the $1331$ cells of $n=17,\dots,22$, columns the $448$ symmetrised
monomials) has rank $448$ modulo $q=33554393$, hence rank $448$ over $\QQ$. Consequently
$\ker(\text{design})\cap\{\text{pointwise zero}\}=0$: every one of the $135$ kernel directions
is a genuine \emph{summation} identity $\sum_{k,l}T\cdot z=0$, of exactly the same difficulty as
the target. \emph{Certifying one representative does not certify another for free}, and the
representative one certifies must be the representative one uses.
\end{proposition}

%% FIXED [M-6]: relabelled $\PROVED$ negative -> $\CERT$ negative. This is a mod-$q$ rank
%% FIXED computation, exactly what \S\ref{sec:evidence} defines as $\CERT$, and exactly how
%% FIXED Theorem \ref{thm:DEPTHplus} (the same kind of computation) is labelled. Inconsistency
%% FIXED over $\FF_q$ does not formally imply inconsistency over $\QQ$. Also flagged in the
%% FIXED statement that the third alphabet was run at one prime only, and in \S16's table.
\begin{theorem}[Structural obstruction at weight 5]\label{thm:degfit}
$\CERT$ negative. The weight-$5$ fitting system $P_n=\sum_{k,l}T\cdot w$ is
\emph{inconsistent} when the support is restricted to letter monomials of degree $\le3$, and
the obstruction is the \emph{fit identity alone} --- not the $p$-integrality conditions --- so
no pole-cap regime can rescue it. Verified in three alphabets at $N=600$ equations,
$q=33554393$ (and re-run at $q=33554467$ for the two depth-$1$ alphabets, identical verdicts):
\[
\begin{array}{lccl}
\text{alphabet} & \text{degree cap} & \text{cols},\ \rank(\mathrm{fit}) & \text{fit alone}\\\hline
A_r,B_r,C_r,N_r & \le2 & 44,\ 41 & \text{INCONSISTENT}\\
 & \le3 & 176,\ 151 & \text{INCONSISTENT}\\
 & \le4 & 338,\ 269 & \text{INCONSISTENT}\\
 & \le5 & 448,\ 313 & \text{consistent}\\
+\,R_r(k)\ \text{(Ap\'ery)} & \le3 & 369,\ 329 & \text{INCONSISTENT}\\
+\,Y_{ab},V_{ab},Z_{ab}\ \text{(nested)} & \le3 & 488,\ 327 & \text{INCONSISTENT}
\end{array}
\]
(The third alphabet was run at $q=33554393$ only.)
The harness reproduces both a known positive ($\rank(\mathrm{fit})=313$; the $w_5^{\mathrm I}$
system with $212$ condition rows, $\rank(\mathrm{joint})=342$, consistent) and a known negative.
Every negative entry is immune to the ``vacuous consistency'' caveat: at degree $\le3$ the ranks
are $151/329/327$ against $N=600$, i.e.\ the systems are overdetermined by $449/271/273$
independent equations and fail.
\end{theorem}

\begin{remark}[what the label means here]
Inconsistency over $\FF_q$ does not formally imply inconsistency over $\QQ$: reduction can lower
both $\rank(A)$ and $\rank([A\mid b])$, and a rational solution can hide $q$ in a denominator.
Two auxiliary primes with identical verdicts make this overwhelming evidence, which is what
$\CERT$ means in \S\ref{sec:evidence}; it is not a proof. What would upgrade it to $\PROVED$ is
a finite, exactly checkable object: a rational annihilator $y$ with $y^\top A=0$ and
$y^\top b\ne0$, verified over $\QQ$. We have not produced one. Nothing else in this paper
depends on Theorem~\ref{thm:degfit}; it is an explanation of where the certification cost sits,
not a step in any proof.
\end{remark}

The consequence for certification is concrete. Weight-lowering by the certified $\rho,\sigma$
turns the target into a creative telescoping whose cost is governed by the \emph{degree} of the
monomials of $w$, because a squarefree degree-$d$ monomial contributes $2^d-1$ sub-monomials
under the shift. Measured supports of $E(w)/T$: $\mathbf 6$ for $\wthree$ (whose folded form has
degree $\le2$), $208$ for $w_5^{\mathrm{allp}}$, $220$ for $w_5^{\mathrm I}$. Degree histograms
of the known representatives ($w_5^{\mathrm I}$: $1{:}3,2{:}19,3{:}59,4{:}73,5{:}53$;
$w_5^{\mathrm{allp}}$: $1{:}2,2{:}14,3{:}52,4{:}66,5{:}44$; $\wthree$ folded: $1{:}4,2{:}2$) sit
exactly where Theorem~\ref{thm:degfit} says they must. \emph{So $\wthree$'s degree-$\le2$ folded
form --- the single structural fact that makes the weight-$3$ certificate route exist --- has no
weight-$5$ analogue.} We record this as a mathematical fact about $\wfive$, not a resource
shortfall: the cheapest certified route to (T1-top) is a creative telescoping on a
$\partial$-finite module of rank $100$ to $220$, against rank $3$ for Theorem~\ref{thm:B}, and
nothing short of a new idea about the shape of $\wfive$ changes this. Adding depth would not
have helped even had the degree test passed: a single nested letter is degree $1$ but its
$\partial$-module is not rank $2$ (shifting it produces $A_a(k),A_b(k),H^{(a)}_k,H^{(b)}_k$), so
the $2^d$ support model becomes roughly $4^d$.

\subsection{Theorem B}\label{sec:certsB}

\begin{evidence}\label{ev:B}
$\VERIFIED$ exact over $\QQ$ for $n\le40$ (and independently re-derived and re-checked for
$n\le30$, and re-implemented directly from the displayed formula \eqref{eq:w3} for $n\le16$),
$0$ discrepancies; $340$ excess equations satisfied modulo $q$ in a $93$-element basis at
$N=400$ (with a negative control: the same basis against the target $P$ is inconsistent, as it
must be); and $\hat P_n-\sum_{k,l}T\wthree=0$ exactly for $n=0,\dots,300$ --- $301$ values.
\end{evidence}

Theorem~\ref{thm:B} is a \emph{compute} node, not a structural one. It has been reduced by proof
to $\sum_{k,l}E(v)=0$ for an explicit $E(v)$; the boundary and telescoping steps are proved from
the pole structure of $\rho,\sigma$ (using $\rho(n,0,l)=\sigma(n,k,0)=0$ and the double zeros of
$T$), the regularised boundary lemma is $\VERIFIED$ exactly for $n_0=1,\dots,6$, and $E(v)/T$ is
a single hypergeometric term times a rank-$3$ factor ($E(v)/T=c_0+\beta A_2(k)+\alpha\Psi$ after
the proved relations $\gamma=3\alpha$, $\delta=\tfrac32\alpha$, $\varepsilon=\tfrac12\alpha$).
Splitting $E(v)=\sum_\tau F_\tau$ over the five shift terms before telescoping --- a split that
is itself $\CERTIFIED$, symbolically and without loading any package --- drops the weights from
$132\,917$ leaves to $\{84,86,66,12471,2255\}$ at unchanged rank $\le3$; splitting further by
letter reduces every piece to a rational multiple of $T$, i.e.\ to rank $1$. Since
$\hat P_n-\sum T\wthree=0$ is $\VERIFIED$ for $n=0,\dots,300$, any operator of order $\le298$
finishes the identity by matching initial values.

%% FIXED [Mo-1]: snapshot brought up to the 06:15 state (PHASE2\_CERTS \S\S18.18--18.19):
%% FIXED $R4$ \texttt{DFiniteTimes} has cleared, so four of the five stages are clear on every
%% FIXED $\tau$ and $ct_2$ is the sole wall, with zero returns under unconstrained search;
%% FIXED \S18.13's \texttt{guessrec} measurement is now acknowledged against the ``compute node''
%% FIXED claim; and the loose ``the same object before the split'' has been corrected.
At the date of this preprint no complete $M_\tau$ has been produced. What has changed, and is
worth recording because it inverts the earlier diagnosis, is where the cost sits. The pipeline
has five stages --- annihilator, first telescoping (eliminate $l$), Gr\"obner normalisation,
$\partial$-finite product, second telescoping --- and after the letter-split \emph{four of the
five are clear on every $\tau$ and every piece tried}. On the largest piece the annihilator
returns in $34$~s and the first telescoping in $33$~s (the $12489$-leaf term; the $13069$-leaf
unsplit object of which it is one piece had failed with an out-of-memory at $7.8$~GB after $85$
minutes), the Gr\"obner step in $1$--$5$~s, and the $\partial$-finite product --- until recently
the second wall --- in $500$~s at no measurable memory, returning $4$ generators on the largest
cofactor in the problem. The \emph{second} telescoping is now the sole remaining obstacle, and it
returns no telescoper by any method: not under any support bound attempted (a full ladder
$d=0,\dots,5$ on two pieces), and not under unconstrained search with no ansatz box at all
($421$~s on the largest piece, $600$~s on another). We therefore do not quote a kernel-hour
estimate: earlier ones were withdrawn, and we prefer to state the position as measured.

That the residual difficulty is exactly one stage sharpens, but also qualifies, the description
of Theorem~\ref{thm:B} as a compute node. A direct measurement (\texttt{guessrec} over
$\mathrm{order}\le12$, $\deg\le30$) finds that \emph{no} single-letter component sum satisfies
an operator in that box, while their combination satisfies the unique minimal one, of order $3$
and degree $9$ --- namely $\LBZ$. So the second telescoping, applied to split pieces, has been
searching an empty box: the split that made the first four stages cheap is the same split that
puts the target out of range.

%% FIXED [P1e-s7]: the structural finding of PHASE2\_CERTS \S19.1--19.2 added (the $E$-route is
%% FIXED blocked for a reason internal to $E$, measured at $7$ letters as well as at $9$); the
%% FIXED closing sentence no longer says ``algebraic rather than computational'' without
%% FIXED qualification, because the direct route is a live computation, not a dead one.
A subsequent measurement locates a second obstruction one stage earlier, and this one is
structural rather than a matter of resources. The cost model above --- support size, hence
monomial degree --- is only one of two axes: \texttt{Annihilator} needs the summand small both
in the number of distinct harmonic letters it carries \emph{and} in the size of the rational
coefficients multiplying them, and refolding $w$ can address only the first. Carried out, it
does: a refolded $\tilde v$, equal to $v$ modulo proved kernel identities of the
Lemma~\ref{lem:Phi2} family together with the $k\leftrightarrow l$ symmetry, reduces $E$ from
$9$ distinct letters to $7$. The annihilator
nevertheless becomes \emph{more} expensive, not less --- it aborts against an $8.5$~GB memory
cap after $536$~s (peak $8.4$~GB) with memory rising a straight $3.6$~GB per minute, where the
$13069$-leaf $9$-letter object above reached $7.8$~GB only after $85$ minutes. The reason is
visible in $E$ itself. Writing $E(w)=\sum_\tau F_\tau$ with $F_\tau=G_\tau\,(\tau.w-w)$ over
the five shifts, the two one-variable shifts have $G_{kk}=-(\rho\,T)|_{k\to k+1}$ and
$G_{ll}=-(\sigma\,T)|_{l\to l+1}$: the coefficients of $E(w)$ carry the very cofactors
$\rho,\sigma$ of $\LBZ\cdot T=\Delta_k(\rho T)+\Delta_l(\sigma T)$ --- the two largest rational
functions in the problem --- unless
$\tau.w-w=0$ for $\tau=kk$ and $\tau=ll$ both, i.e.\ unless $w$ is a function of $n$ alone,
which the fit for $\wthree$ excludes. \emph{So no choice of representative makes $E(w)$ cheap.}
The obstruction lives in the certificate for $Q_n$ that produced $E$, not in the weight $w$,
and it is invariant under every refolding of $w$.

Splitting further cannot help either, so what remains is the one route that never forms $E$ at
all: telescope the summand $T\,w$ directly. Its coefficients are polynomial rather than
$\Theta(\rho,\sigma)$, it needs no weight-lowering step and hence no boundary lemma for $E$,
and its second telescoping searches the support $\{1,S_n,S_n^2,S_n^3\}$ --- the one box in this
problem known to contain its answer, since $\LBZ\hat P_n=0$ is $\PROVED$ and
$\sum_{k,l}T\wthree=\hat P_n$ is $\VERIFIED$ to $n=300$. Every side condition such a certificate
would need is discharged in advance: $\sum_{k,l}T\tilde v=\hat
P_n$ exactly for $n\le33$, $\wthree-\tilde v$ is an explicit combination of proved summation
identities, the far-edge boundary
terms vanish ($T$ has a double zero at each integer $k>n$ where $\tilde v$ has at worst a simple
pole, and none in $l$), and $301$ exact initial values are banked. The next move on
Theorem~\ref{thm:B} is therefore a single annihilator on a $91$-leaf object, or else an
algebraic one; in neither case is it a larger version of the computation already attempted.

%% FIXED [M-1]: (DEPTH) restored here. A (T1-top) certificate alone does NOT make Theorems
%% FIXED \ref{thm:main}, \ref{thm:baseintro}, \ref{thm:51intro} unconditional --- (DEPTH) is
%% FIXED [CERT] and is consumed by all three.
\begin{remark}[what would flip]\label{rem:flip}
If a certificate for (T1-top) lands, Theorem~\ref{thm:main} and Theorems~\ref{thm:baseintro},
\ref{thm:51intro} rest on the single remaining input (DEPTH); they become
\emph{unconditional} only when (DEPTH) is in addition promoted from $\CERT$ to a proof, which
is a separate and much smaller task --- a rational consistency witness for a finite linear
system, of the kind described after Theorem~\ref{thm:degfit}. If a certificate for
Theorem~\ref{thm:B} lands, Theorem~\ref{thm:Eintro} becomes unconditional outright, since it
does not consume (DEPTH). No other statement in this paper changes; in particular a certificate
for Theorem~\ref{thm:B} changes nothing about Theorem~\ref{thm:main}, which does not use it.
\end{remark}

%%%%%%%%%%%%%%%%%%%%%%%%%%%%%%%%%%%%%%%%%%%%%%%%%%%%%%%%%%%%%%%%%%%%%%%%%%%%%%%%
\section{The primes 2 and 3}\label{sec:twothree}
%%%%%%%%%%%%%%%%%%%%%%%%%%%%%%%%%%%%%%%%%%%%%%%%%%%%%%%%%%%%%%%%%%%%%%%%%%%%%%%%

The constant $12=2^2\cdot3$ is not covered by anything above: $p\ge5$ is forced by the single
use of Wolstenholme inside Lemma~\ref{lem:W}, and again by $\sum_{j<p}j^{-2}\equiv0$ at the
nucleus. The residual statement is:

\begin{quote}
\textbf{(Sharp-$12$, $p\in\{2,3\}$ part).}
$\ord_2(P_n)\ge-2-5\lfloor\log_2n\rfloor$ and $\ord_3(P_n)\ge-1-5\lfloor\log_3n\rfloor$,
i.e.\ the \emph{constant} $12=2^2\cdot3$ and nothing worse.
\end{quote}

$\VERIFIED$ $361/361$, $n\le360$ (Evidence~\ref{ev:law12}).  There is now a
proved summand-level explanation of both primes.  Put
\[
 \alpha=(H_{n+k}-H_k)-(H_{n+l}-H_l),\qquad
 \beta=(H_{n-k}-H_k)-(H_{n-l}-H_l)
\]
and
\begin{equation}\label{eq:compact-low-prime}
 \omega_5(n,k,l)=H_{n+k}^{(5)}+\frac{\alpha-\beta}{2}H_{n+k}^{(4)}
 +\frac{A_2(k)+A_2(l)-\alpha^2-2\alpha\beta}{4}H_{n+k}^{(3)},
\end{equation}
where $A_2(x)=H_{n+x}^{(2)}-H_x^{(2)}$, and set
$S_n=\sum_{k,l=0}^nT(n,k,l)\omega_5(n,k,l)$.

\begin{proposition}[Endpoint-residue theorem]\label{prop:endpoint23}
$\PROVED$ For every $n\ge1$,
\[
 2^{2+5\lfloor\log_2n\rfloor}S_n\in\mathbb Z_2,
 \qquad
 3^{1+5\lfloor\log_3n\rfloor}S_n\in\mathbb Z_3.
\]
More precisely, if $N=p^{\lfloor\log_pn\rfloor}$, the only binary cells below
the termwise bound are $(N,0)$ and $(N,N)$; after multiplication by
$2^{5\lfloor\log_2n\rfloor+4}$ their residues are $-1,+1\pmod4$.  In the
ternary case there are no deficient cells when the leading ternary digit of
$n$ is $1$.  If $n=2N+r$, their number is
\[
 3^{1+e_1(r)},\qquad
 e_1(r)=\#\{\hbox{ternary digits of $r$ equal to $1$}\},
\]
and every normalized residue is $-1\pmod3$.
\end{proposition}

\begin{proof}
For $M=pN$ use the exact digit split
\[
 H_{pm+d}^{(s)}=p^{-s}H_m^{(s)}+
 \sum_{a<m}\sum_{u=1}^{p-1}(pa+u)^{-s}
 +\sum_{u=1}^{d}(pm+u)^{-s}.
\]
Writing $e_x=\ii{n+x\ge M}$, its principal part in
\eqref{eq:compact-low-prime} is determined by
\[
 M\alpha\equiv e_k-e_l,\quad M\beta\equiv0,\quad
 M^2(A_2(k)+A_2(l)-\alpha^2-2\alpha\beta)
 \equiv e_k+e_l-(e_k-e_l)^2\pmod p.
\]
The last expression is even at $p=2$.  At $p=3$ the coefficient of
$M^{-5}$ is
\[
 1+\frac{1-e_l}{2}+\frac{1+e_l-(1-e_l)^2}{4}\equiv0\pmod3
\]
when $e_k=1$; the case $e_k=0$ is one order easier.  Thus
$v_2(\omega_5)\ge-5L-6$ and $v_3(\omega_5)\ge-5L-4$.

Kummer's theorem now leaves a finite carry table.  On the minimal shell its
leading triples $(n_L,k_L,l_L)$ and lower triples are
\[
\begin{array}{c|c|c}
p&\text{leading}&\text{lower}\ \\ \hline
2&(1,1,0),(1,1,1)&(0,0,0),(1,0,0)\\
3&(2,1,0),(2,1,2),(2,2,0),(2,2,1)&
 (0,0,0),(1,0,0),(1,0,1),(1,1,0),(2,0,0).
\end{array}
\]
Equivalently, at a lower ternary digit $n_i=0,2$ one must have
$(k_i,l_i)=(0,0)$, while at $n_i=1$ the possibilities are
$(0,0),(0,1),(1,0)$.  The factorial-unit recurrences
\[
 u_p((pm+d)!)\equiv(-1)^mu_p(m!)d!\pmod p\quad(p\text{ odd}),
\]
\[
 u_2((2m+d)!)\equiv u_2(m!)\prod_{j=1}^{m+d}(2j-1)\pmod4
\]
show by substitution that every allowed lower transition has normalized
multiplier $1$.  The leading residues are $(-1,+1)$ in the two binary states;
in the four ternary states they are $(-1,-1,-1,0)$, the last state gaining an
extra factor $3$.  This gives the two binary endpoints and, digit by digit,
$3\prod_i(1,3,1)_{r_i}=3^{1+e_1(r)}$ ternary endpoints.  Their stated residue
sums vanish modulo $4$ and $3$, respectively.  All other carry shells supply
the missing valuation directly.
\end{proof}

Consequently the low-prime sharp-$12$ statement follows from the single
identity
\begin{equation}\label{eq:compact-P-bridge}
 P_n=S_n.
\end{equation}
Identity~\eqref{eq:compact-P-bridge} is $\VERIFIED$ exactly over $\mathbb Q$
on the currently tested range, but its all-$n$ residue/creative-telescoping
certificate remains open.  Thus Proposition~\ref{prop:endpoint23} is a theorem;
its transfer from $S_n$ to the recurrence-defined $P_n$ is not yet a theorem.
The distinction is important: the local denominator problem has been solved,
and the remaining obstruction is one global summation identity.

We state the older geometric provenance as well, since it explains why this
arithmetic mechanism was not visible there.

The mechanism proposed for the factor $2$ in the earlier campaign of this program is an index-$2$
computation on the Betti side of the period matrix, via the bicomplex of \cite{Du18}; its own
recorded bottom line is that
\begin{quote}
``the claim `index $2$ derived from geometry alone' is \textbf{not} yet earned; the honest status
is `index $2$ is the unique value consistent with all closed sub-computations, with the sole
open input being the integral bicomplex differential at two explicitly identified strata'.''
\end{quote}
Of the three sub-results there, two are proved and one is not: it is \emph{proved} that the
factor $2$ is absent from the arrangement combinatorics on the other side (all Smith elementary
divisors are $1$, under every orientation convention); it is \emph{proved} that only the prime
$2$ can arise from those incidence matrices; and the factor $2$ itself is \emph{not closed},
being localised to the integral weight spectral-sequence differential across the uncoloured
generic facets on two explicitly identified strata --- i.e.\ to an integral refinement of the
$\QQ$-only bicomplex of \cite{Du18} that is not in the literature.

The honest consequence is that the geometric route cannot produce a $3$.
Proposition~\ref{prop:endpoint23} identifies its arithmetic source instead: it is the
divisibility by $3$ of the ternary endpoint-fibre cardinality
$3^{1+e_1(r)}$.  What remains open is not the cancellation mechanism but the bridge
\eqref{eq:compact-P-bridge} connecting that compact harmonic representative to the
recurrence-defined top row.

The general-parameter conjecture of the earlier campaign, which our totally symmetric statement
specialises, is
\[
12\cdot D_\nu(\boldsymbol a,n)\cdot I'(\boldsymbol a\cdot n)\in\ZZ\zeta(5)+\ZZ
\]
for all admissible $\boldsymbol a$ and $n\ge1$, where $D_\nu$ is \cite{BZ}'s $\Phi_n$-sharpened
denominator with the exponents $\nu_p$ of \cite[(nu\_p)]{BZ} extended from $p>\sqrt{m_1n}$ to
%% FIXED [M-7]: the three numbers came from three different audits. Now quoting one consistent
%% FIXED pair --- the 164-cell audit of the source's updated header, 46x / 30x. (The 28x / 14x
%% FIXED figures belong to an earlier 103-cell audit; ``~215 cells'' is the count in the later
%% FIXED slack-trajectory section, a different sweep again.)
every prime \cite{ProgCONJ}. Its evidence is $164$ audited cells: excess at $p=2$ is at
most $+2$ (attained $46$ times), at $p=3$ at most $+1$ (attained $30$ times). One verified cell
at $n=1$ shows excess $+1$ at $p=7=m_1$, the boundary prime of the multiset --- a rare,
orbit-coherent, non-persistent fluctuation which a re-audit at $n=2$ flips to slack at both
partner points. That effect is outside the totally symmetric family treated here, but it is a
reminder that the general-parameter statement has a boundary-prime subtlety the symmetric one
does not.

%%%%%%%%%%%%%%%%%%%%%%%%%%%%%%%%%%%%%%%%%%%%%%%%%%%%%%%%%%%%%%%%%%%%%%%%%%%%%%%%
\section{Summary of the dependency structure}\label{sec:summary}
%%%%%%%%%%%%%%%%%%%%%%%%%%%%%%%%%%%%%%%%%%%%%%%%%%%%%%%%%%%%%%%%%%%%%%%%%%%%%%%%

\begin{center}
\footnotesize
\begin{tabular}{@{}lll@{}}
\toprule
node & class & where\\
\midrule
%% FIXED [M-1]: root rows now name both unproved inputs, (T1-top) and (DEPTH).
Theorem~\ref{thm:main} (Sharp-$12$, $p\ge5$) & $\PROVED$ modulo (T1-top) $+$ (DEPTH) & \S\ref{sec:gapdesc}\\
\quad Theorem~\ref{thm:51intro} ($-5L-1$, all $n$) & $\PROVED$ modulo (T1-top) $+$ (DEPTH) & Cor.~\ref{cor:51}\\
\quad (IND) step $\mathcal I(L-1)\Rightarrow\mathcal I(L)$ & $\PROVED$ & Cor.~\ref{cor:step}\\
\qquad fibre term (II) & $\PROVED$, slack $0$ & Prop.~\ref{prop:ledgerII}\\
\qquad\quad (DEPTH-gen) & $\PROVED$ + $\CERT$ & Cor.~\ref{cor:depthgen}\\
\qquad\qquad Prop.\ LIFT & $\PROVED$ & Prop.~\ref{prop:LIFT}\\
\qquad\quad Lemma $F$-gen & $\PROVED$ & Lem.~\ref{lem:Fgen}\\
\qquad\qquad Lemma $B$ (general $a$) $\to$ Lemma $W$ & $\PROVED$ & Lem.~\ref{lem:B}, \ref{lem:W}\\
\qquad descent (I), in-regime & $\PROVED$ & Prop.~\ref{prop:inregime}\\
\qquad descent (I), off-regime & $\PROVED$ & Thm~\ref{thm:gapdesc}\\
\qquad\quad Lemma DK, Lemma D1 & $\PROVED$ & Lem.~\ref{lem:DK}, \ref{lem:D1}\\
\quad base case $\mathcal I(0)$ & $\PROVED$ modulo (T1-top) $+$ (DEPTH) & Thm~\ref{thm:baseintro}\\
\qquad one-band reduction & $\PROVED$ + $\CERT$ & Thm~\ref{thm:oneband}\\
\qquad corner sum $=-14s_2$ & $\PROVED$ & Thm~\ref{thm:nucleus}\\
\qquad apparency of the $a_0$ steps; $\tilde L$ & $\PROVED$ & Lem.~\ref{lem:APP}, Prop.~\ref{prop:Ltilde}\\
\midrule
%% FIXED [M-1/M-2]: (T1-top) is not the only open node --- (DEPTH) is a second, and Theorem B
%% FIXED a third (consumed only by the middle row).
(T1-top) $P_n=\sum T\wfive$ & $\VERIFIED$ \textbf{-- open} & \S\ref{sec:certsT1}\\
(DEPTH) depth linear system & $\CERT$ \textbf{-- open} & Thm~\ref{thm:depthcons}\\
\midrule
%% FIXED [m-2]: the $\LEAN$ tag covers only the $\bmod\ p$ statement and its digit-product form
%% FIXED (LEAN\_LUCAS\_STATUS: Q\_lucas, Q\_lucas\_digits); the $\bmod\ p^2$ supercongruence is not
%% FIXED formalised. The row now says so, as Evidence \ref{ev:A} already did.
Theorem~\ref{thm:Aintro} (Lucas, $\bmod\ p$) & $\PROVED$, $\LEAN$ & \S\ref{sec:Qrow}\\
\quad the $\bmod\ p^2$ supercongruence & $\PROVED$ (not formalised) & Cor.~\ref{cor:Qsuper}\\
Theorem~\ref{thm:C} ($H$-layer) & $\PROVED$ & \S\ref{sec:Hlayer}\\
Theorem~\ref{thm:Eintro} (middle row) & $\PROVED$ modulo Theorem~\ref{thm:B} & \S\ref{sec:middle}\\
Theorem~\ref{thm:B} $\hat P_n=\sum T\wthree$ & $\VERIFIED$ & \S\ref{sec:certsB}\\
Theorem~\ref{thm:wall} ($\theta$ wall) & $\PROVED$ & \S\ref{sec:wall}\\
Theorem~\ref{thm:degfit} (degree obstruction) & $\CERT$ negative & \S\ref{sec:certsT1}\\
Endpoint theorem at $p\in\{2,3\}$ & $\PROVED$ for $S_n$; bridge $S_n=P_n$ open & \S\ref{sec:twothree}\\
\bottomrule
\end{tabular}
\end{center}

\medskip
\noindent
Two remarks close the summary. First, \emph{the ledger has no slack anywhere}: (DEPTH) is
attained, Lemma~\ref{lem:F} is attained, Lemma~\ref{lem:Fgen} is attained,
Lemma~\ref{lem:DK} is attained, and the induction consumes exactly what they provide, digit by
digit. That is not bookkeeping tidiness; it is why the exponent in Theorem~\ref{thm:main} is
$5$ and not $6$.
%% FIXED [m-1]: ``exactly two arithmetic inputs'' was not as stated --- Lemma \ref{lem:Tcorner}
%% FIXED uses Wilson's theorem (and Lucas and Kummer), and \S\ref{sec:wall}'s de-$p$-adicisation
%% FIXED remark uses Wilson and the reflections. The defensible claim, now made, is that exactly
%% FIXED two of them constrain the prime.
Second, \emph{the whole chain has exactly two arithmetic inputs that constrain the prime} ---
Wolstenholme's congruence, used once inside Lemma~\ref{lem:W}, and
$\sum_{j<p}j^{-2}\equiv0\pmod p$, used once at the nucleus. Both fail at $p=3$, and both are
exactly why the hypothesis of Theorem~\ref{thm:main} is $p\ge5$. The chain does use three
further classical facts --- Lucas's theorem, Kummer's theorem and Wilson's theorem (the last in
Lemma~\ref{lem:Tcorner} and in the de-$p$-adicisation remark of \S\ref{sec:wall}) --- but all
three hold at every prime and so cost nothing.

%%%%%%%%%%%%%%%%%%%%%%%%%%%%%%%%%%%%%%%%%%%%%%%%%%%%%%%%%%%%%%%%%%%%%%%%%%%%%%%%
\appendix
\section{Reproduction}\label{app:repro}
%%%%%%%%%%%%%%%%%%%%%%%%%%%%%%%%%%%%%%%%%%%%%%%%%%%%%%%%%%%%%%%%%%%%%%%%%%%%%%%%

All numerical statements in this paper are exact. The Python computations use exact rational
arithmetic, arbitrary-precision integers, $\FF_q$ linear algebra, or tracked-precision $p$-adic
arithmetic; the symbolic computations use exact rational-function
arithmetic. Design matrices over $\FF_q$ use a $13$-bit \texttt{float64} split (two
\texttt{dgemm} calls), exact for all sizes used and cross-checked against a direct
implementation. The creative-telescoping work uses the package of \cite{Ko10}; every
certificate quoted as $\CERTIFIED$ was re-verified by exact arithmetic in a session that does
not trust the search, at least once in a kernel that never loaded that package.

\medskip
\noindent The principal artefacts, by section:

\begin{center}
\small
\begin{tabular}{@{}ll@{}}
\toprule
what it settles & where\\
\midrule
consistency of (DEPTH): $\rank(\mathrm{joint})=\rank(\mathrm{aug})=324$, two primes & Thm~\ref{thm:depthcons}\\
$(\mathrm{DEPTH}^+)$, $(\mathrm{DEPTH}^{++})$ inconsistent & Thm~\ref{thm:DEPTHplus}\\
$w_5^{\mathrm{allp}}$: ladder identity, depth sweep, cell-by-cell test & Evidence~\ref{ev:depth}\\
$w_5^{\mathrm I}$: ladder identity, one-band integrality & Evidence~\ref{ev:oneband}\\
(DEPTH-gen) $+$ Kummer floor, digit levels $L\le2$ & Evidence~\ref{ev:depthgen}\\
Lemma $F$-gen at multi-digit $a$ & Evidence~\ref{ev:Fgen}\\
Lemma $B$ for general $a$ & Evidence~\ref{ev:LemB}\\
Lemma DK, its slot-wise ingredients, Theorem~\ref{thm:gapdesc} & Evidence~\ref{ev:DK}\\
base case over every prime $5\le p\le367$ & Evidence~\ref{ev:base}\\
corner identity, apparency sweep, $\tilde L$ & Evidence~\ref{ev:corner}, \ref{ev:above}\\
$(\mathrm{REC}\text{-}\star)$ normalised sweep & Evidence~\ref{ev:recstar}\\
the $\theta$ wall trials & Prop.~\ref{prop:thetanonzero}\\
the degree-cap experiment & Thm~\ref{thm:degfit}\\
(Sharp-$12$) sweeps $n\le360$ and $n\le3000$ & Evidence~\ref{ev:sweep}\\
\bottomrule
\end{tabular}
\end{center}

\begin{evidence}[the base case, directly]\label{ev:base}
%% FIXED [m-7]: the range is $1\le n$; $n=0$ included would give 11955 cells, not 11884.
$\VERIFIED$ $11\,884/11\,884$ cells, $0$ failures: $\ord_p(P_n)\ge0$ for every prime
$5\le p\le367$ and every $1\le n<\min(p,361)$, with minimum value exactly $0$.
\end{evidence}

\begin{evidence}[the main theorem, directly]\label{ev:sweep}
$\VERIFIED$ $0$ failures:
\[
\small
\begin{array}{llrl}
\text{sweep} & \text{range} & \text{cells} & \min\bigl(\ord_p(P_n)+5L\bigr)\\\hline
\text{exact ladders} & n\le360,\ p\in\{5,\dots,31\} & 3\,240 & 0\ (\text{at } p=5,\ n=1)\\
\text{same, as } \ord_p(p_n)\ge\kappa(n)-5L & n\le360,\ p\le31 & 3\,240 & \text{---}\\
\text{certified order-3 recurrence, exact }\QQ & n\le3000,\ p\in\{5,\dots,31\} & 27\,000 & 0\\
\text{recurrence vs.\ exact ladder} & n\le360 & 361 & \text{0 mismatches}
\end{array}
\]
where $p_n=\binom{2n}{n}P_n$ and $\kappa(n)=v_p\binom{2n}{n}$, so that
$\ord_p(p_n)\ge\kappa(n)-5L\iff\ord_p(P_n)\ge-5L$. The minimum slack is $0$: the bound is
attained.
\end{evidence}

%%%%%%%%%%%%%%%%%%%%%%%%%%%%%%%%%%%%%%%%%%%%%%%%%%%%%%%%%%%%%%%%%%%%%%%%%%%%%%%%
\section*{Acknowledgments}
%%%%%%%%%%%%%%%%%%%%%%%%%%%%%%%%%%%%%%%%%%%%%%%%%%%%%%%%%%%%%%%%%%%%%%%%%%%%%%%%

%% METHODOLOGY SENTENCE --- FLAGGED FOR RIVER: keep, edit, or delete before circulation.
The exploratory computations, the search for the weight-$5$ representatives, and a first
assembly of several of the arguments above were carried out with the assistance of AI language
models; every statement labelled $\PROVED$ has a human-readable proof given here, and every
numerical claim was re-verified by exact arithmetic in an independent implementation.

%%%%%%%%%%%%%%%%%%%%%%%%%%%%%%%%%%%%%%%%%%%%%%%%%%%%%%%%%%%%%%%%%%%%%%%%%%%%%%%%
\begin{thebibliography}{99}

\bibitem[Ap79]{Ap79}
R.~Ap\'ery, \emph{Irrationalit\'e de $\zeta(2)$ et $\zeta(3)$}, Ast\'erisque \textbf{61} (1979),
11--13.

\bibitem[Be79]{Be79}
F.~Beukers, \emph{A note on the irrationality of $\zeta(2)$ and $\zeta(3)$}, Bull.\ London Math.\
Soc.\ \textbf{11} (1979), 268--272.

\bibitem[Br09]{Br09}
F.~C.~S.~Brown, \emph{Multiple zeta values and periods of moduli spaces
$\overline{\mathfrak M}_{0,n}$}, Ann.\ Sci.\ \'Ec.\ Norm.\ Sup\'er.\ (4) \textbf{42} (2009),
no.~3, 371--489.

\bibitem[Br16]{Br16}
F.~Brown, \emph{Irrationality proofs for zeta values, moduli spaces and dinner parties}, Mosc.\
J. Comb.\ Number Theory \textbf{6} (2016), no.~2--3, 102--165.

\bibitem[BZ]{BZ}
F.~Brown and W.~Zudilin, \emph{On cellular rational approximations to $\zeta(5)$}, preprint
\texttt{arXiv:2210.03391v3} (2022).

\bibitem[Du18]{Du18}
C.~Dupont, \emph{Odd zeta motive and linear forms in odd zeta values}, with a joint appendix
with D.~Zagier, Compos.\ Math.\ \textbf{154} (2018), no.~2, 342--379.

\bibitem[Ko10]{Ko10}
C.~Koutschan, \emph{HolonomicFunctions (user's guide)}, RISC Report 10-01 (2010).

\bibitem[MZ20]{MZ20}
R.~Marcovecchio and W.~Zudilin, \emph{Hypergeometric rational approximations to $\zeta(4)$},
Proc.\ Edinburgh Math.\ Soc.\ \textbf{63} (2020), 374--397.

\bibitem[RV96]{RV96}
G.~Rhin and C.~Viola, \emph{On a permutation group related to $\zeta(2)$}, Acta Arith.\
\textbf{77} (1996), no.~1, 23--56.

\bibitem[RV01]{RV01}
G.~Rhin and C.~Viola, \emph{The group structure for $\zeta(3)$}, Acta Arith.\ \textbf{97}
(2001), no.~3, 269--293.

\bibitem[WZ92]{WZ92}
H.~S.~Wilf and D.~Zeilberger, \emph{An algorithmic proof theory for hypergeometric (ordinary and
``q'') multisum/integral identities}, Invent.\ Math.\ \textbf{108} (1992), no.~3, 575--633.

\bibitem[Ze91]{Zei91}
D.~Zeilberger, \emph{The method of creative telescoping}, J. Symbolic Comput.\ \textbf{11}
(1991), no.~3, 195--204.

\bibitem[Zu02]{Zu02a}
W.~Zudilin, \emph{Irrationality of values of the Riemann zeta function}, Izv.\ Ross.\ Akad.\
Nauk Ser.\ Mat.\ \textbf{66} (2002), no.~3, 49--102; English transl., Russian Acad.\ Sci.\ Izv.\
Math.\ \textbf{66} (2002), no.~3, 489--542.

\bibitem[Zu02b]{Zu02b}
W.~Zudilin, \emph{A third-order Ap\'ery-like recursion for $\zeta(5)$}, Mat.\ Zametki
\textbf{72} (2002), no.~5, 796--800; English transl., Math.\ Notes \textbf{72} (2002),
no.~5--6, 733--737.

\bibitem[Zu04]{Zu04}
W.~Zudilin, \emph{Arithmetic of linear forms involving odd zeta values}, J. Th\'eor.\ Nombres
Bordeaux \textbf{16} (2004), 251--291.

\bibitem[P-C]{ProgCONJ}
Internal report \texttt{CONJECTURE.md} of the present programme, 16 July 2026 (unpublished):
the general-parameter sharp-$12$ conjecture and its audit data.

\bibitem[P-M]{ProgA1MID}
Internal report \texttt{PHASE2\_A1\_MIDPOINT\_THEOREM.md} of the present programme,
20 July 2026 (unpublished): the complete generic midpoint theorem for the $a=1$ band.

\end{thebibliography}

\end{document}
